% !TeX root = ../Book2.tex
% 纲要标题：# 第四编　向量丛、示性类与曲率代表
% 纲要位置：几何学计划纲要.md，第 1733 行。
% 正文状态：本编导读初稿；尚未正式审核。

% 纲要内容开始
% 
% 纲要内容结束

\BookPart{向量丛、示性类与曲率代表}{g2:part:04}
\BookPartGuide
向量丛在每个小邻域上都是乘积，整体上却可能没有一组处处连续的标架。要比较两张丛，可以直接比较过渡函数，也可以把它们送到同一个参数空间，或计算由它们产生的上同调类。若丛带有联络，还能通过曲率得到微分形式。本编依次建立这些描述之间的联系，使拓扑分类、截面零点和曲率积分可以在同一张丛上相互对照。

\ProseParagraphBreak

\chapref{g2:ch:19}从局部平凡化与拉回出发，证明向量丛随底空间映射的同伦保持同构类型。将各局部标架用单位分解组合，就能把纤维嵌入一个固定线性空间，得到以 Grassmann 空间为目标的分类映射。分类映射的同伦唯一性说明这种表示如何摆脱局部选择。本章同时区分固定秩的分类、添加平凡丛后的稳定分类，以及带联络时需要保留的额外数据。

\ProseParagraphBreak

接下来的构造从每根纤维的方向开始。\chapref{g2:ch:20}把盘与边界球面的相对生成元粘成 Thom 类，再证明与它作杯积给出次数平移的同构。沿零截面拉回产生 Euler 类，把它放回盘球丛对的长正合列则得到 Gysin 序列。截面的横截零点使这些上同调类具有几何代表；局部系数保留绕行时的方向变化。Euler 类所记录的首障碍，与更高维胞腔上可能继续出现的障碍，在具体例子中分别考察。

\ProseParagraphBreak

\chapref{g2:ch:21}继续构造不同次数的示性类。射影丛的上同调关系给出实丛的 Stiefel--Whitney 类和复丛的陈类，实丛的复化再产生 Pontryagin 类。分裂原理把计算化到线丛，同时保证所得等式能回到原来的底空间。Thom 平方及 Wu 公式连接这些类与前面的上同调运算；复 Grassmann 流形的 Schubert 计算则把示性类落实为子空间满足相交条件的计数。

\ProseParagraphBreak

\chapref{g2:ch:22}转向光滑丛及联络。不变多项式作用于曲率后得到闭形式，改变联络时，转移形式给出两份代表之间的恰当差。与拓扑示性类比较以后，曲率的积分才具有确定的归一化和拓扑含义。陈--Simons 形式进一步记录联络之间的差异；三维模型把规范变换、整数模糊和边界条件写成可以核对的公式。

\ProseParagraphBreak

阅读\chapref{g2:ch:19}需要第一卷的向量丛、连续单位分解，以及本卷\chapref{g2:ch:01}的 CW 构造；丛的分类可以从这里单独学习。\chapref{g2:ch:20}、\chapref{g2:ch:21}需要本卷\chapref{g2:ch:04}至\chapref{g2:ch:06}的上同调乘积、局部系数和对偶，零点解释另用\chapref{g2:ch:14}至\chapref{g2:ch:16}的横截性与相交。与首障碍、平方运算有关的比较分别回到\chapref{g2:ch:12}、\chapref{g2:ch:13}。\chapref{g2:ch:22}则需第一卷第18章的联络和曲率、本卷\chapref{g2:ch:07}的 de Rham 比较及本编已建的示性类。计算时宜同时注明工作系数：模二类、整数类和曲率的实上同调类回答不同的问题，其中的整数挠信息应由拓扑构造继续保存。

% !TeX root = ../../Book2.tex
% 纲要标题：## 第19章　向量丛及分类空间
% 纲要位置：几何学计划纲要.md，第 1735 行。
% 正文状态：未撰写；本文件目前是框架。

% 纲要内容开始
% 
% > 本章目标与证明范围：在仿紧、可数可平凡化或 CW 等明确条件下建立分类定理。区分有限维 Grassmann 近似、稳定分类和含联络的几何分类。
% 
% 纲要内容结束

\chapter{向量丛及分类空间}
\label{g2:ch:19}
\BookChapterNumber{g2:ch:19}

\begin{chapterabstract}
待撰写
\end{chapterabstract}

\begin{learninggoals}
待撰写
\end{learninggoals}

待撰写

% !TeX root = ../../Book2.tex
% 纲要标题：### 19.1 丛的运算与同伦不变性
% 纲要位置：几何学计划纲要.md，第 1739 行。
% 正文状态：初稿；尚未正式审核。

% 纲要内容开始
% 
% * 过渡函数、拉回和丛同构
% * 区间上的平凡化延拓
% * 同伦拉回丛的同构
% * 球面向量丛的粘合分类
% 
% 纲要内容结束

\section{丛的运算与同伦不变性}
\label{g2:sec:1901}

把圆周上的每一点换成一条直线，既可以得到平凡线丛，也可以得到 Möbius 线丛。纤维维数相同，区别发生在绕行一周后的粘合。分类向量丛首先要确定允许怎样改变局部坐标，以及底空间中的连续变形会不会改变所得的丛。本节把第一卷的光滑丛语言转到连续范畴，再建立分类映射将要使用的同伦不变性。

\subsection{过渡函数、拉回和丛同构}
\label{g2:subsec:1901:01}

以下 \(\mathbb K\) 表示 \(\mathbb R\) 或 \(\mathbb C\)。第一卷第12章第1节中，向量丛的局部平凡化把每根纤维线性识别为同一个向量空间。本章允许一般拓扑基空间，要求平凡化及换架连续；当底空间和过渡函数都光滑时，便恢复那里的光滑向量丛。

\begin{definition}\label{g2:def:topological-vector-bundle}
	设 \(X\) 为拓扑空间。一个秩为 \(n\) 的 \(\mathbb K\)-向量丛是连续映射 \(\pi:E\to X\)，其每根纤维 \(E_x=\pi^{-1}(x)\) 是 \(n\) 维 \(\mathbb K\)-向量空间，并且每点附近存在开集 \(U\) 及同胚
	\[
		\phi_U:E|_U\longrightarrow U\times\mathbb K^n,
	\]
	使第一分量为 \(\pi\)，第二分量在各纤维上线性可逆。这种同胚称为局部平凡化。覆盖恒等映射、逐纤维线性的连续映射称为丛同态；若它是逐纤维可逆的，则称为丛同构。
\end{definition}

逐纤维可逆已经保证逆映射连续：局部表示是连续矩阵 \(A(x)\in\operatorname{GL}_n(\mathbb K)\)，而逆矩阵的各元素由行列式和余子式连续给出。因此，同构要求的不只是逐点存在某个线性同构，而是这些同构随基点连续变化。记同构类的集合为 \(\operatorname{Vect}_n^{\mathbb K}(X)\)，平凡秩 \(n\) 丛记为 \(\varepsilon_X^n=X\times\mathbb K^n\)。

\ProseParagraphBreak

在覆盖 \(\{U_i\}\) 上选平凡化 \(\phi_i\)，坐标列向量满足 \(v_i=g_{ij}(x)v_j\)，其中
\[
	g_{ij}:U_i\cap U_j\longrightarrow\operatorname{GL}_n(\mathbb K),
	\qquad g_{ij}g_{jk}=g_{ik},\quad g_{ii}=I.
\]
这些连续矩阵就是\term[guodu hanshu]{过渡函数}[transition functions]。反过来，满足上述余循环等式的矩阵把各 \(U_i\times\mathbb K^n\) 粘合起来：在重叠处将第 \(j\) 张图的 \((x,v)\) 与第 \(i\) 张图的 \((x,g_{ij}(x)v)\) 识别。余循环等式恰好保证识别的传递性。取商拓扑后，各张图的饱和开集给出局部乘积，因而所得对象确实是向量丛。

\begin{proposition}\label{g2:prop:bundle-transition-isomorphism}
	设 \(E,F\) 是同一基空间上的有限秩 \(\mathbb K\)-向量丛，在共同开覆盖上的过渡函数分别为 \(g_{ij},h_{ij}\)。丛同态 \(E\to F\) 等价于连续矩阵族 \(A_i\)，满足
	\[
		A_i g_{ij}=h_{ij}A_j.
	\]
	同秩时，它是丛同构当且仅当各 \(A_i(x)\) 可逆。特别地，将局部坐标改为 \(v'_i=b_i(x)v_i\)，过渡函数变为 \(g'_{ij}=b_i g_{ij}b_j^{-1}\)。
\end{proposition}
\begin{proof}
	局部同态写成 \((x,v_i)\mapsto(x,A_i(x)v_i)\)。同一个向量也可用第 \(j\) 张图表示；两种输出相同恰好要求 \(A_ig_{ij}v_j=h_{ij}A_jv_j\) 对所有 \(v_j\) 成立。这就是所列等式，它保证局部连续映射粘成整体同态。逆矩阵连续说明同构判据，最后的换架公式由 \(v'_i=b_ig_{ij}b_j^{-1}v'_j\) 得到。
\end{proof}

丛的运算由纤维运算和过渡函数共同确定。例如 \(E\oplus F\)、\(E\otimes F\) 和 \(E^*\) 的过渡矩阵依次为
\[
	\begin{pmatrix}g_{ij}&0\\0&h_{ij}\end{pmatrix},
	\qquad g_{ij}\otimes h_{ij},
	\qquad(g_{ij}^{-1})^{\mathsf T}.
\]
这里复丛的 \(E^*\) 取复线性对偶，所以出现转置而不是共轭转置。矩阵乘法、张量积及逆转置保留余循环等式，故这些纤维运算都得到向量丛；同样有 \(\operatorname{Hom}(E,F)=E^*\otimes F\) 与外幂 \(\bigwedge^kE\)。

\ProseParagraphBreak

对连续映射 \(f:Y\to X\)，定义\term[lahui cong]{拉回丛}[pullback bundle]
\[
	f^*E=\{(y,e)\in Y\times E:f(y)=\pi(e)\},\qquad (y,e)\longmapsto y.
\]
其局部平凡化来自 \(E\) 在 \(U_i\) 上的平凡化，过渡函数为 \(g_{ij}\circ f\)。下面的方块既交换，也刻画了这个构造的泛性质：
\[
	\begin{tikzcd}
		f^*E\arrow[r]\arrow[d]&E\arrow[d,"\pi"]\\
		Y\arrow[r,"f"']&X.
	\end{tikzcd}
\]
具体说，若 \(q:F\to Y\) 是另一向量丛，\(\Phi:F\to E\) 是覆盖 \(f\) 的逐纤维线性连续映射，则唯一的分解为 \(v\mapsto(q(v),\Phi(v))\in f^*E\)。连续性来自积与子空间拓扑。由此得到自然同构 \((f\circ h)^*E\cong h^*(f^*E)\)，以及拉回与直和、张量积、对偶交换；逐点写下对应的有序组或张量即可核对它们与过渡函数相容。

\subsection{区间上的平凡化延拓}
\label{g2:subsec:1901:02}

为比较同伦的两端，先考察 \(E\to X\times I\)，其中 \(I=[0,1]\)。局部坐标能把相邻时间的纤维对应起来，但一张坐标通常覆盖不了整个区间。区间的紧性提供有限次接续；接缝处需要调整坐标，使两段的对应真正一致。

\begin{lemma}\label{g2:lem:bundle-cylinder-local-trivialization}
	设 \(E\to X\times I\) 为有限秩实或复向量丛。对每个 \(x\in X\)，存在开邻域 \(U\)，使 \(E|_{U\times I}\) 平凡。若已给定 \(E|_{U\times\{0\}}\) 的平凡化，可以在缩小 \(U\) 后将其延拓到整条柱体。
\end{lemma}
\begin{proof}
	沿紧集 \(\{x\}\times I\) 取有限个乘积邻域上的平凡化。取这有限个基空间邻域的交，并对时间作细分
	\[
		0=t_0<t_1<\cdots<t_m=1,
	\]
	使每个 \(U\times[t_{j-1},t_j]\) 都落在其中一张平凡化内。这样的细分存在，因为有限个时间开区间覆盖紧区间，具有正的 Lebesgue 数。

	设前一段已使用坐标 \(\psi_{j-1}\)，后一段原来的坐标为 \(\phi_j\)。在接缝 \(t=t_{j-1}\) 上，二者相差一个连续可逆矩阵 \(B_j(y)\)：\(\psi_{j-1}(y,t_{j-1})=B_j(y)\phi_j(y,t_{j-1})\)。将后一段坐标改为 \(\psi_j(y,t)=B_j(y)\phi_j(y,t)\)，就在接缝处相同。有限个闭时间段的粘合引理使这些坐标拼成一个连续平凡化；逆坐标也按同样方式粘合，因此确为同胚。在第一段先乘一个只依赖 \(y\) 的可逆矩阵，就能匹配任意指定的起点平凡化，再依次接续。
\end{proof}

固定这样的柱体平凡化，保持坐标列向量不变，便得到局部运输
\[
	T_U(y;s,t):E_{(y,s)}\longrightarrow E_{(y,t)},
	\qquad 0\le s,t\le1.
\]
它连续依赖 \((y,s,t)\)，并满足 \(T_U(y;t,t)=I\) 和 \(T_U(y;t,r)T_U(y;s,t)=T_U(y;s,r)\)。不同柱体平凡化所得运输可以不同。下一步用单位分解决定每次使用哪一张图及使用多长的时间，使这些局部运输有次序地接成整体同构。

\subsection{同伦拉回丛的同构}
\label{g2:subsec:1901:03}

此处使用第一卷第4章第4节的连续单位分解定理：仿紧 Hausdorff 空间的任意开覆盖都有从属连续单位分解 \(\{\rho_i\}\)，支集族局部有限，\(\operatorname{supp}\rho_i\subset U_i\)，并且 \(\sum_i\rho_i=1\)。局部有限性使每一点附近的求和实际只有有限项。下面用这些函数分配时间，而不把可逆矩阵作平均，因为可逆矩阵的平均未必仍可逆。

\begin{theorem}[丛的同伦不变性]\label{g2:thm:bundle-homotopy-invariance}
	设 \(X\) 为仿紧 Hausdorff 空间，\(E\to X\times I\) 为有限秩实或复向量丛，\(E_0=E|_{X\times\{0\}}\)。则存在覆盖 \(X\times I\) 恒等映射的丛同构
	\[
		\operatorname{pr}_X^*E_0\cong E,
	\]
	它在 \(t=0\) 上为恒等。因此 \(E_0\cong E_1\)。若 \(f_0,f_1:X\to Y\) 连续同伦，\(V\to Y\) 为有限秩向量丛，则 \(f_0^*V\cong f_1^*V\)。
\end{theorem}
\begin{proof}
	由\cref{g2:lem:bundle-cylinder-local-trivialization}取开覆盖 \(\{U_i\}_{i\in A}\)，使各 \(U_i\times I\) 上都有平凡化及运输 \(T_i\)。取从属的局部有限单位分解 \(\rho_i\)，再给指标集选一个良序。良序只用来固定有限个实际运输的先后，最终同构不要求与这项选择无关。

	令
	\[
		a_i(x)=\sum_{j<i}\rho_j(x),\qquad b_i(x)=a_i(x)+\rho_i(x).
	\]
	这些函数连续，因为每点附近只有有限个支集。对目标时间 \(t\in I\)，第 \(i\) 次运输从高度 \(t a_i(x)\) 到 \(t b_i(x)\)。当 \(x\in U_i\) 时，使用
	\[
		L_i(x,t)=T_i\bigl(x;t a_i(x),t b_i(x)\bigr).
	\]
	当 \(x\notin\operatorname{supp}\rho_i\) 时，两高度相同，令 \(L_i\) 为该纤维上的恒等映射。两种规定在交叠处一致；支集包含于 \(U_i\)，所以它们合成两个图拉回丛之间的连续同构。这个支集条件同时保证，在需要使用局部运输的地方，所选坐标始终有定义。

	现在固定 \(x\)。只有有限个 \(\rho_i(x)\) 非零，按既定顺序复合相应 \(L_i(x,t)\)，得到
	\[
		L(x,t):E_{(x,0)}\longrightarrow E_{(x,t)}.
	\]
	这些时间段首尾相接，因为累积和从零增加到一；零权重给出的段只是恒等。要证明 \(L\) 连续，取 \(x\) 的一个只遇到有限多个支集的邻域 \(W\)。在 \(W\times I\) 上，将这有限个指标全部按顺序列入复合，包括可能权重为零的项，便得到一个固定的有限连续复合。它与逐点定义一致，所以 \(L\) 连续。反向复合这些运输给出连续逆，故 \(L\) 是所需丛同构。\(t=0\) 时每一段都是恒等。

	若 \(H:X\times I\to Y\) 是从 \(f_0\) 到 \(f_1\) 的同伦，将前一结论用于 \(H^*V\)。它的两端分别是 \(f_0^*V\)、\(f_1^*V\)，故得到最后的断言。
\end{proof}

这项构造处理了非紧底空间：不同点可以需要不同数量、不同位置的时间坐标图；局部有限性保证附近的搬运仍由有限个连续步骤给出。同构依赖平凡化及单位分解，定理确定的是同构类。

\begin{corollary}\label{g2:cor:contractible-base-trivial-bundle}
	仿紧 Hausdorff 可缩空间上的每个有限秩实或复向量丛都平凡。更一般地，若 \(f:X\to Y\) 是仿紧 Hausdorff 空间之间的同伦等价，则拉回给出 \(\operatorname{Vect}_n^{\mathbb K}(Y)\to\operatorname{Vect}_n^{\mathbb K}(X)\) 的双射。
\end{corollary}
\begin{proof}
	可缩时，恒等映射与某个常值映射 \(c_{x_0}\) 同伦，而 \(c_{x_0}^*E=X\times E_{x_0}\)，故由定理得到平凡性。一般地，若 \(g\) 为 \(f\) 的同伦逆，则 \(g^*f^*=(f\circ g)^*\) 与 \(f^*g^*=(g\circ f)^*\) 在同构类上分别为恒等，故互为逆映射。
\end{proof}

例如，圆环 \(S^1\times I\) 上的向量丛都来自核心圆周；从内圆周限制和从外圆周限制所得的丛同构。这样的同构并不规定绕圆周的平行移动，后者还需要联络。

\begin{corollary}\label{g2:cor:principal-bundle-homotopy}
	设 \(X\) 为仿紧 Hausdorff 空间，\(G\) 为拓扑群，\(P\to X\times I\) 为局部平凡的右主 \(G\) 丛。则 \(P\cong\operatorname{pr}_X^*(P|_{X\times\{0\}})\)，且同构可在零端取为恒等。因而同伦映射拉回主 \(G\) 丛得到同构主丛。
\end{corollary}
\begin{proof}
	右主丛的局部模型是 \(U\times G\)，右作用为第二分量的右乘；换平凡化在第二分量上左乘一个连续的群值函数。因而柱体平凡化仍可逐段调整并粘合，保持第二坐标的局部运输是右 \(G\) 等变的。上一定理只复合这些运输，没有使用纤维加法。以相同的时间分配与局部有限复合构造同构，它自动右 \(G\) 等变。
\end{proof}

\subsection{球面向量丛的粘合分类}
\label{g2:subsec:1901:04}

把 \(S^m\)、\(m\ge1\) 分成两个闭半球 \(D_+^m,D_-^m\)，交集为赤道 \(S^{m-1}\)。半球可缩，其上的向量丛平凡；全部粘合信息因而集中到赤道上的一族可逆线性变换。

\begin{definition}\label{g2:def:clutching-bundle}
	对连续映射 \(g:S^{m-1}\to\operatorname{GL}_n(\mathbb K)\)，将两个平凡丛 \(D_+^m\times\mathbb K^n\)、\(D_-^m\times\mathbb K^n\) 按
	\[
		(x,v)_+\sim(x,g(x)v)_-\qquad(x\in S^{m-1})
	\]
	粘合，所得丛记为 \(E_g\to S^m\)。映射 \(g\) 称为这个构造的\term[clutching hanshu]{粘合函数}[clutching function]。
\end{definition}

赤道附近的局部平凡性可以直接看见。取领圈 \(S^{m-1}\times(-\varepsilon,\varepsilon)\)，在正半侧保持原坐标，在负半侧乘 \(g(x)^{-1}\)；两侧在赤道上给出相同的坐标。因而商空间在每个赤道点附近仍为乘积，远离赤道则沿用半球内的平凡化。

\begin{theorem}[球面上的粘合分类]\label{g2:thm:clutching-classification}
	设 \(m,n\ge1\)。每个 \(S^m\) 上的秩 \(n\) 实或复向量丛同构于某个 \(E_g\)。两粘合函数 \(g,h\) 给出同构丛，当且仅当存在连续映射
	\[
		A_\pm:D_\pm^m\longrightarrow\operatorname{GL}_n(\mathbb K)
	\]
	使赤道上有 \(h=A_-gA_+^{-1}\)。同伦的粘合函数给出同构丛。

	复丛的同构类因而与 \([S^{m-1},\operatorname{GL}_n(\mathbb C)]\) 双射；若只考虑带指定取向的实丛及保向同构，则相应地与 \([S^{m-1},\operatorname{GL}_n^+(\mathbb R)]\) 双射。
\end{theorem}
\begin{proof}
	由\cref{g2:cor:contractible-base-trivial-bundle}，给定丛在两个半球上都可平凡化；沿赤道比较这两套坐标就得到 \(g\)，原丛由同样的识别关系恢复。若有覆盖球面恒等映射的同构 \(E_g\to E_h\)，在两半球内分别表示为 \(A_+,A_-\)。赤道上的同一点先粘合再映射或先映射再粘合，给出
	\[
		A_-(x)g(x)v=h(x)A_+(x)v.
	\]
	这就是所述关系。反过来，满足关系的两半球同构在接缝处相容，因而粘成整体同构。

	若 \(g_t\) 是粘合函数间的同伦，就在赤道乘区间上用 \(g_t\) 粘合两个 \(D_\pm^m\times I\times\mathbb K^n\)。刚才的领圈坐标同时依赖 \(t\)，所以得到 \(S^m\times I\) 上的向量丛，其两端为 \(E_{g_0},E_{g_1}\)。\cref{g2:thm:bundle-homotopy-invariance}说明两端同构。

	最后，\(\operatorname{GL}_n(\mathbb C)\) 与 \(\operatorname{GL}_n^+(\mathbb R)\) 都道路连通。这里可用极分解先将可逆矩阵连续变为酉矩阵或正向正交矩阵，再由酉矩阵的特征角、实正交矩阵的二维旋转块把它连到单位矩阵。\(A_\pm\) 可沿半球收缩同伦到一个常值，而这个常值再连到单位元。因此 \(h=A_-gA_+^{-1}\) 蕴含 \(h\simeq g\)。带取向时先在两半球选保向标架，所有矩阵都落在正行列式群中。于是上述满射在相应同伦类上也单射。
\end{proof}

常用计算还可以把群缩小。极分解给出形变收缩
\[
	A\longmapsto A(A^*A)^{-s/2},\qquad 0\le s\le1,
\]
复数情形终点在 \(\operatorname U(n)\)，实数情形终点在 \(\operatorname O(n)\)，正行列式分支则在 \(\operatorname{SO}(n)\)。正定矩阵的实幂由正交或酉对角化定义，并连续依赖矩阵；在任一固定紧谱区间上用多项式一致逼近相应正实函数，也能直接验证这种连续性。因此上述球面分类可以分别使用 \(\operatorname U(n)\) 和 \(\operatorname{SO}(n)\)。对 \(m\ge2\)，归一化映射 \(g(x)\mapsto g(x_0)^{-1}g(x)\) 将自由同伦类识别为这些连通群的 \(\pi_{m-1}\)：给定自由同伦 \(H(x,t)\)，乘上 \(H(x_0,t)^{-1}\) 就得到定基点同伦。

\ProseParagraphBreak

未指定取向的实丛需要保留换架关系中的两个分支。当 \(m\ge2\) 时，赤道连通，\(\det g\) 符号恒定；必要时在一个半球改用反向标架，可使 \(g\) 落在正行列式群中。因此每个这样的实丛都可定向。遗忘取向后，还须将 \(g\) 与 \(JgJ^{-1}\) 识别，其中 \(J\) 是一个固定反射矩阵。确实，两半球同时反向换架改变了总取向；反之，在 \(h,g\) 均正行列式的等式 \(h=A_-gA_+^{-1}\) 中，\(A_-\) 与 \(A_+\) 的行列式符号必须相同，故两者可同时同伦到 \(I\) 或同时同伦到 \(J\)。普通实丛的分类便是相应 \(\pi_{m-1}\operatorname{SO}(n)\) 对这个反射共轭作用的轨道集。

\ProseParagraphBreak

\(m=1\) 时，赤道 \(S^0\) 有两个点，不能用连通性将二者的行列式符号合并。沿两个闭弧平凡化后，可以把一个接缝矩阵化为单位矩阵，另一个接缝留下一个矩阵 \(B\)。\(\det B>0\) 时实丛平凡，\(\det B<0\) 时同构于 Möbius 线丛与 \(n-1\) 条平凡线的直和：把 \(B\) 分别沿其所在连通分支变到 \(I\) 或 \(\operatorname{diag}(-1,1,\ldots,1)\) 即得。复矩阵群只有一个连通分支，所以圆周上的复向量丛全部平凡。这个一维例子显示，换架的连通分支本身就是分类数据。

% !TeX root = ../../Book2.tex
% 纲要标题：### 19.2 通用丛和 Grassmann 模型
% 纲要位置：几何学计划纲要.md，第 1746 行。
% 正文状态：初稿；尚未正式审核。

% 纲要内容开始
% 
% * 有限维 Grassmann 典范丛
% * 由截面构造分类映射
% * 维数增大与相容嵌入
% * 无限 Grassmann 空间的拓扑
% 
% 纲要内容结束

\section{通用丛和 Grassmann 模型}
\label{g2:sec:1902}

若能把一张秩 \(n\) 向量丛的每根纤维都线性嵌入同一个 \(\mathbb K^N\)，那么每个基点就确定一个 \(n\) 维子空间。原来的丛因而可以用一张到“所有 \(n\) 维子空间组成的空间”的映射来描述。这里 \(\mathbb K\) 表示 \(\mathbb R\) 或 \(\mathbb C\)，秩始终按 \(\mathbb K\) 计算。本节把这一直观构造写成局部坐标，并说明当不同丛需要不同的 \(N\) 时，如何把这些模型放在同一个空间中。

\subsection{有限维 Grassmann 典范丛}
\label{g2:subsec:1902:01}
第1章已用正交投影矩阵定义 Grassmann 空间，并在\cref{g2:thm:grassmann-schubert-cw}中构造其胞腔分解。这里记为
\[
	\operatorname{Gr}_n(\mathbb K^N)=\{V\subseteq\mathbb K^N:\dim_{\mathbb K}V=n\},\qquad 1\le n\le N.
\]
拓扑由 \(V\mapsto P_V\) 给出，其中 \(P_V\) 是标准内积下的正交投影。矩阵条件 \(P=P^*\)、\(P^2=P\)、\(\operatorname{tr}P=n\) 是闭条件，且 \(\|P\|\le1\)，所以该空间紧且 Hausdorff。现在需要一种能同时描述子空间及其内部向量的坐标。

\ProseParagraphBreak

固定正交分解 \(\mathbb K^N=A\oplus A^\perp\)，其中 \(\dim A=n\)。令 \(U_A\) 为投影 \(\pi_A|_V:V\to A\) 可逆的那些 \(V\)。这样的子空间唯一写成线性映射 \(L:A\to A^\perp\) 的图：
\[
	V=\operatorname{graph}L=\{a+L(a):a\in A\}.
\]
设 \(B_L\) 表示矩阵 \(\binom{I}{L}\)，则
\[
	P_V=B_L(I+L^*L)^{-1}B_L^*.
\]
反过来，若按 \(A\oplus A^\perp\) 把 \(P_V\) 分块，则 \(L=P_{21}P_{11}^{-1}\)。这些公式证明图参数与投影矩阵相互连续，且在有限维矩阵变量中光滑。可逆条件是开条件，故 \(U_A\) 开。任意子空间 \(V\) 都属于 \(U_V\)，这些图坐标覆盖整个空间，坐标变换也是由矩阵乘法及可逆矩阵求逆组成。因此实 Grassmann 空间的实维数为 \(n(N-n)\)，复 Grassmann 空间的复维数为 \(n(N-n)\)。

\begin{definition}\label{g2:def:finite-tautological-bundle}
	对 \(\mathbb K=\mathbb R\) 或 \(\mathbb C\) 及 \(1\le n\le N\)，定义\term[dianfan-xiangliangcong]{典范向量丛}[tautological vector bundle]
	\[
		\gamma_n^N=\{(V,v)\in\operatorname{Gr}_n(\mathbb K^N)\times\mathbb K^N:v\in V\},\qquad \pi(V,v)=V,
	\]
	总空间取积空间的子空间拓扑。其在 \(V\) 上的纤维就是 \(V\) 本身。
\end{definition}

它确实是秩 \(n\) 向量丛，而不只是附带线性纤维的映射。在 \(U_A\) 上，局部平凡化及逆为
\[
	(V,v)\longmapsto(V,\pi_Av),\qquad
	(V,a)\longmapsto(V,a+L_Va).
\]
二者连续且逐纤维线性。特别地，选定 \(A\) 的一组基 \(a_1,\ldots,a_n\) 后，\(a_j+L_Va_j\) 就给出 \(U_A\) 上的一组连续局部标架。\(n=1\) 时，基底分别是 \(\mathbb RP^{N-1}\)、\(\mathbb CP^{N-1}\)，所得典范线丛正是前面射影空间中的模型。

\ProseParagraphBreak

一张映射 \(c:X\to\operatorname{Gr}_n(\mathbb K^N)\) 的拉回丛可以直接写成
\[
	c^*\gamma_n^N=\{(x,v)\in X\times\mathbb K^N:v\in c(x)\}.
\]
因此有限维典范丛统一描述了平凡丛 \(X\times\mathbb K^N\) 的所有秩 \(n\) 子丛。这里的 \(N\) 控制可容纳的子丛；下一步证明紧基底上的每张有限秩丛都能放入某个这样的平凡丛。

\subsection{由截面构造分类映射}
\label{g2:subsec:1902:02}
一张非平凡丛通常没有一组处处独立的全局标架，但可以有更多全局截面，在每一点共同张成纤维。多出来的截面允许我们先使用局部标架，再把每个标架在其定义域以外补成零。单位分解恰好提供所需权重。

\begin{proposition}\label{g2:prop:bundle-embedding-classifier}
	设 \(X\) 为仿紧 Hausdorff 空间，\(E\to X\) 为秩 \(n\) 的实或复向量丛，并有有限平凡化覆盖 \(U_1,\ldots,U_m\)。则存在逐纤维线性的丛嵌入
	\[
		J:E\longrightarrow X\times\mathbb K^{mn}
	\]
	和连续映射 \(c_J:X\to\operatorname{Gr}_n(\mathbb K^{mn})\)，满足 \(c_J(x)=J(E_x)\) 及 \(E\cong c_J^*\gamma_n^{mn}\)。特别地，紧 Hausdorff 基底上的每张有限秩向量丛都有这样的有限维模型。
\end{proposition}
\begin{proof}
	取从属连续单位分解 \(\rho_1,\ldots,\rho_m\)，满足 \(\operatorname{supp}\rho_i\subseteq U_i\)。所用的是第一卷第4.4节的连续单位分解定理，其条件正是仿紧与 Hausdorff。将各局部内积加权求和，得到 \(E\) 上的连续正定内积。对每个 \(U_i\) 的局部标架作逐向量正交化，得到正交单位标架 \(e_{i1},\ldots,e_{in}\)；归一化的分母始终为正，故标架连续。

	令 \(s_{ia}=\sqrt{\rho_i}\,e_{ia}\) 于 \(U_i\)，在其余地方补零。支集闭且包含于 \(U_i\)，所以每个不属于 \(U_i\) 的点都有一个邻域使该截面恒为零，补零连续。它们定义满射
	\[
		A_x:\mathbb K^{mn}\longrightarrow E_x,\qquad
		(z_{ia})\longmapsto\sum_{i,a}z_{ia}s_{ia}(x).
	\]
	至少一个 \(\rho_i(x)>0\)，对应的 \(n\) 个截面已张成 \(E_x\)，所以确实满射。令 \(\tau_i(v)\in\mathbb K^n\) 为 \(v\in E_x\) 在第 \(i\) 组正交标架下的坐标。伴随映射给出
	\begin{equation}\label{g2:eq:weighted-bundle-embedding}
		J_x(v)=A_x^*(v)
		=\bigl(\sqrt{\rho_1(x)}\,\tau_1(v),\ldots,
		\sqrt{\rho_m(x)}\,\tau_m(v)\bigr).
	\end{equation}
	不在 \(U_i\) 时相应坐标仍取零。正交性与 \(\sum_i\rho_i=1\) 给出
	\[
		\|J_x(v)\|^2=\sum_i\rho_i(x)\|v\|^2=\|v\|^2.
	\]
	故 \(J_x\) 是等距单射，且 \(A_xJ_x=I_{E_x}\)。其像的正交投影为 \(P_x=J_xA_x\)，矩阵元连续依赖于 \(x\)，所以 \(c_J(x)=\operatorname{im}P_x\) 连续。

	映射 \(v\mapsto(x,J_xv)\) 是 \(E\) 到 \(c_J^*\gamma_n^{mn}\) 的逐纤维线性同构，逆为 \((x,z)\mapsto A_xz\)，故是丛同构。这给出拉回图
	\[
		\begin{tikzcd}[column sep=large]
			E\arrow[r,"{v\mapsto(c_J(x),J_xv)}"]\arrow[d]&\gamma_n^{mn}\arrow[d]\\
			X\arrow[r,"c_J"']&\operatorname{Gr}_n(\mathbb K^{mn}).
		\end{tikzcd}
	\]
	最后，紧 Hausdorff 空间仿紧，且任意局部平凡化覆盖有有限子覆盖，所以满足已用条件。
\end{proof}

这个构造也解释了映射为何取值于 Grassmann 空间：局部标架经加权提供了多个坐标，真正与原纤维对应的是这些坐标共同张成的 \(n\) 维像空间。\term[fenlei-yingshe]{分类映射}[classifying map]是给定一个拉回同构 \(E\cong c^*\gamma\) 的这种映射；不同权重与标架可能给出不同的 \(c\)。下一节将证明，在统一的无限模型中，它们具有同一个同伦类。

\subsection{维数增大与相容嵌入}
\label{g2:subsec:1902:03}
增大环境维数时，把 \(\mathbb K^N\) 嵌入 \(\mathbb K^{N+1}\) 的前 \(N\) 个坐标，便得到闭嵌入
\[
	i_N:\operatorname{Gr}_n(\mathbb K^N)\longrightarrow\operatorname{Gr}_n(\mathbb K^{N+1}),\qquad V\longmapsto V\oplus0.
\]
投影矩阵只是在右下角添零，所以它连续；紧源到 Hausdorff 目标的单射给出嵌入。典范丛相容地满足 \(i_N^*\gamma_n^{N+1}\cong\gamma_n^N\)，同构在纤维上就是 \(v\mapsto(v,0)\)。因此把一个分类映射移入更大的环境，不会改变它表示的秩 \(n\) 丛。

\ProseParagraphBreak

另一个有用构造是正交补丛。在 \(\operatorname{Gr}_n(\mathbb K^N)\) 上令
\[
	(\gamma_n^N)^\perp=\{(V,w):w\in V^\perp\}.
\]
它由连续投影矩阵 \(I-P_V\) 描述，秩为 \(N-n\)，图坐标同样给出局部平凡化。逐纤维相加得到
\begin{equation}\label{g2:eq:tautological-orthogonal-sum}
	\gamma_n^N\oplus(\gamma_n^N)^\perp
	\cong\operatorname{Gr}_n(\mathbb K^N)\times\mathbb K^N.
\end{equation}
拉回这个同构，便知满足\cref{g2:prop:bundle-embedding-classifier}条件的 \(E\) 有某张补丛 \(F\)，使 \(E\oplus F\) 平凡。补丛本身可以非平凡。

\begin{example}\label{g2:exm:sphere-tangent-stable-triviality}
	把单位球面 \(S^m\) 嵌入 \(\mathbb R^{m+1}\)，其切空间为 \(x^\perp\)。于是 \(x\mapsto x^\perp\) 将 \(\mathrm TS^m\) 表示为 \(\gamma_m^{m+1}\) 的拉回。正交补直线由处处非零的法向量 \(x\) 张成，故为平凡线丛。具体同构为
	\[
		\mathrm TS^m\oplus\varepsilon^1\longrightarrow\varepsilon^{m+1},
		\qquad(x,v,t)\longmapsto(x,v+tx),
	\]
	其中 \(\varepsilon^r=S^m\times\mathbb R^r\)。逆映射将 \(z\) 分解为 \(z-\langle z,x\rangle x\) 与 \(\langle z,x\rangle\)。

	然而 \(\mathrm TS^2\) 不平凡。若有连续非零切向量场，将它归一化为 \(v(x)\)，则
	\[
		H(x,t)=\cos(\pi t)x+\sin(\pi t)v(x)
	\]
	始终为单位向量，会给出恒等映射与对径映射之间的同伦。两者在 \(S^2\) 上的度分别为 \(1\)、\(-1\)，与\cref{g2:prop:degree-basic,g2:prop:reflection-antipodal-degree}矛盾。因此加一条平凡线后得到的同构，并不保证原秩下也有同构。
\end{example}

这里有两种不同的增大：\(N\) 增大只改变装载纤维的环境，\(E\mapsto E\oplus\varepsilon^r\) 则实际增加纤维的秩。前者用于固定秩的分类空间，后者产生稳定分类。例子中的切丛说明，稳定操作可能遗失原秩下的信息。

\subsection{无限 Grassmann 空间的拓扑}
\label{g2:subsec:1902:04}
把 \(\mathbb K^\infty\) 定义为有限支撑序列构成的向量空间，即 \(\bigcup_N\mathbb K^N\)。它的每个有限维子空间都包含于某个 \(\mathbb K^N\)：只须对一组有限基向量的支撑取并。因而所有 \(n\) 维子空间可以由前面的有限模型穷尽。

\begin{definition}\label{g2:def:infinite-grassmann}
	对固定 \(n\ge1\)，定义
	\[
		G_n=\operatorname{Gr}_n(\mathbb K^\infty)
		=\bigcup_{N\ge n}\operatorname{Gr}_n(\mathbb K^N).
	\]
	赋予它\term[zhijixian-ruotuopu]{直极限弱拓扑}[weak direct-limit topology]：集合 \(U\subseteq G_n\) 开，当且仅当每个 \(U\cap\operatorname{Gr}_n(\mathbb K^N)\) 都在该有限模型中开。同样，令 \(\gamma_n=\bigcup_N\gamma_n^N\)，总空间也取直极限弱拓扑。
\end{definition}

这些弱拓扑的作用是把“由各有限阶段定义的映射”粘合成连续映射。有限阶段的原有拓扑保留不变：一个有限阶段中的闭集，在所有更大的阶段中仍闭，所以在整个直极限中闭。投影矩阵的各个坐标函数在所有阶段连续，故在 \(G_n\) 上连续，并能分离不同子空间，因此 \(G_n\) 是 Hausdorff 空间。第1章的 Schubert 胞腔与上述嵌入相容：在 \(\mathbb K^N\) 中的胞腔恰好是最大枢轴指标不超过 \(N\) 的那些胞腔，闭包关系保持不变。所以这些有限子复形取并也给出 \(G_n\) 的 CW 结构，其弱拓扑就是这里规定的拓扑。

\begin{lemma}\label{g2:lem:grassmann-compact-stage}
	固定 \(n\ge1\) 及 \(\mathbb K=\mathbb R\) 或 \(\mathbb C\)。在上述弱拓扑下，\(\operatorname{Gr}_n(\mathbb K^\infty)\) 的每个紧子集都包含于某个有限阶段 \(\operatorname{Gr}_n(\mathbb K^N)\)。因此紧空间上的映射及其同伦都可在足够大的有限阶段中进行。
\end{lemma}
\begin{proof}
	若紧集 \(K\) 不包含于任一阶段，可递归选择点 \(x_j\in K\) 与严格递增的整数 \(N_j\)，使第 \(N_j\) 阶段包含全部先前选过的点，而 \(x_j\) 不在该阶段。于是点集 \(A=\{x_j:j\ge1\}\) 与每个有限阶段都只有有限交。对任意子集 \(A'\subseteq A\)，该交仍有限，因有限阶段 Hausdorff而闭。弱拓扑遂说明 \(A'\) 在 \(G_n\) 中闭。特别地，\(A\) 是 \(K\) 的闭子集，且其每个子集都闭，所以 \(A\) 是紧的无限离散空间。这与单点开覆盖没有有限子覆盖矛盾。

	连续映射的紧像满足第一项。若定义域 \(X\) 紧，则 \(X\times I\) 也紧，将第一项用于同伦映射即可。
\end{proof}

判断一张从弱直极限出发的同伦是否连续，还需要同时处理时间参数。下面把这一点明确下来，后面的坐标移动将反复使用它。

\begin{lemma}\label{g2:lem:grassmann-parameter-continuity}
	设 \(Y=\bigcup_NY_N\) 取由各 \(Y_N\) 规定的直极限拓扑，\(L\) 为局部紧 Hausdorff 空间。映射 \(F:Y\times L\to Z\) 连续，当且仅当各限制 \(F|_{Y_N\times L}\) 连续。特别地，这适用于 \(Y=G_n\)、\(L=I\)，以及局部丛坐标中的 \(L=\mathbb K^n\)。
\end{lemma}
\begin{proof}
	必要性直接来自限制映射。反过来，给连续映射集 \(C(L,Z)\) 取紧开拓扑，其子基本开集为
	\[
		[C,O]=\{u:u(C)\subseteq O\},\qquad C\subseteq L\ \text{紧},\quad O\subseteq Z\ \text{开}.
	\]
	每个连续映射 \(F_N:Y_N\times L\to Z\) 的伴随映射 \(y\mapsto F_N(y,\cdot)\) 连续：若 \(F_N(\{y\}\times C)\subseteq O\)，紧性允许把 \(C\) 上的有限个乘积邻域合并，得到 \(y\) 的一个邻域，使同样的包含对该邻域内全部点成立。由 \(Y\) 的终拓扑性质，这些伴随映射合成连续映射 \(Y\to C(L,Z)\)。

	局部紧 Hausdorff 条件保证求值 \(C(L,Z)\times L\to Z\) 连续。具体地，若 \(u(t)\in O\)，可取 \(t\) 的开邻域 \(V\)，使 \(\overline V\) 紧且包含于 \(u^{-1}(O)\)；于是 \([\overline V,O]\times V\) 的求值落在 \(O\) 内。将求值与前述伴随映射复合，便得到 \(F\) 的连续性。
\end{proof}

现在验证无限典范丛仍然局部平凡。对任意有限坐标集 \(J\subset\mathbb N\)、\(|J|=n\)，记 \(A_J\) 为相应的坐标 \(n\) 平面，\(U_J\subseteq G_n\) 为 \(\pi_{A_J}|_V\) 可逆的子空间集合。它与各有限阶段的交都是开集，所以 \(U_J\) 开。任一 \(n\) 个线性无关向量构成的矩阵都有一个非零 \(n\) 阶子式，故这些 \(U_J\) 覆盖 \(G_n\)。有限阶段的平凡化
\[
	(V,v)\longmapsto(V,\pi_{A_J}v)
\]
相容地取并，给出 \(\gamma_n|_{U_J}\to U_J\times A_J\)。它的连续性由总空间的弱拓扑得到，逆映射在每个有限阶段乘 \(A_J\) 上连续，再由\cref{g2:lem:grassmann-parameter-continuity}得到整体连续。这里 \(U_J\) 本身取其与各阶段交所给的弱拓扑；这是开子空间上的终拓扑。于是 \(\gamma_n\) 是一张真正的秩 \(n\) 向量丛，且具有可数平凡化覆盖。

\begin{lemma}\label{g2:lem:universal-numerable-cover}
	无限典范丛 \(\gamma_n\to G_n\) 的可数坐标平凡化覆盖 \(\{U_J\}\) 有从属的局部有限连续单位分解。此外，在每个有限阶段中，只有有限个权重可能非零。
\end{lemma}
\begin{proof}
	枚举全部 \(n\) 元坐标集为 \(J_1,J_2,\ldots\)。对 \(V\in G_n\)，置
	\[
		f_j(V)=\det\bigl(P_V[J_j,J_j]\bigr).
	\]
	这是投影矩阵相应主子式。在一个包含 \(V\) 的有限阶段选正交单位标架矩阵 \(Q\)，则 \(P_V=QQ^*\)，所以 \(f_j(V)=|\det Q_{J_j}|^2\ge0\)。对 \(\det(Q^*Q)\) 作矩阵乘积的子式展开，重复行指标的项由行列式的交替性抵消，剩下每个 \(n\) 元坐标集的两个共轭子式之积，得到
	\[
		\sum_j f_j(V)=\det(Q^*Q)=1.
	\]
	在第 \(N\) 阶段，只有 \(J_j\subseteq\{1,\ldots,N\}\) 时该函数才可能非零；这样的 \(j\) 只有有限个。并且 \(f_j>0\) 的集合正是 \(U_{J_j}\)。

	原函数族的支集未必局部有限，因此先削去每一项的一小部分。令
	\[
		a_j=2^{-j-1},\qquad g_j(V)=\max\{0,f_j(V)-a_j\}.
	\]
	有 \(\operatorname{supp}g_j\subseteq\{f_j\ge a_j\}\subseteq U_{J_j}\)。对固定 \(V\)，取 \(j_0\) 包含所有在 \(V\) 处非零的 \(f_j\)。集合
	\[
		C=\bigcup_{j>j_0}\{f_j\ge a_j\}
	\]
	与每个有限阶段的交只是有限个闭集的并，故 \(C\) 在弱拓扑下闭；且 \(V\notin C\)。开邻域 \(G_n\setminus C\) 不遇到任何 \(j>j_0\) 的支集，证明了局部有限性。

	又因 \(\sum_j a_j=1/2\)，有 \(\sum_jg_j(V)\ge1/2>0\)。局部有限性使这个和连续。于是
	\[
		\rho_j(V)=\frac{g_j(V)}{\sum_i g_i(V)}
	\]
	就是所需单位分解。在每个有限阶段只有有限个 \(f_j\) 非零，相同结论也适用于 \(\rho_j\)。
\end{proof}

最后把有限嵌入构造用于可数局部有限数据。设 \(E\to X\) 有可数平凡化覆盖及从属局部有限单位分解，则式\eqref{g2:eq:weighted-bundle-embedding}按 \(i=1,2,\ldots\) 排列坐标，给出 \(E\to X\times\mathbb K^\infty\) 的逐纤维单射。在每个基点的某个邻域内，只有固定有限个权重可能非零，所以其分类映射局部落在一个固定有限阶段；前面的矩阵公式因此直接保证它连续。若 \(X\) 仿紧 Hausdorff 且有可数平凡化覆盖，先取局部有限细化与从属单位分解，再把属于同一原覆盖成员的权重相加，即得到这样的可数局部有限数据。局部有限支集的并仍闭，故分组后支集仍包含于对应平凡化域。这样，非紧基底也得到无限典范丛的拉回表示。

% !TeX root = ../../Book2.tex
% 纲要标题：### 19.3 分类定理的证明
% 纲要位置：几何学计划纲要.md，第 1753 行。
% 正文状态：未撰写；本文件目前是框架。

% 纲要内容开始
% 
% * 丛到分类映射的构造
% * 同伦类与选择无关性
% * 同伦映射给出同构丛
% * 适用基空间与数值型限制
% 
% 纲要内容结束

\section{分类定理的证明}
\label{g2:sec:1903}

待撰写

\subsection{丛到分类映射的构造}
\label{g2:subsec:1903:01}
待撰写

\subsection{同伦类与选择无关性}
\label{g2:subsec:1903:02}
待撰写

\subsection{同伦映射给出同构丛}
\label{g2:subsec:1903:03}
待撰写

\subsection{适用基空间与数值型限制}
\label{g2:subsec:1903:04}
待撰写
% 纲要条目排为编号小节标题，并进入目录与书签；说明段仍保留于纲要注释中。

% !TeX root = ../../Book2.tex
% 纲要标题：### 19.4 主丛、结构群约化与例子
% 纲要位置：几何学计划纲要.md，第 1760 行。
% 正文状态：未撰写；本文件目前是框架。

% 纲要内容开始
% 
% * 主丛分类和 BG
% * 约化与齐性纤维丛的截面
% * 实线丛、复线丛和低秩球面丛
% * 联络不是纯拓扑分类数据
% 
% ---
% 
% 纲要内容结束

\section{主丛、结构群约化与例子}
\label{g2:sec:1904}

待撰写

\subsection{主丛分类和 BG}
\label{g2:subsec:1904:01}
待撰写

\subsection{约化与齐性纤维丛的截面}
\label{g2:subsec:1904:02}
待撰写

\subsection{实线丛、复线丛和低秩球面丛}
\label{g2:subsec:1904:03}
待撰写

\subsection{联络不是纯拓扑分类数据}
\label{g2:subsec:1904:04}
待撰写
% 纲要条目排为编号小节标题，并进入目录与书签；说明段仍保留于纲要注释中。


\begin{chaptersummary}
待撰写
\end{chaptersummary}

\BookExercises
待撰写

% !TeX root = ../../Book2.tex
% 纲要标题：## 第20章　Thom 同构、Euler 类与 Gysin 序列
% 纲要位置：几何学计划纲要.md，第 1769 行。
% 正文状态：初稿；尚未正式审核。

% 纲要内容开始
% 
% > 本章目标与证明范围：从局部纤维定向构造 Thom 类，在相对上同调中证明同构。随后以零截面与球丛解释 Euler 类和 Gysin，不把这些对象分别当作独立定义表。
% 
% 纲要内容结束

\chapter{Thom 同构、Euler 类与 Gysin 序列}
\label{g2:ch:20}
\BookChapterNumber{g2:ch:20}

\begin{chapterabstract}
	纤维的局部方向可以组成总空间上的相对上同调类。本章从盘球丛、Thom 空间及定向局部系统出发，证明任意 CW 基空间上带局部系数的 Thom 同构，说明有限覆盖粘合与无限胞腔的次数控制如何配合。零截面拉回得到 Euler 类，横截截面的零点流形及孤立零点指标给出它的几何解释，并与前文的球丛传递规范比较。最后由盘球丛对的长正合列推导 Gysin 序列，计算 Hopf 丛和圆丛，再用微分形式建立沿纤维积分及其符号约定。
\end{chapterabstract}

\begin{learninggoals}
	\begin{itemize}
		\item 比较盘球丛对、穿孔总空间对与 Thom 商空间，辨认非紧基底及扭系数时各模型的适用范围。
		\item 用纤维的局部同调描述方向输运，解释模二系数和定向局部系统各自如何处理方向反转。
		\item 完整证明局部乘积与整体 Thom 同构，说明无限 CW 基空间为何由有限个维数层控制固定次数。
		\item 从零截面构造 Euler 类，证明它与首障碍、横截零点类及局部指标和的关系。
		\item 用一个具体三维丛辨认 Euler 类消失与非零截面存在之间的区别。
		\item 从长正合列推导 Gysin 序列，计算球丛上同调，并将其积分映射与微分形式的纤维积分对应。
	\end{itemize}
\end{learninggoals}

球面上的向量场为什么不能总是避开零点？第12章把问题放在逐胞腔延拓中，用边界球面的映射度记录首个障碍；第16—17章则把孤立零点附近的映射度相加，得到 Euler 示性数。两种计算都出现同一种局部整数：绕着一个零点走一圈，向量方向绕纤维球面多少次。本章将这个整数放进一张向量丛的相对上同调，从而把延拓、零点和整体上同调联系起来。

\ProseParagraphBreak

先看平凡丛 \(B\times\mathbb R^r\)。每根纤维都带有圆盘对 \((D^r,S^{r-1})\) 的正生成元 \(u_r\)，于是底空间上的类 \(a\) 可以变成相对类 \(a\times u_r\)，次数升高 \(r\)。若丛有扭转，局部正生成元可能在交叠处变号；选择全局方向，或让系数也按同一符号运输，便能把它们拼成 Thom 类。Thom 同构断言，这个升次数构造已经给出盘球丛对的全部上同调。

\ProseParagraphBreak

随后只需对同一个 Thom 类作两种操作。沿零截面拉回得到 Euler 类，它记录一个截面的零点；把相对上同调放回盘球丛对的长正合列，得到 Gysin\footnote{Werner Gysin（1915--1998），瑞士数学家，研究流形纤维化的同调，提出后来以其名字命名的映射与正合列。}序列，其中一支箭头正是乘 Euler 类。这样，三项构造共享同一个局部生成元，而每个映射都能追溯到一个具体的空间对或截面。

\ProseParagraphBreak

本章主要对象是 CW 复形上的有限秩实向量丛；光滑零点与纤维积分部分另说明紧性、横截性及定向条件。第5章的局部系数、第6章的对偶与帽积、第7章的 de Rham 比较、第12章的首障碍及第19章的丛同伦不变性会在用到时逐项回引。Thom 构造与平方运算的联系可参阅\cite[第23章第5--6节]{May1999}，Euler 类与 Thom 同构的经典组织可参阅\cite[第9节，Oriented Bundles and the Euler Class；第10节，The Thom Isomorphism Theorem]{MilnorStasheff1974}。第21章将继续比较由 Thom 平方产生的类与其他示性类构造。

% !TeX root = ../../Book2.tex
% 纲要标题：### 20.1 Thom 空间及定向
% 纲要位置：几何学计划纲要.md，第 1773 行。
% 正文状态：初稿；尚未正式审核。

% 纲要内容开始
% 
% * 盘丛、球丛与相对模型
% * 局部纤维生成元
% * 系数环及扭定向系统
% 
% 纲要内容结束

\section{Thom 空间及定向}
\label{g2:sec:2001}

一根实 \(r\) 维纤维本身可缩，它的普通上同调看不出维数；把零点以外的部分作为相对项后，便出现一个 \(r\) 维生成元。本章把这个局部生成元沿整张向量丛组织起来。首先需要说明采用哪一个相对空间，以及纤维坐标反转方向时生成元如何变化。

\subsection{盘丛、球丛与相对模型}
\label{g2:subsec:2001:01}
本章设 \(B\) 为 CW 复形，\(p:E\to B\) 为秩 \(r\ge1\) 的实向量丛，零截面记为 \(s_0\)。\cref{g2:lem:cw-paracompact-partitions}保证 \(B\) 仿紧；将局部 Euclid 内积按从属单位分解求和，就得到连续纤维内积。这是\secref{g2:sec:1902}构造丛嵌入时已经使用的方法，局部有限性使它同样适用于无限覆盖。以下固定这样一个内积，并记 \(\|v\|\) 为相应范数。

\begin{definition}\label{g2:def:disk-sphere-thom}
	向量丛 \(E\) 的\term[pancong]{盘丛}[disk bundle]、\term[qiucong]{球丛}[sphere bundle]分别为
	\[
		D(E)=\{v\in E:\|v\|\le1\},\qquad S(E)=\{v\in E:\|v\|=1\}.
	\]
	把整个 \(S(E)\) 压成同一个基点 \(*\)，所得带基点空间
	\[
		\operatorname{Th}(E)=D(E)/S(E)
	\]
	称为 \(E\) 的\term[thom-kongjian]{Thom 空间}[Thom space]。另记 \(E^\times=E\setminus s_0(B)\)。
\end{definition}

在局部正交标架中，盘丛和球丛分别是 \(U\times D^r\) 与 \(U\times S^{r-1}\)。Thom 商把所有这些纤维的边界共同送到 \(*\)；各纤维内部的点仍然保留。所以它同时记住纤维的维数及纤维之间的粘合方式。对基底只有一点的情形，这一构造就是 \(D^r/S^{r-1}\cong S^r\)。

\begin{proposition}\label{g2:prop:thom-relative-models}
	对上述向量丛，\((D(E),S(E))\) 是 NDR 对。包含映射
	\[
		(D(E),S(E))\longrightarrow(E,E^\times)
	\]
	在同调和上同调中诱导同构，系数可以取从 \(B\) 拉回的局部系统。对常系数环 \(R\)，商映射还诱导自然同构
	\begin{equation}\label{g2:eq:thom-relative-quotient}
		\widetilde H^k(\operatorname{Th}(E);R)\xrightarrow{\ \cong\ }H^k(D(E),S(E);R).
	\end{equation}
	约化同调也有相同的商空间比较。
\end{proposition}
\begin{proof}
	先给出控制球面附近的径向同伦。对 \(v\in D(E)\)，令 \(\rho=\|v\|\)，并置
	\[
		u(v)=\min\{1,2(1-\rho)\},\qquad
		H(v,t)=\left(1-t+\frac{t}{\max\{\rho,1/2\}}\right)v.
	\]
	分母处处为正，故包括零截面在内该式连续。其范数始终不超过一；当 \(\rho=1\) 时同伦固定 \(v\)，当 \(u(v)<1\) 时有 \(\rho>1/2\)，从而 \(\|H(v,1)\|=1\)。又 \(u^{-1}(0)=S(E)\)。这逐项满足\cref{g2:def:NDR-pair}的条件，所以球丛包含是\cref{g2:thm:NDR-cofibration}所述的闭余纤维化。

	径向映射 \(v\mapsto v/\max\{1,\|v\|\}\) 给出 \(E\) 到 \(D(E)\) 的强形变收缩；线性插值保持每根纤维，且固定盘丛。同样，\(v\mapsto v/\|v\|\) 给出 \(E^\times\) 到 \(S(E)\) 的强形变收缩。因此包含映射在两个空间和两个相对子空间的同调、上同调上分别诱导同构。把这些映射放入对的长正合列，五引理给出所需的相对同构。所有同伦均保持 \(p(v)\) 不变，故拉回系数的识别也是固定的；局部系数情形使用\cref{g2:thm:twisted-homotopy-invariance}的同伦不变性。

	最后核对商空间比较。取球面附近的开领圈
	\[
		N=\{v\in D(E):\|v\|>1/2\}.
	\]
	上面的同伦限制为 \(N\) 到 \(S(E)\) 的强形变收缩，且在商上把 \(N/S(E)\) 收缩到 \(*\)。诱导同伦连续，是因为商映射与局部紧区间 \(I\) 的乘积仍为商映射，见\cref{g2:lem:quotient-interval}。于是长正合列给出
	\[
		\begin{aligned}
			H^k(D(E),N;R)&\xrightarrow{\ \cong\ }H^k(D(E),S(E);R),\\
			H^k(\operatorname{Th}(E),N/S(E);R)&\cong\widetilde H^k(\operatorname{Th}(E);R).
		\end{aligned}
	\]
	商映射在 \(D(E)\setminus S(E)\) 上是到 \(\operatorname{Th}(E)\setminus\{*\}\) 的同胚。分别从上述两个对中切除 \(S(E)\) 和 \(\{*\}\)，所得对因而同胚。切除条件成立：这两个闭集分别包含于开集 \(N\) 和 \(N/S(E)\) 内。由\cref{g2:prop:relative-cohomology-exactness}的上同调切除，两个对的上同调相同。合并这些同构便得到式\eqref{g2:eq:thom-relative-quotient}，而且合成映射正是商映射诱导的映射。同调的证明逐步相同，所用切除是\cref{g2:thm:singular-excision}。
\end{proof}

盘球丛模型适合使用长正合列，\((E,E^\times)\) 适合讨论零截面附近的几何，而 Thom 空间把相对问题写成一个带基点空间的约化上同调。以上论证说明这些模型为何给出同一份常系数信息。局部系数将在相对模型上保留：从 \(B\) 拉回的局部系统一般不能在压缩整个 \(S(E)\) 后成为商空间上的局部系统。

\ProseParagraphBreak

内积的选择不会改变 Thom 空间的同胚类型。若两组范数为 \(\|\cdot\|_0,\|\cdot\|_1\)，则
\[
	v\longmapsto
	\begin{cases}
		\dfrac{\|v\|_0}{\|v\|_1}v,&v\ne0,\\
		0,&v=0
	\end{cases}
\]
把第一组盘球丛送到第二组盘球丛，交换下标便得逆。在任何局部标架中，两组正定内积在每个基点附近具有统一的范数比较常数，所以该映射在零截面也连续。它因此诱导 Thom 空间的同胚。若 \(B\) 紧，则 \(D(E)\) 紧，且 \(v\mapsto v/(1-\|v\|)\) 将开盘丛同胚到 \(E\)，这时 \(\operatorname{Th}(E)\) 正是总空间 \(E\) 的一点紧化。非紧基底时使用上述盘球丛商定义；\cref{g2:ex:noncompact-thom-not-compactification}将具体比较两种构造。

\subsection{局部纤维生成元}
\label{g2:subsec:2001:02}
固定 \(b\in B\)，记 \(D_b=D(E)\cap E_b\)、\(S_b=S(E)\cap E_b\)。相对球的计算给出
\begin{equation}\label{g2:eq:fiber-orientation-group}
	H_j(D_b,S_b;\mathbb Z)\cong
	\begin{cases}
		\mathbb Z,&j=r,\\
		0,&j\ne r.
	\end{cases}
\end{equation}
这也可由 \(D_b/S_b\cong S^r\) 及\cref{g2:thm:sphere-homology}得到。选一组有序线性坐标，就选出顶维相对基本类 \(o_b=[D_b,S_b]\)；反转坐标方向时它变成 \(-o_b\)。由上一小节的相对比较，它也确定 \(H_r(E_b,E_b\setminus\{0\};\mathbb Z)\) 的一个生成元。

\ProseParagraphBreak

对于任意 Abel 群 \(A\)，上同调万有系数定理给出自然识别
\begin{equation}\label{g2:eq:fiber-top-cohomology-hom}
	H^r(D_b,S_b;A)\cong\operatorname{Hom}_{\mathbb Z}\bigl(H_r(D_b,S_b;\mathbb Z),A\bigr),
\end{equation}
其余次数的相对上同调为零。这里顶维相对同调是自由群，且其余相对同调消失，所以万有系数定理中的 \(\operatorname{Ext}\) 项为零，见式\eqref{g2:eq:uct-cohomology}。当 \(A=\mathbb Z\) 时，与 \(o_b\) 对偶的生成元 \(u_b\) 由
\[
	\langle u_b,o_b\rangle=1
\]
刻画。方向改变使 \(o_b\) 和 \(u_b\) 同时变号。这一配对将是把局部方向转成 Thom 类的标准化条件。

\ProseParagraphBreak

球面取向采用外法向在前的边界约定。长正合列的边界映射于是满足 \(\partial o_b=[S_b]\)。秩一时这个公式读成
\[
	\partial[-1,1]=[1]-[-1]\in\widetilde H_0(S^0;\mathbb Z).
\]
若球面的约化顶维上同调类 \(\omega_b\) 在 \([S_b]\) 上取值一，则对的连接映射满足 \(\delta\omega_b=u_b\)：在链上有 \(\langle\delta\omega_b,o_b\rangle=\langle\omega_b,\partial o_b\rangle=1\)。这个约定固定了后面球丛长正合列中纤维积分的正负号。

\begin{definition}\label{g2:def:bundle-orientation-system}
	向量丛 \(E\) 的\term[xiangliangcong-dingxiang-jubuxitong]{定向局部系统}[orientation local system]记为 \(\mathcal O_E\)，其在 \(b\) 处的群为
	\[
		\mathcal O_{E,b}=H_r(E_b,E_b\setminus\{0\};\mathbb Z).
	\]
	局部平凡化给出这些群之间的识别，沿路径依次使用局部识别便得到输运。若可为每个 \(b\) 选取生成元 \(o_b\in\mathcal O_{E,b}\)，使它们在局部平凡化中相容，则称 \(E\) \term[xiangliangcong-kedingxiang]{可定向}[orientable]；这样一组选择称为 \(E\) 的一个定向。
\end{definition}

这一定向描述每根向量空间纤维的方向，基底本身不需要是流形，也不需要已有定向。它与\chapref{g2:ch:05}中流形的定向局部系统使用同一种局部同调思想；此处取的是纤维在零点的局部同调。

\begin{proposition}\label{g2:prop:bundle-orientation-transport}
	若两组局部坐标的过渡函数为 \(g_{\alpha\beta}:U_\alpha\cap U_\beta\to\operatorname{GL}_r(\mathbb R)\)，则 \(\mathcal O_E\) 的相应过渡函数为
	\[
		\operatorname{sgn}\det g_{\alpha\beta}\in\{1,-1\}.
	\]
	这些函数定义秩一的整数局部系统，其路径输运只依赖端点固定的路径同伦类。在每个路径连通分支上，\(E\) 可定向当且仅当所有闭路在 \(\mathcal O_E\) 上的输运都是恒等。
\end{proposition}
\begin{proof}
	先考察一个可逆线性映射 \(A:\mathbb R^r\to\mathbb R^r\)。极分解把它写成正交矩阵与正定矩阵的乘积；正定矩阵可沿 \((1-t)P+tI\) 连到恒等。行列式为正的正交矩阵由二维平面旋转组成，可以把各旋转角连续缩到零。因此 \(\det A>0\) 时，\(A\) 与恒等映射通过可逆线性映射同伦，在相对顶维同调上作用为 \(+1\)。若 \(\det A<0\)，先复合一次坐标反射便化为前一种情形；反射在相对基本类上作用为 \(-1\)，与\cref{g2:prop:reflection-antipodal-degree}的球面度计算一致。故所求作用正是 \(\operatorname{sgn}\det A\)。

	行列式连续且不为零，所以其符号在交集上局部常值。过渡函数的乘法关系经行列式符号仍成立，因而确实得到局部系统。用这些局部常值识别沿路径输运，细分路径或改变局部平凡化均由过渡关系保证不改变合成；对路径同伦，把参数方形细分到每个小方形落入一个平凡化域，小方形内的输运相消，得到端点固定的同伦不变性。这也就是\secref{g2:sec:0501}中局部系统的路径描述。

	若已有相容生成元，任何闭路都把它送回自身，故输运恒等。反过来，在一个分支中固定 \(b_0\) 及一个生成元 \(o_{b_0}\)，沿从 \(b_0\) 到 \(b\) 的路径输运它。两条路径之差是一条闭路，所以假设保证所得 \(o_b\) 与路径无关；局部平凡化中的输运恒定，又保证这些生成元相容。CW 复形局部道路连通，因此上述逐分支的构造给出所需定向。
\end{proof}

\begin{example}\label{g2:exm:complex-bundle-real-orientation}
	每张复秩 \(m\) 向量丛的底层实秩 \(2m\) 丛都有自然定向。对一组复基 \(e_1,\ldots,e_m\)，取实有序基
	\[
		(e_1,\mathrm i e_1,e_2,\mathrm i e_2,\ldots,e_m,\mathrm i e_m).
	\]
	若复坐标变换矩阵为 \(A\)，其底层实矩阵的行列式为 \(\det_{\mathbb R}A=|\det_{\mathbb C}A|^2>0\)。这个公式可在复化后看出：实矩阵的复化分解为 \(A\) 与 \(\overline A\) 两块，行列式为 \(\det A\cdot\overline{\det A}\)。因此更换复标架不会改变上述实方向，局部方向自动粘合。
\end{example}

\subsection{系数环及扭定向系统}
\label{g2:subsec:2001:03}
取交换有幺环 \(R\)，并令
\[
	\mathcal O_{E,R}=\mathcal O_E\otimes_{\mathbb Z}R.
\]
这是秩一自由 \(R\) 模的局部系统，过渡函数仍为 \(\pm1\)，现将它们视作 \(R\) 中的单位。一个 \(R\) 定向是 \(\mathcal O_{E,R}\) 与常局部系统 \(\underline R\) 的同构，等价地，是一组相容的纤维 \(R\) 模基元。整数定向经系数变换给出 \(R\) 定向；若 \(2=0\) 于 \(R\)，则两个整数方向都给出同一个模 \(R\) 基元，故每张实向量丛都有规范的 \(R\) 定向。特别地，模二系数完全消除了方向反转造成的符号。

\begin{example}\label{g2:exm:mobius-orientation-system}
	设 \(L\to S^1\) 为 Möbius 线丛，其矩形模型沿两端按 \((0,t)\sim(1,-t)\) 粘合。沿基圆一周，纤维的相对基本类变成其相反数，所以 \(\mathcal O_L\) 的单值作用为 \(-1\)，它没有整数定向。若一个类 \(u\in H^1(D(L),S(L);\mathbb Z)\) 在每根纤维上都是生成元，则沿基圆输运后其限制必须仍为自身，同时又要变号。这要求一个整数生成元等于其相反数，矛盾。用 \(\mathbb F_2\) 系数时两个符号相同，方向障碍消失。\secref{g2:sec:1304}已经利用这一模二 Thom 类研究了平方运算，见\cref{g2:exm:mobius-thom-square}。
\end{example}

为了对所有实丛统一处理整数系数，可以让系数也随纤维方向一起变化。固定 \(b\)，式\eqref{g2:eq:fiber-top-cohomology-hom}给出规范识别
\begin{equation}\label{g2:eq:twisted-fiber-endomorphism}
	H^r(D_b,S_b;\mathcal O_{E,R,b})
	\cong\operatorname{Hom}_{\mathbb Z}(\mathcal O_{E,b},\mathcal O_{E,R,b})
	\cong\operatorname{End}_R(\mathcal O_{E,R,b}).
\end{equation}
第二个同构把群同态唯一延拓为 \(R\) 线性映射。右端有一个不依赖任何基元的元素：恒等映射。整数系数时它就是 \(\operatorname{id}_{\mathcal O_{E,b}}\)；一般系数时，在中间一项对应 \(m\mapsto m\otimes1\)。

\begin{definition}\label{g2:def:twisted-thom-class}
	类
	\[
		U_E\in H^r(D(E),S(E);p^*\mathcal O_{E,R})
	\]
	若在每个纤维上的限制都对应式\eqref{g2:eq:twisted-fiber-endomorphism}中的恒等映射，就称为 \(E\) 的\term[niu-thom-lei]{扭 Thom 类}[twisted Thom class]。需要区分系数环时记为 \(U_{E,R}\)。若选定 \(E\) 的 \(R\) 定向 \(\mathcal O_{E,R}\cong\underline R\)，由它得到的常系数类记为 \(u_E\)，称为这个定向的 Thom 类；其纤维限制在选定的相对基本类上取值一。
\end{definition}

这一标准化把两次变号配成了方向无关的条件。具体地，选一个整数方向 \(o_b\)，并令 \(u_b(o_b)=1\)，则局部扭类写成 \(u_b\otimes(o_b\otimes1)\)。更换方向时两个因子都变号，张量保持不变；在式\eqref{g2:eq:twisted-fiber-endomorphism}中，这正是同一个恒等映射。在 Möbius 例中，它也沿闭路保持不变。局部条件能否确定一个全局上同调类，需要进一步的粘合证明；下一节的\cref{g2:thm:thom-isomorphism}将证明这样的类存在、唯一，并说明乘以它为何给出次数平移。

\begin{proposition}\label{g2:prop:bundle-orientation-square}
	存在规范的局部系统同构
	\begin{equation}\label{g2:eq:bundle-orientation-square}
		\mathcal O_{E,R}\otimes_R\mathcal O_{E,R}\xrightarrow{\ \cong\ }\underline R.
	\end{equation}
	在任一整数方向 \(o_b\) 下，它将 \((o_b\otimes a)\otimes(o_b\otimes c)\) 送到 \(ac\)。
\end{proposition}
\begin{proof}
	每个纤维都是秩一自由 \(R\) 模，所以该规则给出一个同构。把 \(o_b\) 换成 \(-o_b\) 时，同一元素的两个系数都变成相反数，乘积不变，故映射与方向选择无关。沿路径输运时，两个定向因子具有同一符号，其乘积为 \((\pm1)^2=1\)，因而这些纤维同构与输运相容，给出局部系统同构。
\end{proof}

这与\cref{g2:prop:orientation-square}中流形定向系统的平方配对是同一种构造。它解释了后面公式中为什么会出现一个或两个定向因子：乘一个扭 Thom 类会添上 \(\mathcal O_{E,R}\)，两个相同的定向因子则可按式\eqref{g2:eq:bundle-orientation-square}规范消去。此后的杯积按“底空间类在前、Thom 类在后”排列，记作 \(p^*a\smile U_E\)。这一顺序同时确定次数交换可能带来的符号。
% 纲要条目排为编号小节标题，并进入目录与书签；说明段仍保留于纲要注释中。

% !TeX root = ../../Book2.tex
% 纲要标题：### 20.2 Thom 同构的证明
% 纲要位置：几何学计划纲要.md，第 1729 行。
% 正文状态：未撰写；本文件目前是框架。

% 纲要内容开始
% 
% * 局部乘积上的杯积同构
% * Mayer–Vietoris 粘合
% * 一般基空间的有限维或穷竭控制
% * 自然性与直和相容
% 
% 纲要内容结束

\section{Thom 同构的证明}
\label{g2:sec:2002}

待撰写

\subsection{局部乘积上的杯积同构}
\label{g2:subsec:2002:01}
待撰写

\subsection{Mayer–Vietoris 粘合}
\label{g2:subsec:2002:02}
待撰写

\subsection{一般基空间的有限维或穷竭控制}
\label{g2:subsec:2002:03}
待撰写

\subsection{自然性与直和相容}
\label{g2:subsec:2002:04}
待撰写
% 纲要条目排为编号小节标题，并进入目录与书签；说明段仍保留于纲要注释中。

% !TeX root = ../../Book2.tex
% 纲要标题：### 20.3 Euler 类与零点
% 纲要位置：几何学计划纲要.md，第 1736 行。
% 正文状态：未撰写；本文件目前是框架。

% 纲要内容开始
% 
% * 零截面拉回
% * 横截截面的零点类
% * 局部指标与 Euler 数
% * 零 Euler 类不足以消除所有高阶障碍
% 
% 纲要内容结束

\section{Euler 类与零点}
\label{g2:sec:2003}

待撰写

\subsection{零截面拉回}
\label{g2:subsec:2003:01}
待撰写

\subsection{横截截面的零点类}
\label{g2:subsec:2003:02}
待撰写

\subsection{局部指标与 Euler 数}
\label{g2:subsec:2003:03}
待撰写

\subsection{零 Euler 类不足以消除所有高阶障碍}
\label{g2:subsec:2003:04}
待撰写
% 纲要条目排为编号小节标题，并进入目录与书签；说明段仍保留于纲要注释中。

% !TeX root = ../../Book2.tex
% 纲要标题：### 20.4 Gysin 序列及积分沿纤维
% 纲要位置：几何学计划纲要.md，第 1793 行。
% 正文状态：未撰写；本文件目前是框架。

% 纲要内容开始
% 
% * 从盘球丛对的长正合列推导
% * 乘 Euler 类的映射
% * Hopf 丛及球丛计算
% * 定向纤维积分的上同调模型
% 
% ---
% 
% 纲要内容结束

\section{Gysin 序列及积分沿纤维}
\label{g2:sec:2004}

待撰写

\subsection{从盘球丛对的长正合列推导}
\label{g2:subsec:2004:01}
待撰写

\subsection{乘 Euler 类的映射}
\label{g2:subsec:2004:02}
待撰写

\subsection{Hopf 丛及球丛计算}
\label{g2:subsec:2004:03}
待撰写

\subsection{定向纤维积分的上同调模型}
\label{g2:subsec:2004:04}
待撰写
% 纲要条目排为编号小节标题，并进入目录与书签；说明段仍保留于纲要注释中。


\begin{chaptersummary}
	盘球丛对把每根纤维的方向记录为一个最高次相对生成元；定向局部系统同时保存这些生成元沿道路的符号变化。局部乘积给出次数平移，有限覆盖通过 Mayer--Vietoris 粘合，一般 CW 基空间则由相对胞腔复形控制。所得 Thom 类唯一，并与拉回、直和及系数变换相容；任意局部系数的上同调都可由与它作杯积移入盘球丛对。

	\ProseParagraphBreak

	零截面拉回 Thom 类得到 \(e_{\mathrm{Th}}(E)\)。它等于球丛的首障碍，也等于横截零点流形的对偶类；秩与底流形维数相等时，求值就是孤立零点的指标和。按前文已固定的传递符号，它与球丛的 \(e(S(E))\) 相差一个负号。Euler 类消失只消除首障碍，高维胞腔还可能阻止非零截面存在。

	\ProseParagraphBreak

	盘球丛对的长正合列经 Thom 同构改写为 Gysin 序列：忘去相对条件对应乘 Euler 类，球丛上的类经连接映射再作 Thom 逆同构，得到降低纤维维数的积分映射。Hopf 丛与二球面上的圆丛展示了它的计算作用。按底空间微分在前、纤维微分在后的顺序，形式积分满足无附加符号的投影公式，并与规范化后的拓扑积分一致；整数序列还保留实微分形式看不到的扭信息。
\end{chaptersummary}

\BookExercises

\begin{exercise}\label{g2:ex:noncompact-thom-not-compactification}
	取 \(\mathbb R\) 上的平凡实线丛 \(E=\mathbb R\times\mathbb R\)。直接从盘球丛模型证明 \(\operatorname{Th}(E)\simeq S^1\)，并计算其约化整数上同调。再与总空间 \(E\) 的一点紧化比较，说明非紧基底时两种构造确实可能不同。
\end{exercise}
\begin{solution}
	盘丛和球丛分别是 \(\mathbb R\times[-1,1]\) 与 \(\mathbb R\times\{-1,1\}\)。同伦
	\[
		F_s(x,t)=((1-s)x,t),\qquad 0\le s\le1,
	\]
	始终把球丛映入球丛，固定中间纤维 \(\{0\}\times[-1,1]\)，所以下降为 Thom 空间上的同伦。下降后的同伦连续，由\cref{g2:lem:quotient-interval}的积商性质保证。终映射的像为中间纤维的商 \([-1,1]/\{-1,1\}\cong S^1\)，因而得到一个强形变收缩。于是
	\[
		\widetilde H^k(\operatorname{Th}(E);\mathbb Z)\cong
		\begin{cases}
			\mathbb Z,&k=1,\\
			0,&k\ne1.
		\end{cases}
	\]
	总空间则同胚于 \(\mathbb R^2\)，其一点紧化为 \(S^2\)，约化整数上同调的非零项在二次。因此它与 \(\operatorname{Th}(E)\) 连同伦等价都不成立。盘球丛商只将纤维的单位边界共同压成一点；总空间的一点紧化还把沿非紧基底逃向无穷远的方向纳入无穷远点的邻域条件。
\end{solution}

\begin{exercise}\label{g2:ex:doubled-bundle-orientation}
	设 \(E\to B\) 为实秩 \(r\) 向量丛。对任意局部实标架 \(e_1,\ldots,e_r\)，用有序标架
	\[
		((e_1,0),\ldots,(e_r,0),(0,e_1),\ldots,(0,e_r))
	\]
	给 \(E\oplus E\) 定向。证明该定向不依赖局部标架，所以即使 \(E\) 不可定向，\(E\oplus E\) 仍有这一规范定向。再用 \((v,w)\mapsto v+\mathrm i w\) 识别它与复化丛 \(E\otimes_{\mathbb R}\mathbb C\) 的底层实丛，比较题给定向与复结构的自然定向。
\end{exercise}
\begin{solution}
	若原标架的变换矩阵是 \(A\in\operatorname{GL}_r(\mathbb R)\)，则直和中这组有序标架的变换矩阵为 \(\operatorname{diag}(A,A)\)，行列式是 \((\det A)^2>0\)。所以所有局部方向相容，且更换任何局部标架都不改变定向。这一论证不要求 \(\det A\) 本身为正。

	在题给识别下，直和定向对应实有序基
	\[
		(e_1,\ldots,e_r,\mathrm i e_1,\ldots,\mathrm i e_r),
	\]
	而复结构的自然定向对应交错排列 \((e_1,\mathrm i e_1,\ldots,e_r,\mathrm i e_r)\)。把前一个排列变为后一个时，\(\mathrm i e_j\) 需要跨过 \(r-j\) 个实基向量，总交换次数为
	\[
		(r-1)+(r-2)+\cdots+1=\frac{r(r-1)}2.
	\]
	故两个定向的比较符号为 \((-1)^{r(r-1)/2}\)：\(r\equiv0,1\pmod4\) 时相同，\(r\equiv2,3\pmod4\) 时相反。两种定向都与标架无关，但所采用的有序基排列不同。
\end{solution}

\begin{exercise}\label{g2:ex:orientation-over-z-four}
	令 \(L\to S^1\) 为 Möbius 线丛。分别对 \(R=\mathbb Z\)、\(\mathbb Z/4\)、\(\mathbb F_2\)，计算 \(H^0(S^1;\mathcal O_{L,R})\)，并判断 \(L\) 是否 \(R\) 可定向。说明存在非零平行截面为何还不足以给出 \(R\) 定向。
\end{exercise}
\begin{solution}
	这三个定向系统沿正向圆周的单值作用都是乘以 \(-1\)。\cref{g2:prop:circle-local-coefficients}给出圆周局部系数零次上同调为单值作用的不变量，因此
	\[
		H^0(S^1;\mathcal O_{L,R})=\ker(-1-1:R\longrightarrow R)=\{a\in R:2a=0\}.
	\]
	在三种系数中依次得到 \(0\)、\(\{0,2\}\subset\mathbb Z/4\)、\(\mathbb F_2\)。一个 \(R\) 定向要求平行截面的值为秩一自由 \(R\) 模的基元，选定一个局部基后，这等价于该值是 \(R\) 中的单位。

	整数系数没有非零平行截面，所以不可定向。\(\mathbb Z/4\) 系数虽有非零平行截面 \(2\)，但它不是单位，所生成的子模只有 \(\{0,2\}\)，不能成为纤维模的一组基；故仍不可定向。\(\mathbb F_2\) 系数的平行截面 \(1\) 是单位，因而给出规范定向。第二个例子说明，在带有零因子的系数环上，“非零”与“生成整个秩一自由模”是两个不同的条件。
\end{solution}

\begin{exercise}\label{g2:ex:mobius-thom-cohomology}
	设 \(L\to S^1\) 为 Möbius 线丛。分别计算 \(H^k(D(L),S(L);\mathbb Z)\)、\(H^k(D(L),S(L);p^*\mathcal O_L)\) 与 \(H^k(D(L),S(L);\mathbb F_2)\)。再用盘球丛对的长正合列独立验证第一组结果，并解释其二次整数类为什么不能充当 Thom 类。
\end{exercise}
\begin{solution}
	由式\eqref{g2:eq:twisted-thom-isomorphism}及 \(\mathcal O_L\otimes\mathcal O_L\cong\mathbb Z\)，普通整数相对上同调对应 \(H^{k-1}(S^1;\mathcal O_L)\)，扭系数相对上同调对应 \(H^{k-1}(S^1;\mathbb Z)\)。定向系统的圆周胞腔上链为
	\[
		\mathbb Z\xrightarrow{\,-2\,}\mathbb Z,
	\]
	所以它的零次上同调为零、一次为 \(\mathbb Z/2\)。于是
	\[
	\begin{array}{c|ccc}
		&k=1&k=2&\text{其余 }k\\\hline
		H^k(D(L),S(L);\mathbb Z)&0&\mathbb Z/2&0\\
		H^k(D(L),S(L);p^*\mathcal O_L)&\mathbb Z&\mathbb Z&0\\
		H^k(D(L),S(L);\mathbb F_2)&\mathbb F_2&\mathbb F_2&0
	\end{array}
	\]
	最后一行使用模二方向系统平凡及圆周的模二上同调。

	独立验证时，\(D(L)\) 收缩到基圆，\(S(L)\) 也是一个圆，而其投影绕基圆两次。故限制 \(H^1(D(L);\mathbb Z)\to H^1(S(L);\mathbb Z)\) 为乘二；两者的零次限制为恒等。长正合列于是给出
	\[
		0\longrightarrow H^1(D(L),S(L);\mathbb Z)\longrightarrow\mathbb Z
		\xrightarrow{\,2\,}\mathbb Z\longrightarrow H^2(D(L),S(L);\mathbb Z)\longrightarrow0,
	\]
	以及其他次数消失，与上表相符。秩一丛的 Thom 类须是一次类，并在每根区间对上取生成元；这里普通整数一次相对群已经为零。二次的扭元素限制到纤维后为零，因为纤维对的二次上同调为零，所以它不能替代 Thom 类。
\end{solution}

\begin{exercise}\label{g2:ex:infinite-thom-cochains}
	令 \(B=\bigvee_{j\ge1}S^1_j\) 具有 CW 弱拓扑，\(E=B\times\mathbb R^2\) 取标准定向。计算 \(\widetilde H^k(\operatorname{Th}(E);\mathbb Z)\)。找出一个三次类，使它在每个圆周上限制后都非零，并说明用有限子楔和的上同调按补零映射取直极限会漏掉什么。
\end{exercise}
\begin{solution}
	\(B\) 有一个零胞腔及可数多个一胞腔，胞腔边界全为零。因此
	\[
		H^0(B;\mathbb Z)=\mathbb Z,\qquad H^1(B;\mathbb Z)=\prod_{j\ge1}\mathbb Z,
	\]
	更高次为零。Thom 同构与常系数商比较给出
	\[
		\widetilde H^k(\operatorname{Th}(E);\mathbb Z)=
		\begin{cases}
			\mathbb Z,&k=2,\\
			\prod_{j\ge1}\mathbb Z,&k=3,\\
			0,&\text{其他次数}.
		\end{cases}
	\]
	取一次胞腔上链 \(a=(1,1,\ldots)\)。它是上循环且不为余边界，\(\Phi_E(a)\) 在限制到任意 \(S^1_j\) 上的丛时都是三次群的生成元，这由 Thom 自然性可知。

	前 \(N\) 个圆周的上同调为 \(\mathbb Z^N\)。若按题给补零映射取直极限，得到的是 \(\bigoplus_{j\ge1}\mathbb Z\)，只含有限支撑序列；它不含 \((1,1,\ldots)\)。实际由子空间包含诱导的上同调映射是反向的限制，而非补零映射。胞腔计算使用全部胞腔系数的直积，因而保留了这个类，也说明固定次数的有限骨架控制不等于有限胞腔支撑。
\end{solution}

\begin{exercise}\label{g2:ex:thom-space-product-euler}
	设 \(E\to B\) 为已定向实秩 \(r\ge1\) 丛，取系数环 \(R\)，记 \(e=e_{\mathrm{Th}}(E)\)。证明 Thom 同构一般只是 \(H^*(B;R)\) 模的次数平移同构；其乘法满足
	\[
		\Phi_E(a)\smile\Phi_E(b)=(-1)^{r\deg b}\Phi_E(a\smile b\smile e).
	\]
	特别计算 \(u_E^2\)，并说明为什么平凡正秩丛的 Thom 空间约化上同调中任意两个类的乘积都为零。
\end{exercise}
\begin{solution}
	记 \(j^*:H^*(D(E),S(E);R)\to H^*(D(E);R)\) 为忘记相对条件的映射。零截面与投影互为同伦逆，而 Euler 类定义给出 \(s_0^*j^*u_E=e\)，故
	\[
		j^*u_E=p^*e.
	\]
	相对杯积与绝对—相对杯积相容：两个相对上链都在 \(S(E)\) 上为零，将其中一个视为绝对上链，不改变它们的杯积。因此
	\[
		u_E\smile u_E=j^*u_E\smile u_E=p^*e\smile u_E=\Phi_E(e).
	\]
	对一般齐次类，把次数 \(r\) 的 \(u_E\) 越过 \(p^*b\)，得到
	\[
		(p^*a\smile u_E)\smile(p^*b\smile u_E)
		=(-1)^{r\deg b}p^*(a\smile b\smile e)\smile u_E,
	\]
	即题给公式。非齐次类按齐次分量展开即可。

	平凡正秩丛有常值非零截面，由\cref{g2:prop:euler-basic-properties}得 \(e=0\)，所以上式恒为零。经商比较，这就是 Thom 空间约化上同调的乘法；它不意味着底空间本身的上同调乘法为零。Thom 同构保留底空间类作用在相对类上的模结构，约化环的内部乘法则额外记住 Euler 类。
\end{solution}

\begin{exercise}\label{g2:ex:odd-rank-euler-torsion}
	设 \(E\to B\) 是奇数秩的定向实向量丛。证明 \(2e_{\mathrm{Th}}(E)=0\)。据此说明：若 \(H^{\operatorname{rank}E}(B;\mathbb Z)\) 没有非零的二阶扭元，则 Euler 类为零；但这一点仍不保证存在非零截面。再用这个结论与\cref{g2:cor:tangent-euler-characteristic}证明闭定向奇维流形的 Euler 示性数为零。
\end{exercise}
\begin{solution}
	记 \(E^{\mathrm{op}}\) 为反转纤维定向后的同一个丛。逐纤维映射 \(-I:E\to E^{\mathrm{op}}\) 保持定向：它相对于原定向的行列式符号为 \((-1)^{\operatorname{rank}E}=-1\)，而目标又反转了一次定向。Euler 类对保持定向的丛同构不变，因此
	\[
		e_{\mathrm{Th}}(E)=e_{\mathrm{Th}}(E^{\mathrm{op}})=-e_{\mathrm{Th}}(E).
	\]
	于是 \(2e_{\mathrm{Th}}(E)=0\)。若相应上同调群没有非零二阶扭元，就有 \(e_{\mathrm{Th}}(E)=0\)。\cref{g2:ex:zero-euler-no-section}中的 \(S^4\) 上定向三维丛满足这个条件，却没有非零截面，故消失条件不充分。

	若 \(M\) 是闭定向奇维流形，则整数求值满足 \(2\langle e_{\mathrm{Th}}(\mathrm TM),[M]\rangle=0\)，而整数群没有非零二阶扭元，所以这个求值为零。由\cref{g2:cor:tangent-euler-characteristic}，它恰是 \(\chi(M)\)。
\end{solution}

\begin{exercise}\label{g2:ex:product-tangent-euler}
	设 \(M,N\) 是闭定向光滑流形，维数分别为正整数 \(m,n\)。乘积按 \(M\) 在前、\(N\) 在后定向。用切丛分解和 Euler 类证明
	\[
		e_{\mathrm{Th}}(\mathrm T(M\times N))
		=\pi_M^*e_{\mathrm{Th}}(\mathrm TM)\smile\pi_N^*e_{\mathrm{Th}}(\mathrm TN),
		\qquad\chi(M\times N)=\chi(M)\chi(N).
	\]
	进一步，若 \(X,Y\) 分别是 \(M,N\) 上仅有非退化零点的向量场，计算乘积向量场 \((X,Y)\) 在 \((x,y)\) 处的指标。
\end{exercise}
\begin{solution}
	乘积坐标给出保持所选定向的丛同构
	\[
		\mathrm T(M\times N)\cong\pi_M^*\mathrm TM\oplus\pi_N^*\mathrm TN.
	\]
	由拉回自然性和直和公式，得到所求的第一式。乘积基本类为 \([M]\times[N]\)，而按同一因子次序的外积满足
	\[
		\bigl\langle\pi_M^*a\smile\pi_N^*b,[M]\times[N]\bigr\rangle
		=\langle a,[M]\rangle\langle b,[N]\rangle
	\]
	对 \(a\in H^m(M;\mathbb Z)\)、\(b\in H^n(N;\mathbb Z)\) 成立。这里没有交换两个因子，故没有附加符号。将两个 Euler 类代入，并在三个流形上使用\cref{g2:cor:tangent-euler-characteristic}，便得第二式。

	乘积向量场的零点恰为两个零点集的笛卡儿积。在正向乘积坐标下，其线性化为分块对角矩阵
	\[
		\begin{pmatrix}\mathrm DX_x&0\\0&\mathrm DY_y\end{pmatrix}.
	\]
	因此每个零点非退化，且
	\[
		\operatorname{ind}_{(x,y)}(X,Y)
		=\operatorname{sgn}(\det\mathrm DX_x\det\mathrm DY_y)
		=\operatorname{ind}_x(X)\operatorname{ind}_y(Y).
	\]
	对所有零点求和，得到两个指标和的乘积，也从零点计数直接解释了 Euler 数的乘法公式。
\end{solution}

\begin{exercise}\label{g2:ex:pullback-tangent-sphere-euler}
	设 \(f:S^2\to S^2\) 是次数为 \(d\in\mathbb Z\) 的光滑映射，给 \(E=f^*\mathrm TS^2\) 取拉回定向。计算 \(e_{\mathrm{Th}}(E)\)，判定 \(E\) 何时存在非零截面。若 \(s\) 是与零截面横截的截面，证明它至少有 \(2|d|\) 个零点，并且零点总数为偶数。
\end{exercise}
\begin{solution}
	令 \(\omega\in H^2(S^2;\mathbb Z)\) 在正基本类上取值一。由自然性、映射度的定义和\cref{g2:ex:euler-tangent-spheres-surfaces}，
	\[
		e_{\mathrm{Th}}(E)=f^*(2\omega)=2d\omega.
	\]
	若有非零截面，Euler 类必须为零，故 \(d=0\)。反过来，若 \(d=0\)，Euler 类为零，而底空间维数等于丛的秩；\cref{g2:thm:euler-primary-comparison}与\cref{g2:thm:sphere-primary-euler}说明首障碍消失后即可得到全空间上的单位截面。因此非零截面存在当且仅当 \(d=0\)。这里先得到连续截面；若要求光滑，可将其充分接近地光滑逼近，紧性保证逼近后仍非零。

	设横截零点中指标为 \(+1,-1\) 的个数分别为 \(N_+,N_-\)。\cref{g2:thm:euler-index-evaluation}给出 \(N_+-N_-=2d\)，故
	\[
		N_++N_-\ge|N_+-N_-|=2|d|,
		\qquad N_++N_-=2d+2N_-.
	\]
	第一式给出下界，第二式说明总数为偶数。\(d=0\) 时下界允许没有零点，也允许某个截面有符号相消的零点对。
\end{solution}

\begin{exercise}\label{g2:ex:disk-relative-euler-index}
	给闭圆盘 \(D^2\subset\mathbb C\) 取标准定向，在平凡定向二维丛上指定一个非零边界截面 \(v:S^1\to\mathbb C^*\)，令 \(k=\deg(v/|v|)\)。证明 \(v\) 可以延拓为处处非零的连续截面当且仅当 \(k=0\)。若一个光滑延拓 \(F:D^2\to\mathbb C\) 的零点全部孤立，证明其指标和为 \(k\)。当 \(v(z)=z^k\) 时，对任意整数 \(k\) 写出达到这个指标和的具体光滑延拓，并解释为何这不与平凡丛的绝对 Euler 类为零矛盾。
\end{exercise}
\begin{solution}
	若有非零延拓，将其除以长度便得到 \(D^2\to S^1\) 的延拓，故边界映射零伦，次数为零。反之，\(k=0\) 时 \(v/|v|:S^1\to S^1\) 零伦。用这个零伦沿半径填入圆盘，可得连续映射 \(g:D^2\to S^1\)，其边界值为 \(v/|v|\)。边界上的正函数 \(|v|\) 也有正延拓，例如取
	\[
		h(re^{\mathrm i\theta})=(1-r)+r|v(e^{\mathrm i\theta})|,
		\qquad h(0)=1.
	\]
	于是 \(hg\) 是所求非零延拓。

	设光滑延拓的零点为 \(z_1,\ldots,z_\ell\)，在每个零点周围取小闭圆盘 \(D_i\)，互不相交且不碰外边界。在挖去这些小圆盘内部的区域上，\(F/|F|\) 是到 \(S^1\) 的连续映射。该区域正向边界的同调类为外边界减去各小圆盘的正向边界；边界总类在区域的一维同调中为零。映入 \(H_1(S^1;\mathbb Z)\) 后便得
	\[
		k-\sum_{i=1}^\ell\deg\bigl((F/|F|)|_{\partial D_i}\bigr)=0.
	\]
	每个内边界次数按小圆盘本身的正向计算，正是对应零点的指标，故总和为 \(k\)。

	对 \(v(z)=z^k\)，若 \(k>0\)，取 \(F(z)=z^k\)；若 \(k<0\)，取 \(F(z)=\overline z^{\,-k}\)；若 \(k=0\)，取 \(F(z)=1\)。前两种延拓只有原点一个零点，其小圆周映射次数分别为 \(k\)，后一种没有零点。绝对 Euler 类属于 \(H^2(D^2;\mathbb Z)=0\)，只描述是否能自由选择一个非零截面；本题要求保持指定边界值，障碍属于 \(H^2(D^2,S^1;\mathbb Z)\cong\mathbb Z\)，恰由边界绕数 \(k\) 给出。
\end{solution}

\begin{exercise}\label{g2:ex:gysin-product-normalization}
	设 \(B\) 满足本章 Thom 同构的条件，\(r\ge2\)，并取平凡定向丛 \(E=B\times\mathbb R^r\)。令 \(u\) 是球面 \(S^{r-1}\) 的正向上同调生成元，仍用 \(u\) 表示它在 \(B\times S^{r-1}\) 上的拉回。证明
	\[
		H^k(B\times S^{r-1};R)=\pi^*H^k(B;R)\oplus\bigl(\pi^*H^{k-r+1}(B;R)\smile u\bigr),
	\]
	并计算两部分的沿纤维积分。写出上同调环的乘法约定，解释 Gysin 序列在这个例子中如何分裂。
\end{exercise}
\begin{solution}
	球面只有零次及 \(r-1\) 次上同调，且都是自由秩一 \(R\) 模。第4章的乘积比较因此给出所述直和；也可在球面的一份零胞腔、一份 \(r-1\) 胞腔的链模型上直接得到同一分解。正向生成元满足 \(u^2=0\)，因为它来自球面，而 \(H^{2r-2}(S^{r-1};R)=0\)。若 \(a\in H^\ell(B;R)\)，乘积中还有
	\[
		u\smile\pi^*a=(-1)^{(r-1)\ell}\pi^*a\smile u.
	\]
	这两条关系连同基底的乘法，给出分次张量积环 \(H^*(B;R)\otimes H^*(S^{r-1};R)\)。

	由正合性，\(\pi_!\pi^*a=0\)；由纤维归一化及投影公式，\(\pi_!(\pi^*b\smile u)=b\)。故
	\[
		\pi_!(\pi^*a+\pi^*b\smile u)=b.
	\]
	平凡丛有非零截面，所以 Euler 类为零，Gysin 序列分解成短正合列
	\[
		0\longrightarrow H^k(B;R)\xrightarrow{\pi^*}H^k(B\times S^{r-1};R)\xrightarrow{\pi_!}H^{k-r+1}(B;R)\longrightarrow0.
	\]
	映射 \(b\mapsto\pi^*b\smile u\) 是右逆，给出所述分裂。这个例子也检验了符号：基底类写在 \(u\) 前面时，积分结果就是 \(b\)。
\end{solution}

\begin{exercise}\label{g2:ex:gysin-sphere-product-base}
	在 \(B=S^2\times S^2\) 上，取两因子的正向生成元 \(x,y\)，于是
	\[
		H^*(B;\mathbb Z)=\mathbb Z[x,y]/(x^2,y^2),\qquad\deg x=\deg y=2.
	\]
	设定向实二平面丛 \(E\to B\) 满足 \(e_{\mathrm{Th}}(E)=ax+by\)，其中 \((a,b)\ne(0,0)\)，令 \(d=\gcd(|a|,|b|)>0\)。计算 \(H^*(S(E);\mathbb Z)\) 的各次群，并用 \(X=\pi^*x\)、\(Y=\pi^*y\) 以及适当的三次、五次生成元描述乘法。
\end{exercise}
\begin{solution}
	所有非零的 Euler 乘法由下面两个整数矩阵给出：
	\[
		H^0(B)=\mathbb Z\longrightarrow H^2(B)=\mathbb Z^2,\quad 1\longmapsto(a,b),
	\]
	\[
		H^2(B)=\mathbb Z^2\longrightarrow H^4(B)=\mathbb Z,\quad (s,t)\longmapsto bs+at.
	\]
	第一个映射单射，其余核为 \(\mathbb Z\oplus\mathbb Z/d\)，可由 Bézout 等式把列向量 \((a,b)\) 通过整数可逆行变换化成 \((d,0)\) 得到。第二个映射的像为 \(d\mathbb Z\)，核由 \((a/d,-b/d)\) 生成。基底的奇数次上同调为零，故 Gysin 序列在各次直接给出
	\[
		H^j(S(E);\mathbb Z)=
		\begin{cases}
			\mathbb Z,&j=0,3,5,\\
			\mathbb Z\oplus\mathbb Z/d,&j=2,\\
			\mathbb Z/d,&j=4,\\
			0,&\text{其余次数}.
		\end{cases}
	\]
	这里 \(\mathbb Z/1=0\)。二次群由 \(X,Y\) 生成，唯一加法关系为 \(aX+bY=0\)。四次群由 \(XY=\pi^*(xy)\) 生成，阶恰为 \(d\)，并有 \(X^2=Y^2=0\)。

	积分把三次群同构到上述第二个矩阵的核，把五次群同构到 \(H^4(B)\)。因此可选择唯一的生成元 \(z,w\)，使
	\[
		\pi_!z=\frac adx-\frac bdy,\qquad \pi_!w=xy.
	\]
	投影公式给出
	\[
		\pi_!(Xz)=x\smile\left(\frac adx-\frac bdy\right)=-\frac bdxy,\qquad
		\pi_!(Yz)=\frac adxy.
	\]
	由于五次积分是同构，得到 \(Xz=-(b/d)w\)、\(Yz=(a/d)w\)。其余正次数乘积由这些关系、分次交换律及五次以上群为零决定。这也展示了为什么只列出加法群还不够：同样的次数中，乘法仍保留 Euler 类两个分量的具体数值。
\end{solution}

\begin{exercise}\label{g2:ex:gysin-fiber-degree}
	设 \(E,E'\to B\) 为同秩 \(r\ge2\) 的定向实向量丛，\(B\) 连通且满足 Thom 同构的条件。设 \(F:S(E)\to S(E')\) 覆盖 \(\operatorname{id}_B\)，并在每条正向球面纤维上具有同一个整数映射度 \(m\)。证明
	\[
		e_{\mathrm{Th}}(E')=m\,e_{\mathrm{Th}}(E),\qquad
		\pi_!F^*b=m\,\pi'_!b.
	\]
	特别说明：纤维反向的同构会怎样改变 Euler 类和沿纤维积分。
\end{exercise}
\begin{solution}
	对 \(0<t\le1\)、\(v\in S(E)\)，定义 \(\overline F(tv)=tF(v)\)，并把各零向量送到相应的零向量。这是从 \((D(E),S(E))\) 到 \((D(E'),S(E'))\) 的连续映射，且覆盖基底恒等。零截面附近的连续性来自 \(\|\overline F(tv)\|=t\)；球面方向无论怎样变化，像的长度都趋于零。

	每条盘纤维上的相对映射度等于边界球面上的映射度 \(m\)，这是盘球对长正合列的自然性。故 \(\overline F^*u_{E'}\) 限制到每条纤维都是 \(m\) 倍正向生成元。由 Thom 同构在零次的唯一性，
	\[
		\overline F^*u_{E'}=m\,u_E.
	\]
	沿零截面拉回，得到 \(e_{\mathrm{Th}}(E')=m\,e_{\mathrm{Th}}(E)\)。对任意基底类 \(c\)，又有
	\[
		\overline F^*\Phi_{E'}(c)=p^*c\smile\overline F^*u_{E'}=m\Phi_E(c).
	\]
	空间对连接映射的自然性给出 \(\delta_EF^*b=\overline F^*\delta_{E'}b\)。两丛秩相同，积分定义中的次数符号也相同，因此代入上述等式便有 \(\pi_!F^*b=m\pi'_!b\)。

	纤维反向同构对应 \(m=-1\)，所以两边按所选方向计算的 Euler 类互为负数，积分也相差负号。若只反转同一丛的定向而保持底层映射为恒等，这正是 Thom 类、Euler 类和纤维积分同时反号的规则。
\end{solution}

\begin{exercise}\label{g2:ex:angular-form-thom}
	设 \(E\to B\) 为光滑定向实二平面丛，\(\pi:S(E)\to B\) 为其单位圆丛。给定一形式 \(\theta\in\Omega^1(S(E))\) 及闭二形式 \(\omega\in\Omega^2(B)\)，满足
	\[
		\int_{S(E)_x}\theta=1,\qquad \mathrm d\theta=\pi^*\omega.
	\]
	在 \(D(E)\setminus s_0(B)\) 上记径向投影为 \(\rho(v)=v/\|v\|\)。取光滑径向函数 \(\chi\)，使其在零截面附近为零，在球丛附近为一。证明
	\[
		U=\mathrm d(\chi\rho^*\theta)-p^*\omega
	\]
	给出实 Thom 类的相对形式代表，并求 \(e_{\mathrm{Th}}(E)\) 的实系数像。最后证明：若这个实 Euler 类非零，则球丛上不存在积分为一的闭一形式。
\end{exercise}
\begin{solution}
	因为 \(\chi\) 在零截面邻域恒为零，\(\chi\rho^*\theta\) 可通过补零光滑延拓到整个盘丛。闭性 \(\mathrm dU=0\) 来自 \(\mathrm d^2=0\) 及 \(\mathrm d\omega=0\)。在球丛上，\(\chi=1\)，所以
	\[
		i^*U=\mathrm d\theta-\pi^*\omega=0.
	\]
	因此 \(U\) 是相对闭形式。限制到单条盘纤维时，\(p^*\omega\) 为零；由外法向优先的 Stokes 公式，
	\[
		\int_{D(E)_x}U=\int_{S(E)_x}\theta=1.
	\]
	每条纤维上都是正向相对生成元，故 Thom 类的唯一性给出 \([U]=u_E\) 的实系数像。

	在零截面附近，\(U=-p^*\omega\)，于是
	\[
		e_{\mathrm{Th}}(E)_{\mathbb R}=[s_0^*U]=-[\omega].
	\]
	这个负号由相对代表在边界必须为零决定，与本节 Gysin 的定向约定一致。若另有闭一形式 \(\eta\) 满足每条纤维积分为一，则实系数 Gysin 序列给出 \(\pi_![\eta]=1\)，从而 \(e_{\mathrm{Th}}(E)_{\mathbb R}=0\)。这与假设矛盾，故所述闭形式不存在。
\end{solution}

% !TeX root = ../../Book2.tex
% 纲要标题：## 第21章　Stiefel–Whitney、陈与 Pontryagin 类
% 纲要位置：几何学计划纲要.md，第 1802 行。
% 正文状态：初稿；尚未正式审核。

% 纲要内容开始
% 
% > 本章目标与证明范围：分别完成整数类与模2类的构造，再证明分裂原理和计算法。拓扑扭信息在此保留，不能用下一章的曲率形式取代。
% 
% 纲要内容结束

\chapter{Stiefel–Whitney、陈与 Pontryagin 类}
\label{g2:ch:21}
\BookChapterNumber{g2:ch:21}

\begin{chapterabstract}
	一张向量丛在不同维数上的扭转，可以由底空间上不同次数的上同调类记录。本章从实、复线丛出发，用射影丛的上同调关系构造 Stiefel--Whitney 类与陈类，再由复化定义 Pontryagin 类。Leray--Hirsch 定理和旗丛提供分裂原理，使线丛计算能够回到原来的向量丛。随后证明 Whitney 公式、顶类与 Euler 类的比较、Thom 平方解释和 Wu 公式，计算射影空间切丛与结构群提升障碍。最后在复 Grassmann 流形上把示性类与 Schubert 条件联系起来，完成四维复空间中二维子空间的上同调环及四条直线的相交计数。
\end{chapterabstract}

\begin{learninggoals}
	\begin{itemize}
		\item 将实线丛的方向输运、复线丛的 Euler 类和过渡函数联系起来，固定第一示性类的符号。
		\item 证明 Leray--Hirsch 定理、射影丛公式和分裂原理，说明形式根计算为何能决定原底空间上的等式。
		\item 构造 Stiefel--Whitney 类、陈类与 Pontryagin 类，掌握自然性、直和及复化公式，并辨认整数二阶扭信息。
		\item 完成 Thom 平方与 Wu 类的比较，用切丛计算检验向量场、几乎复结构和 spin 提升的条件。
		\item 在 Grassmann 流形上解释 Schubert 类及特殊类，计算 \(\operatorname{Gr}_2(\mathbb C^4)\) 的乘法并求解四条直线的相交问题。
	\end{itemize}
\end{learninggoals}

上一章的 Euler 类记录一张丛的截面能否避开零点。它的次数等于丛的实秩，因此只提供一个特定次数的信息：线丛的方向会沿闭路改变，较高秩丛又可能在二维或四维子空间上发生扭转，这些现象需要不同次数的类。示性类把一张丛组织成一组可拉回的上同调数据；同一个定义可以随映射限制到曲线、曲面或更高维的检验空间。

\ProseParagraphBreak

先考虑复秩二丛 \(E=L_1\oplus L_2\)。如果两条线的第一类分别为 \(x_1,x_2\)，它们的和与乘积就是最自然的对称数据：交换两条线不改变它们。但是一般丛并没有给定的线丛分解。我们会将 \(E\) 拉回到旗丛，在那里得到实际分解，并证明上同调拉回单射。于是在线丛上算出的等式能够返回原基空间。读者需要区分这两件事：分解发生在新的基空间，而等式由单射保证仍能检测原来的类。

\ProseParagraphBreak

为了避免先假定乘法公式再用它定义示性类，本章从射影丛的上同调关系读出各阶系数，再证明这些系数具有所需性质。实线丛使用模二一次类，复线丛使用整数二次类，两种构造具有相同的计算形状，系数却保留不同的信息。Pontryagin 类由实丛的复化得到；它的整数直和公式中还可能出现被二消去的修正项。我们将写出这个修正，并在实射影空间上给出非零例子。

\ProseParagraphBreak

本章使用第19章的丛分类、第20章的 Thom 与 Euler 构造，以及第13章的平方运算；所需条件会在具体证明中回引。分裂、平方与三类示性类的组织可参阅\cite[第23章第6--7节]{May1999}。最后一节的 Grassmann 计算把此前的上同调环、对偶和横截相交放在同一个空间里使用，也为后续代数几何中的交叉计算准备一个完整例子。

% !TeX root = ../../Book2.tex
% 纲要标题：### 21.1 线丛与通用生成元
% 纲要位置：几何学计划纲要.md，第 1806 行。
% 正文状态：未撰写；本文件目前是框架。

% 纲要内容开始
% 
% * 实线丛和第一 Stiefel–Whitney 类
% * 复线丛与第一 Chern 类
% * 过渡函数、分类空间和指数序列的比较
% 
% 纲要内容结束

\section{线丛与通用生成元}
\label{g2:sec:2101}

待撰写

\subsection{实线丛和第一 Stiefel–Whitney 类}
\label{g2:subsec:2101:01}
待撰写

\subsection{复线丛与第一 Chern 类}
\label{g2:subsec:2101:02}
待撰写

\subsection{过渡函数、分类空间和指数序列的比较}
\label{g2:subsec:2101:03}
待撰写
% 纲要条目排为编号小节标题，并进入目录与书签；说明段仍保留于纲要注释中。

% !TeX root = ../../Book2.tex
% 纲要标题：### 21.2 射影丛与分裂原理
% 纲要位置：几何学计划纲要.md，第 1812 行。
% 正文状态：未撰写；本文件目前是框架。

% 纲要内容开始
% 
% * Leray–Hirsch 的证明
% * 射影丛上同调
% * 旗丛拉回的单射性
% * 以形式根进行对称多项式计算
% 
% 纲要内容结束

\section{射影丛与分裂原理}
\label{g2:sec:2102}

待撰写

\subsection{Leray–Hirsch 的证明}
\label{g2:subsec:2102:01}
待撰写

\subsection{射影丛上同调}
\label{g2:subsec:2102:02}
待撰写

\subsection{旗丛拉回的单射性}
\label{g2:subsec:2102:03}
待撰写

\subsection{以形式根进行对称多项式计算}
\label{g2:subsec:2102:04}
待撰写
% 纲要条目排为编号小节标题，并进入目录与书签；说明段仍保留于纲要注释中。

% !TeX root = ../../Book2.tex
% 纲要标题：### 21.3 一般示性类的构造
% 纲要位置：几何学计划纲要.md，第 1819 行。
% 正文状态：初稿；尚未正式审核。

% 纲要内容开始
% 
% * 陈类的存在唯一性
% * Stiefel–Whitney 类及实射影丛关系
% * Whitney 和公式与 Wu 公式
% * 对偶、张量和复化
% * Pontryagin 类及其性质
% 
% 纲要内容结束

\section{一般示性类的构造}
\label{g2:sec:2103}

射影丛公式把一张秩 \(r\) 丛的粘合信息写入一个首一多项式关系，旗丛又把这个多项式变成线性因子的乘积。现在用关系的系数定义一般示性类，再利用旗丛拉回的单射性，把线丛上的计算送回原基空间。这一顺序同时解决存在性、选择无关性与计算问题。

以下主要以 CW 复形为基，有限秩丛均带纤维内积。证明中的射影丛及旗丛使用上一节的弱 CW 逼近比较。Thom 类也可沿同一比较运输：盘丛、球丛均为局部平凡丛，拉回到 CW 弱逼近后，总空间映射由纤维化同伦长正合列为弱等价；绝对上同调比较及对的长正合列给出相对上同调同构。因此第20章的 Thom、Euler 及其自然性可用于这些实际拉回丛。

\subsection{陈类的存在唯一性}
\label{g2:subsec:2103:01}
设 \(E\to B\) 为复秩 \(r\ge1\) 丛。在 \(p:\mathbb P(E)\to B\) 上记典范线丛为 \(\lambda\)，置 \(h=c_1(\lambda^*)\)。\cref{g2:thm:projective-bundle}说明 \(1,h,\ldots,h^{r-1}\) 是 \(H^*(B;\mathbb Z)\) 模基。因此 \(h^r\) 可以唯一写成这些幂的线性组合；其各项的总次数必须同为 \(2r\)。

\begin{definition}\label{g2:def:chern-classes-projective}
	复秩 \(r\) 丛 \(E\to B\) 的\term[chen-lei]{陈类}[Chern classes]为唯一满足
	\begin{equation}\label{g2:eq:chern-projective-relation}
		h^r+p^*c_1(E)h^{r-1}+p^*c_2(E)h^{r-2}+\cdots+p^*c_r(E)=0
	\end{equation}
	的类 \(c_i(E)\in H^{2i}(B;\mathbb Z)\)。置 \(c_0(E)=1\)，\(i>r\) 时置 \(c_i(E)=0\)，并称 \(c(E)=1+c_1(E)+\cdots+c_r(E)\) 为总陈类。
\end{definition}

当 \(r=1\) 时，\(\mathbb P(E)=B\)、\(\lambda=E\)，关系变成 \(-c_1(E)+c_1(E)=0\)，故新定义延续上一节的线丛规范。正号排列来自 \(h=c_1(\lambda^*)\)：如果改用 \(c_1(\lambda)=-h\)，关系的系数就要交错变号。

\begin{theorem}\label{g2:thm:characteristic-existence-uniqueness}
	CW 复形上的复有限秩向量丛具有上述自然陈类。若 \(q:Y\to B\) 是使 \(q^*E=L_1\oplus\cdots\oplus L_r\) 的旗丛，记 \(x_j=c_1(L_j)\)，则
	\begin{equation}\label{g2:eq:chern-root-formula}
		q^*c_i(E)=\sum_{j_1<\cdots<j_i}x_{j_1}\cdots x_{j_i},\qquad
		q^*c(E)=\prod_{j=1}^r(1+x_j).
	\end{equation}
	自然性及这个分裂公式唯一确定全部 \(c_i\)。
\end{theorem}
\begin{proof}
	射影丛模基已经保证存在唯一的关系系数。若 \(f:A\to B\)，则拉回方块中的 \(\mathbb P(f^*E)\to\mathbb P(E)\) 将 \(\lambda\) 拉回为新的典范线丛，故也将 \(h\) 拉回为新的 \(h\)。拉回式\eqref{g2:eq:chern-projective-relation}并用新射影丛的模基唯一性，得到 \(c_i(f^*E)=f^*c_i(E)\)。

	为确定关系的系数，先设 \(E=L_1\oplus\cdots\oplus L_r\) 已分裂。在 \(\mathbb P(E)\) 上，丛
	\[
		\operatorname{Hom}_{\mathbb C}(\lambda,p^*E)
		=\bigoplus_{j=1}^r(\lambda^*\otimes p^*L_j)
	\]
	有一个处处非零截面：每点的直线包含映射 \(\lambda_\ell\hookrightarrow E_{p(\ell)}\)。\cref{g2:prop:euler-basic-properties}说明它的底层定向实丛 Euler 类为零；同一命题的直和公式及线丛的 \(c_1=e_{\mathrm{Th}}\) 给出
	\[
		0=\prod_{j=1}^r c_1(\lambda^*\otimes p^*L_j)
		=\prod_{j=1}^r(h+p^*x_j).
	\]
	这里使用的是已证明的 Euler 类乘法，并未预先假定一般陈类的乘法。展开右端，与唯一的首一射影关系比较，得到 \(c_i(E)\) 是 \(x_j\) 的第 \(i\) 个初等对称式。

	一般情形拉回到旗丛，用自然性及刚才的分裂计算得到式\eqref{g2:eq:chern-root-formula}。由\cref{g2:thm:splitting-principle}，\(q^*\) 单射，所以这些等式唯一确定原基空间上的类；任何其他自然赋类若满足同样的分裂公式，拉回后相同，原来也必相同。
\end{proof}

例如秩二时，形式根 \(x_1,x_2\) 给出 \(c_1=x_1+x_2\)、\(c_2=x_1x_2\)，射影丛关系为 \(h^2+c_1h+c_2=0\)。根是旗空间上的实际一次线丛类，未必分别下降到底空间；能下降的这些对称式已经由射影关系独立构造。这是“形式根”计算可以可靠使用的依据。

\subsection{Stiefel–Whitney 类及实射影丛关系}
\label{g2:subsec:2103:02}
对实秩 \(r\) 丛，射影纤维是 \(\mathbb RP^{r-1}\)，用模二系数后它同样有一列由一次类的幂组成的基。这里不需要给原丛或射影纤维选整数方向。

\begin{definition}\label{g2:def:stiefel-whitney-projective}
	设 \(E\to B\) 为实秩 \(r\ge1\) 向量丛，\(p:\mathbb P(E)\to B\) 为实射影丛，\(\lambda\) 为典范实线丛，\(t=w_1(\lambda)\)。唯一满足
	\begin{equation}\label{g2:eq:sw-projective-relation}
		t^r+p^*w_1(E)t^{r-1}+p^*w_2(E)t^{r-2}+\cdots+p^*w_r(E)=0
	\end{equation}
	的类 \(w_i(E)\in H^i(B;\mathbb F_2)\)，称为 \(E\) 的\term[Stiefel-Whitney-lei]{Stiefel--Whitney 类}[Stiefel--Whitney classes]。置 \(w_0=1\)、\(w_i=0\)（\(i>r\)），并记 \(w(E)=\sum_iw_i(E)\)。
\end{definition}

\begin{theorem}\label{g2:thm:stiefel-whitney-construction}
	上述类对 CW 基空间上的实丛自然。若旗丛拉回使 \(q^*E=L_1\oplus\cdots\oplus L_r\)，并记 \(a_j=w_1(L_j)\)，则
	\begin{equation}\label{g2:eq:sw-root-formula}
		q^*w(E)=\prod_{j=1}^r(1+a_j).
	\end{equation}
	自然性及该分裂公式唯一确定这些类；秩一时与\cref{g2:def:first-stiefel-whitney-line}一致。
\end{theorem}
\begin{proof}
	射影丛的模二模基保证关系系数存在唯一，射影拉回方块保证自然性。若 \(E\) 已分裂，\(\operatorname{Hom}_{\mathbb R}(\lambda,p^*E)\) 仍有处处非零的包含截面。它分裂成 \(\lambda^*\otimes p^*L_j\) 的直和。由式\eqref{g2:eq:line-mod-two-euler}及模二 Euler 直和公式，得到
	\[
		0=e_2(\operatorname{Hom}(\lambda,p^*E))
		=\prod_{j=1}^r(t+p^*a_j).
	\]
	与式\eqref{g2:eq:sw-projective-relation}比较唯一系数，即得分裂公式。一般丛再拉回实旗丛；\(q^*\) 在模二上同调上单射，所以相同论证给出唯一性。秩一的关系为 \(t+w_1(E)=0\)，正好恢复线丛类。
\end{proof}

\begin{proposition}\label{g2:prop:first-classes-determinants}
	对 CW 基空间上的实丛 \(E\) 与复丛 \(F\)，有
	\[
		w_1(E)=w_1(\det E),\qquad c_1(F)=c_1(\det F).
	\]
	实丛可定向当且仅当 \(w_1(E)=0\)。
\end{proposition}
\begin{proof}
	在相应旗丛上，行列式线丛是各线因子的张量积，因此其第一类是各根之和；这也是两个总类的一次项。拉回单射使等式在底空间成立。实丛方向的输运为过渡矩阵的行列式符号，恰好是 \(\det E\) 的方向输运，所以可定向性由\cref{g2:prop:real-line-first-class}判定。
\end{proof}

\begin{theorem}\label{g2:thm:top-class-euler}
	若 \(E\to B\) 为 CW 基空间上的实秩 \(r\) 丛，则 \(w_r(E)=e_2(E)\)。若 \(E\) 整数可定向，则 \(w_r(E)=\rho_2 e_{\mathrm{Th}}(E)\)。若 \(F\to B\) 为复秩 \(r\) 丛，则
	\[
		c_r(F)=e_{\mathrm{Th}}(F_{\mathbb R}),
	\]
	右侧采用复自然定向。
\end{theorem}
\begin{proof}
	先在旗丛上分裂成线丛。实情形中，模二 Euler 类是各线的 \(w_1\) 之积，恰为式\eqref{g2:eq:sw-root-formula}的最高项。复情形中，实丛的直和顺序按各复线的正向实二平面排列，正是复自然定向；Euler 直和公式使最高类为 \(x_1\cdots x_r\)，恰为 \(c_r\)。再由拉回单射得到原空间等式。整数定向丛的 Thom 类约化模二后是规范模二 Thom 类，零截面拉回与系数约化交换，得到中间断言。
\end{proof}

\subsection{Whitney 和公式与 Wu 公式}
\label{g2:subsec:2103:03}
两张丛直和后，在线丛分裂模型中只是把两组根放在一起。因此直和公式来自一个简单的乘法，但要让它在原基空间成立，需要同时分裂两张丛并保留拉回单射。

\begin{theorem}[Whitney 和公式]\label{g2:thm:whitney-sum}
	设 \(B\) 为 CW 复形。对其上的复有限秩丛 \(E,F\)，以及实有限秩丛 \(V,W\)，分别有
	\begin{equation}\label{g2:eq:whitney-total-formulas}
		c(E\oplus F)=c(E)c(F),\qquad w(V\oplus W)=w(V)w(W).
	\end{equation}
	两组类都在添上平凡直和因子时不变。自然性、线丛归一化、秩以上的消失及式\eqref{g2:eq:whitney-total-formulas}也唯一刻画各组类。
\end{theorem}
\begin{proof}
	先取 \(E\) 的旗丛，再取拉回 \(F\) 的旗丛。复合拉回仍单射，且两丛同时分裂，设根分别为 \(x_i,y_j\)。直和的根就是它们的合并，所以
	\[
		q^*c(E\oplus F)=\prod_i(1+x_i)\prod_j(1+y_j)
		=q^*(c(E)c(F)).
	\]
	单射性给出复情形。实旗丛及一次模二根给出完全相同的乘法论证。平凡丛分裂成第一类均为零的线丛，其总类为 \(1\)，从而得到稳定性。任何另一组满足所列性质的类，在旗丛上由线丛归一化与直和公式确定为同一个根乘积，故由单射性与本构造一致。
\end{proof}

短正合列 \(0\to E\to G\to F\to0\) 也满足这个公式，因为纤维内积给出 \(G\cong E\oplus E^\perp\)，而 \(E^\perp\to F\) 是丛同构。这里需要的是拓扑分裂，若原丛带有全纯等额外结构，不要求分裂保持那些结构。

\ProseParagraphBreak

现在完成第13章留下的 Thom 平方比较。记实秩 \(r\) 丛的模二 Thom 类为 \(U\)，用第20章同构定义 \(w_i^{\mathrm{Sq}}\) 使
\[
	\operatorname{Sq}^iU=p^*w_i^{\mathrm{Sq}}(E)\smile U.
\]
\cref{g2:prop:thom-square-properties}已经用平方的自然性与 Cartan 公式证明这组类的自然性、直和公式和秩以上消失。第20章提供的 Thom 同构及纤维积直和公式使同一论证适用于本节的基空间，余下只需核对线丛归一化。

\begin{theorem}\label{g2:thm:thom-square-comparison}
	对 CW 复形上的实有限秩丛 \(E\)，射影丛构造的 Stiefel--Whitney 类满足
	\begin{equation}\label{g2:eq:thom-sw-identity-proved}
		\operatorname{Sq}(U)=p^*w(E)\smile U.
	\end{equation}
	因此 \(w_i^{\mathrm{Sq}}(E)=w_i(E)\)。
\end{theorem}
\begin{proof}
	设 \(L\) 为实线丛，\(U\) 是它的一次模二 Thom 类。平方的最高项性质给出 \(\operatorname{Sq}^1U=U\smile U\)。记 \(j^*\) 为忘去相对条件，则零截面与盘丛投影互为同伦逆，故
	\[
		j^*U=p^*e_2(L).
	\]
	相对杯积与绝对—相对杯积相容，所以
	\[
		U\smile U=j^*U\smile U=p^*e_2(L)\smile U.
	\]
	Thom 同构的单射性及式\eqref{g2:eq:line-mod-two-euler}说明 \(w_1^{\mathrm{Sq}}(L)=e_2(L)=w_1(L)\)。其余正次项由秩一消失。于是这组平方类满足刚才唯一性定理的全部条件，必与射影丛构造相同。
\end{proof}

这一步把作用于总空间相对类的平方运算，识别为基空间上的示性类。对流形的对角嵌入应用它，便能把切丛类与第13章已经定义的 Wu 类联系起来。

\begin{theorem}[Wu 公式]\label{g2:thm:wu-formula}
	设 \(M\) 为连通闭光滑 \(n\) 维流形，\(v(M)\) 为\cref{g2:def:wu-class}按模二对偶定义的总 Wu 类。则
	\[
		w(\mathrm TM)=\operatorname{Sq}(v(M)).
	\]
	这个等式不要求 \(M\) 整数可定向。
\end{theorem}
\begin{proof}
	全部系数取 \(\mathbb F_2\)。先说明证明的目标：把对角线的 Thom 类平方化，一边读出法丛的 \(w\)，另一边用对偶基读出 Wu 类的平方。记 \(H=H^*(M;\mathbb F_2)\)，并以 \(\epsilon(a)\) 表示 \(a\) 的 \(n\) 次分量在 \([M]_2\) 上的求值。闭流形有有限 CW 同伦型，故 \(H\) 有限维；Poincaré 对偶给出齐次基 \(b_\alpha\) 及其对偶基 \(b^\alpha\)，满足
	\[
		\epsilon(b_\alpha b^\beta)=\delta_\alpha^\beta.
	\]
	对角嵌入 \(\Delta:M\hookrightarrow M\times M\) 的法丛同构于 \(\mathrm TM\)：商线性映射 \((v,w)\mapsto w-v\) 的核就是对角切空间。由管状邻域的 Thom 类定义推前 \(\Delta_!a\)，即先在邻域内取 \(a\) 与法丛 Thom 类的乘积，再经切除和忘去相对条件得到 \(H^{n+\deg a}(M\times M;\mathbb F_2)\) 中的类。第20章法向局部对偶的同一计算给出
	\[
		\langle\Delta_!a\smile z,[M\times M]_2\rangle
		=\langle a\smile\Delta^*z,[M]_2\rangle.
	\]
	令 \(D=\Delta_!1\)。由域系数 Künneth 同构和上式的完美配对，
	\begin{equation}\label{g2:eq:diagonal-dual-basis}
		D=\sum_\alpha b_\alpha\times b^\alpha.
	\end{equation}
	例如将右边与 \(x\times y\) 相乘后求值，得到 \(\sum_\alpha\epsilon(b_\alpha x)\epsilon(b^\alpha y)=\epsilon(xy)\)，恰与对角拉回相同，故这确是 \(D\)。

	自然性使平方与相对类到绝对类的映射交换；\cref{g2:thm:thom-square-comparison}于是给出
	\[
		\operatorname{Sq}(D)=\Delta_!w(\mathrm TM)
		=D\smile\operatorname{pr}_1^*w(\mathrm TM).
	\]
	最后一个等式是邻域内的杯积公式：\(\operatorname{pr}_1\) 限制在管状邻域后，与收缩到对角线再投影同伦，所以其拉回正是 Thom 同构中的基空间因子。

	在 \(H\otimes H\) 上施加线性映射 \(1\otimes\epsilon\)。对右边，式\eqref{g2:eq:diagonal-dual-basis}给出 \((1\otimes\epsilon)D=1\)，所以结果是 \(w(\mathrm TM)\)。对左边，Cartan 公式和外积相容性给出
	\[
		(1\otimes\epsilon)\operatorname{Sq}(D)
		=\sum_\alpha\operatorname{Sq}(b_\alpha)\,\epsilon(\operatorname{Sq}(b^\alpha)).
	\]
	Wu 类的定义逐次数合起来正是 \(\epsilon(\operatorname{Sq}(x))=\epsilon(v(M)x)\)。因此上式等于
	\[
		\operatorname{Sq}\left(\sum_\alpha b_\alpha\,\epsilon(v(M)b^\alpha)\right)
		=\operatorname{Sq}(v(M)),
	\]
	最后一步是对偶基展开。两边比较即得所求公式，全程只使用模二基本类。
\end{proof}

这就证明了\cref{g2:thm:wu-tangent-preview}的预告。低次时，\(w_1=v_1\)、\(w_2=v_2+\operatorname{Sq}^1v_1\)；后者额外的一项说明 Wu 类本身与切丛类一般不同。第13章中 \(\mathbb RP^2\) 与 \(\mathbb CP^2\) 的计算现在分别给出实际切丛的总类，而不再只是运算产生的候选表达。

\subsection{对偶、张量和复化}
\label{g2:subsec:2103:04}
对偶、张量与复化改变线丛的根，因此可以用同一个分裂过程计算。进行比较时须保留系数：复陈类位于整数上同调，实 Stiefel--Whitney 类位于模二上同调，二者之间使用约化映射 \(\rho_2\)。

\begin{proposition}\label{g2:prop:chern-operations}
	设 \(E,F\) 为 CW 基空间上的复丛，秩分别为 \(r,s\)，形式陈根为 \(x_i,y_j\)。则
	\[
		c_k(E^*)=c_k(\overline E)=(-1)^k c_k(E),\qquad
		c(E\otimes F)=\prod_{i,j}(1+x_i+y_j),
	\]
	第二式在同时分裂空间中计算，再由单射性理解于底空间。特别地，
	\[
		c_1(E\otimes F)=s\,c_1(E)+r\,c_1(F).
	\]
	若 \(L\) 为复线丛、\(\ell=c_1(L)\)，则
	\begin{equation}\label{g2:eq:chern-line-twist}
		c_k(E\otimes L)=\sum_{j=0}^k\binom{r-j}{k-j}c_j(E)\ell^{k-j}.
	\end{equation}
\end{proposition}
\begin{proof}
	在同时分裂空间，\(E^*\)、\(\overline E\) 的线因子第一类为 \(-x_i\)，因为线丛对偶与共轭都变号；第 \(k\) 个初等对称式遂乘以 \((-1)^k\)。张量积实际分裂为全部 \(L_i\otimes K_j\)，每个因子的第一类为 \(x_i+y_j\)，给出乘积公式及其一次项。

	对线丛扭转，把第 \(k\) 个初等对称式写为 \(e_k(x_1+\ell,\ldots,x_r+\ell)\)。从一项的 \(k\) 个因子中选 \(j\) 个贡献 \(x\)，其余贡献 \(\ell\)；固定那 \(j\) 个根后，还可从其余 \(r-j\) 个位置选出 \(k-j\) 个，故系数为 \(\binom{r-j}{k-j}\)。这证明式\eqref{g2:eq:chern-line-twist}。各等式都是对称式，且拉回单射，因此在底空间成立。
\end{proof}

实丛的对偶由内积与自身同构，所以 \(w(E^*)=w(E)\)。若两实丛分别分裂为根 \(a_i,b_j\) 的实线丛，则其张量积满足 \(w(E\otimes F)=\prod_{i,j}(1+a_i+b_j)\)。这些公式的乘积次数不同，不能把总陈类的整数系数式直接改名为总 Stiefel--Whitney 类的式子。

\begin{proposition}\label{g2:prop:real-complex-characteristic}
	对 CW 基空间上的复秩 \(r\) 丛 \(F\)，有
	\begin{equation}\label{g2:eq:complex-real-sw}
		w_{2i}(F_{\mathbb R})=\rho_2c_i(F),\qquad w_{2i+1}(F_{\mathbb R})=0.
	\end{equation}
	对实丛 \(E\) 的复化 \(E_{\mathbb C}=E\otimes_{\mathbb R}\mathbb C\)，有
	\begin{equation}\label{g2:eq:complexification-chern-mod-two}
		\rho_2c_i(E_{\mathbb C})=w_i(E)^2,\qquad 2c_{2j+1}(E_{\mathbb C})=0.
	\end{equation}
\end{proposition}
\begin{proof}
	在复旗丛上分裂 \(F\)。该拉回不仅在整数上同调单射，在模二上同调也单射，因为复射影纤维的幂基约化后仍是一组基。每条复线的底层实二平面定向，故 \(w_1=0\)，而\cref{g2:thm:top-class-euler}给出 \(w_2=\rho_2c_1\)。对各线因子应用 Whitney 公式，得 \(w(F_{\mathbb R})=\prod_i(1+\rho_2x_i)\)，比较次数即得第一组公式。

	映射 \(v+\mathrm i w\mapsto(v,w)\) 给出实丛同构 \((E_{\mathbb C})_{\mathbb R}\cong E\oplus E\)，于是
	\[
		\sum_i\rho_2c_i(E_{\mathbb C})=w(E)^2=\sum_iw_i(E)^2
	\]
	在模二系数下成立；最后一步的混合项成对相消。比较次数得到第二组的第一个式子。复化丛有规范的复线性同构 \(E_{\mathbb C}\cong\overline{E_{\mathbb C}}\)，由共轭的陈公式得 \(c_i=(-1)^ic_i\)，故奇下标项被 \(2\) 消去。
\end{proof}

例如实线丛复化后虽然只有一个陈类，这个类也未必为零：式\eqref{g2:eq:line-complexification-bockstein}精确把它写成方向类的整数 Bockstein。其模二约化为 \(\operatorname{Sq}^1w_1=w_1^2\)，与式\eqref{g2:eq:complexification-chern-mod-two}一致。复化保留了某些实方向信息，但把它从一次模二群移入二次整数扭群。

\subsection{Pontryagin 类及其性质}
\label{g2:subsec:2103:05}
实丛复化以后得到整数陈类。奇下标类可能只剩二阶扭信息，偶下标类则给出位于 \(4i\) 次的 Pontryagin 类。定义中的交错符号使定向实二平面的第一 Pontryagin 类成为 Euler 类的平方。

\begin{definition}\label{g2:def:pontryagin-classes}
	实有限秩丛 \(E\to B\) 的\term[Pontryagin-lei]{Pontryagin 类}[Pontryagin classes]定义为
	\begin{equation}\label{g2:eq:pontryagin-definition}
		p_i(E)=(-1)^i c_{2i}(E_{\mathbb C})\in H^{4i}(B;\mathbb Z),\qquad p(E)=1+p_1(E)+p_2(E)+\cdots.
	\end{equation}
	若 \(2i>\operatorname{rank}_{\mathbb R}E\)，则该项为零。
\end{definition}

Pontryagin\footnote{\OriginalName{Лев Семёнович Понтрягин}（Lev Semenovich Pontryagin，1908--1988），苏联数学家，研究代数拓扑、拓扑群与最优控制理论。}类对实丛定义，不要求它可定向。

\begin{proposition}\label{g2:prop:pontryagin-properties}
	对 CW 基空间上的实丛，Pontryagin 类自然且在加平凡直和因子后不变，并满足
	\[
		\rho_2p_i(E)=w_{2i}(E)^2.
	\]
	若 \(E\) 是定向实秩 \(2r\) 丛，则 \(p_r(E)=e_{\mathrm{Th}}(E)^2\)。若 \(F\) 为复丛，形式陈根为 \(x_j\)，则
	\begin{equation}\label{g2:eq:pontryagin-complex-underlying}
		p(F_{\mathbb R})=\prod_j(1+x_j^2),\qquad
		p_1(F_{\mathbb R})=c_1(F)^2-2c_2(F).
	\end{equation}
\end{proposition}
\begin{proof}
	自然性、稳定性来自复化与陈类的相应性质。模二约化时符号消失，再用式\eqref{g2:eq:complexification-chern-mod-two}即可。

	对定向秩 \(2r\) 丛 \(E\)，实同构 \((E_{\mathbb C})_{\mathbb R}\cong E\oplus E\) 所比较的两种方向须明确。若 \(e_1,\ldots,e_{2r}\) 为正向实基，复自然实方向是交错排列 \((e_1,\mathrm i e_1,\ldots,e_{2r},\mathrm i e_{2r})\)，直和方向则先放全部 \(e_j\)，再放全部 \(\mathrm i e_j\)。它们相差 \((-1)^{(2r)(2r-1)/2}=(-1)^r\)。最高陈类等于复方向 Euler 类，因而
	\[
		c_{2r}(E_{\mathbb C})=(-1)^re_{\mathrm{Th}}(E\oplus E)=(-1)^re_{\mathrm{Th}}(E)^2.
	\]
	代入定义便得 \(p_r=e_{\mathrm{Th}}^2\)。

	最后，复线性同构 \((F_{\mathbb R})_{\mathbb C}\cong F\oplus\overline F\) 可在每根纤维由复结构的 \(\pm\mathrm i\) 特征子空间给出。因此其陈总类在分裂空间为
	\[
		\prod_j(1+x_j)(1-x_j)=\prod_j(1-x_j^2).
	\]
	取偶下标项再乘 \((-1)^i\)，得到式\eqref{g2:eq:pontryagin-complex-underlying}。一次项利用 \(\sum_jx_j^2=(\sum_jx_j)^2-2\sum_{i<j}x_ix_j\) 化成所列公式。
\end{proof}

整数 Pontryagin 类的直和关系还要保留奇陈类的乘积。直接从陈类的 Whitney 公式分开偶、奇下标，可得
\begin{equation}\label{g2:eq:pontryagin-sum-torsion}
	\begin{split}
		p_k(E\oplus F)={}&\sum_{i+j=k}p_i(E)p_j(F)\\
		&+(-1)^k\!\sum_{\substack{a+b=2k\\a,b\text{ 为奇数}}}c_a(E_{\mathbb C})c_b(F_{\mathbb C}).
	\end{split}
\end{equation}
第二行被 \(2\) 消去，但未必等于零。因此
\[
	p(E\oplus F)=p(E)p(F)\quad\text{于 }H^*(B;\mathbb Z[1/2])
\]
总是成立；在整数上同调中，若相关群没有二阶扭，或者相关奇陈乘积全部消失，也可使用同一公式。特别是两个丛都来自复丛的底层实丛时，上面根的成对表达使这些奇项确实为零。

\begin{example}\label{g2:exm:pontryagin-whitney-torsion}
	令 \(L\to\mathbb RP^4\) 为典范实线丛，\(a=w_1(L)\)，\(t=c_1(L_{\mathbb C})=\widetilde\beta_2a\)。由第13章 \(\rho_2\widetilde\beta_2=\operatorname{Sq}^1\) 及最高平方性质，\(\rho_2t=a^2\)，所以 \(\rho_2(t^2)=a^4\ne0\)，从而 \(t^2\ne0\)。每条实线丛的 Pontryagin 类在正次全为零，但
	\[
		p_1(L\oplus L)=-c_2(L_{\mathbb C}\oplus L_{\mathbb C})=-t^2\ne0.
	\]
	这里整个差正是式\eqref{g2:eq:pontryagin-sum-torsion}的奇下标乘积。换成实或有理系数后 \(t\) 消失，便看不到这个例子中的区别。
\end{example}

至此，三个系列的次数与系数可以并列为
\[
\begin{array}{c|c|c|c}
	\text{丛的类型}&\text{类}&\text{所在群}&\text{构造所用数据}\\\hline
	\text{实秩 }r&w_i&H^i(B;\mathbb F_2)&\text{实射影丛关系}\\
	\text{复秩 }r&c_i&H^{2i}(B;\mathbb Z)&\text{复射影丛关系}\\
	\text{实秩 }r&p_i&H^{4i}(B;\mathbb Z)&(-1)^ic_{2i}(E_{\mathbb C})
\end{array}
\]
下一节将这些类作用于实际切丛和结构群约化；下一章再在实、复系数中用曲率给出代表。整数类及其扭项已经在拓扑构造中确定，曲率比较须与这里的规范一致。
% 纲要条目排为编号小节标题，并进入目录与书签；说明段仍保留于纲要注释中。

% !TeX root = ../../Book2.tex
% 纲要标题：### 21.4 几何计算与约束
% 纲要位置：几何学计划纲要.md，第 1827 行。
% 正文状态：未撰写；本文件目前是框架。

% 纲要内容开始
% 
% * 射影空间切丛
% * 球丛和低维切丛
% * 特征数与结构群约化
% * spin 条件与更高障碍的区别
% 
% 纲要内容结束

\section{几何计算与约束}
\label{g2:sec:2104}

待撰写

\subsection{射影空间切丛}
\label{g2:subsec:2104:01}
待撰写

\subsection{球丛和低维切丛}
\label{g2:subsec:2104:02}
待撰写

\subsection{特征数与结构群约化}
\label{g2:subsec:2104:03}
待撰写

\subsection{spin 条件与更高障碍的区别}
\label{g2:subsec:2104:04}
待撰写
% 纲要条目排为编号小节标题，并进入目录与书签；说明段仍保留于纲要注释中。

% !TeX root = ../../Book2.tex
% 纲要标题：### 21.5 Grassmann 流形的胞腔、Schubert 类与相交计算
% 纲要位置：几何学计划纲要.md，第 1834 行。
% 正文状态：初稿；尚未正式审核。

% 纲要内容开始
% 
% * 固定完全旗与 Schubert 胞腔的秩条件
% * 行阶梯形坐标、偶维胞腔与整数上同调基
% * 通用子丛、商丛和特殊 Schubert 类
% * 利用旗的移动及相交配对计算低阶乘法
% * $\operatorname{Gr}(2,4)$ 的杯积计算与四条直线问题
% 
%   在复数域模型上计算特殊除子类的四次幂，并由一般位置保证几何计数；完整的方程消去、重数与退化核验放在 V.27.5。本节不以一般 Schubert 簇光滑为前提，也不追加完整旗簇表示论。
% 
% * 与 V.27.5 中 Chow 环及 Plücker 方程计算的比较
% 
% ---
% 
% 纲要内容结束

\section{Grassmann 流形的胞腔、Schubert 类与相交计算}
\label{g2:sec:2105}

示性类已经能把向量丛的运算变成上同调计算。现在反过来问：这些上同调类是否有具体的几何代表，它们的杯积又在数什么？复 Grassmann 流形提供了一个可以同时看见两面的模型。一个子空间与固定旗相交的维数条件给出 Schubert 簇；通用商丛的陈类恰好代表其中最简单的一类条件。最后，\(\operatorname{Gr}(2,4)\) 上一个四次杯积将回答三维射影空间中四条直线的公共截线问题。

以下记 \(G=\operatorname{Gr}(k,n)=\operatorname{Gr}_k(\mathbb C^n)\)，其中 \(1\le k<n\)，并置 \(q=n-k\)、\(N=kq\)。它是复维数 \(N\) 的紧光滑流形，始终采用复结构给出的实定向及整数上同调。第1章的\cref{g2:thm:grassmann-schubert-cw}已经证明胞腔的特征映射及闭包关系；这里恢复坐标含义，再用它们构造同调类并计算乘法。有关胞腔、相交与通用丛的经典叙述，可参阅\cite[第1章第5节，The Cell Decomposition、The Schubert Calculus、Universal Bundles]{GriffithsHarris1994}。

\subsection{固定完全旗与 Schubert 胞腔的秩条件}
\label{g2:subsec:2105:01}
固定一条\term[wanquanqi]{完全旗}[complete flag]
\[
	0=F_0\subset F_1\subset\cdots\subset F_n=\mathbb C^n,
	\qquad\dim_{\mathbb C}F_i=i.
\]
选取与它相容的基 \(e_1,\ldots,e_n\)，使 \(F_i=\operatorname{span}(e_1,\ldots,e_i)\)。对 \(V\in G\)，整数 \(\dim(V\cap F_i)\) 从零增到 \(k\)，每步最多增加一，所以恰有 \(k\) 个跳跃位置
\[
	1\le d_1<d_2<\cdots<d_k\le n.
\]
第1章用这些递增指标编号胞腔。为了让指标直接记录余维，现在改用 \(\lambda_i=q+i-d_i\)。于是 \(q\ge\lambda_1\ge\cdots\ge\lambda_k\ge0\)。这样的整数列称为放在 \(k\times q\) 矩形中的一个分拆，记其大小为 \(|\lambda|=\sum_i\lambda_i\)。它可画成左对齐的方格行，行长依次为 \(\lambda_1,\ldots,\lambda_k\)；方格总数将成为复余维。

\begin{definition}\label{g2:def:schubert-partition-locus}
	给定上述旗与分拆，令 \(d_i=q+i-\lambda_i\)。具有跳跃位置恰为 \(d_1,\ldots,d_k\) 的子空间组成 Schubert 胞腔 \(C_\lambda(F_\bullet)\)。其闭包
	\begin{equation}\label{g2:eq:schubert-rank-conditions}
		\Omega_\lambda(F_\bullet)=\{V\in G:\dim(V\cap F_{d_i})\ge i,\ 1\le i\le k\}
	\end{equation}
	称为相对于该旗的\term[Schubert-cu]{Schubert 簇}[Schubert variety]。
\end{definition}

只写 \(\dim(V\cap F_{d_i})=i\) 还不足以定义开胞腔：还要排除跳跃提前发生。因此等价的精确条件是同时要求
\[
	\dim(V\cap F_{d_i})=i,\qquad\dim(V\cap F_{d_i-1})=i-1.
\]
例如 \(k=1\) 时，\(C_{(j)}=\mathbb P(F_{n-j})\setminus\mathbb P(F_{n-j-1})\)，而 \(\Omega_{(j)}=\mathbb P(F_{n-j})\) 是复射影空间中的余维 \(j\) 线性子空间。闭包中的不等式允许跳跃提前发生；它们比开胞腔的等式更适合描述极限。

\ProseParagraphBreak

这些条件也确实是矩阵的秩条件。若 \(V\) 由一个满秩 \(n\times k\) 矩阵 \(A\) 的列张成，则
\[
	\dim(V\cap F_d)=k-\operatorname{rank}(\mathbb C^k\xrightarrow{A}\mathbb C^n\longrightarrow\mathbb C^n/F_d).
\]
因而 \(\dim(V\cap F_d)\ge i\) 等价于最后 \(n-d\) 行组成的矩阵秩不超过 \(k-i\)，可由全部 \((k-i+1)\) 阶子式为零表达。这里条件与所选 \(V\) 的基无关，因为更换基只在右侧乘一个可逆矩阵。它说明 Schubert 簇在 Grassmann 坐标中由多项式方程给出，也说明闭集中的奇点为何可能出现：一个秩条件的不同阶退化会在同一闭包中相遇。

\subsection{行阶梯形坐标、偶维胞腔与整数上同调基}
\label{g2:subsec:2105:02}
在 \(C_\lambda\) 内，从右向左消元得到唯一的基
\begin{equation}\label{g2:eq:schubert-echelon-coordinates}
	v_i=e_{d_i}+\sum_{\substack{a<d_i\\a\notin\{d_1,\ldots,d_{i-1}\}}}z_{ia}e_a,
	\qquad 1\le i\le k.
\end{equation}
第 \(i\) 个向量有 \(d_i-i=q-\lambda_i\) 个自由复系数，因而
\[
	C_\lambda\cong\mathbb C^{N-|\lambda|},\qquad
	\Omega_\lambda=\coprod_{\mu_i\ge\lambda_i\ \text{对所有 }i}C_\mu.
\]
胞腔的实维数是 \(2(N-|\lambda|)\)。例如在 \(\operatorname{Gr}(2,4)\) 中，\(\lambda=(1,0)\) 对应 \((d_1,d_2)=(2,4)\)，可以把一组基作为矩阵的行写成
\[
	\begin{pmatrix}z_{11}&1&0&0\\z_{21}&0&z_{23}&1\end{pmatrix}.
\]
三个自由复数给出六维实胞腔；分拆只有一个方格，故它在八维实 Grassmann 流形中余维二。

\begin{proposition}\label{g2:prop:schubert-integral-basis}
	复 Grassmann 流形 \(G=\operatorname{Gr}(k,n)\) 的整数同调与上同调没有扭元，奇数次群均为零。每个 \(\Omega_\lambda\) 有一个由其开胞腔复定向确定的类
	\[
		[\Omega_\lambda]\in H_{2(N-|\lambda|)}(\Omega_\lambda;\mathbb Z).
	\]
	把这些类推入 \(G\)，得到各次同调的基；其 Poincar\'e 对偶
	\begin{equation}\label{g2:eq:schubert-class-definition}
		\sigma_\lambda=\mathrm{PD}_G[\Omega_\lambda]\in H^{2|\lambda|}(G;\mathbb Z)
	\end{equation}
	称为\term[Schubert-lei]{Schubert 类}[Schubert class]，并构成整数上同调的加法基。这些类与所选完全旗无关。
\end{proposition}
\begin{proof}
	由\cref{g2:thm:grassmann-schubert-cw}，上述集合给出有限 CW 结构，且全部胞腔的实维数为偶数。胞腔微分降低一次，所以其源或靶必为零，全部微分为零。于是同调在每个偶数维自由生成于该维胞腔，在奇数维为零。万有系数定理中的 Ext 项也为零，得到对应的上同调结论。

	\(\Omega_\lambda\) 是一个子复形，唯一的最高维胞腔为 \(C_\lambda\)，其余胞腔的实维数至少低二。故其最高维同调为 \(\mathbb Z\)，到最高维胞腔相对同调的映射为同构。选择在 \(C_\lambda\) 上给出复定向生成元的唯一类，便定义了 \([\Omega_\lambda]\)。包含诱导的胞腔链映射将这个生成元送到 \(G\) 中同一胞腔的生成元，所以所有这类类组成同调基。这里的构造只用胞腔复形，并未要求整个闭包是流形。

	\(G\) 是闭定向流形，\cref{g2:cor:closed-poincare-duality}将这个同调基送到上同调基。任意两条完全旗由某个 \(g\in\operatorname{GL}_n(\mathbb C)\) 联系。该群道路连通：极分解先把正定因子收缩到恒等，酉因子再逐个缩小其特征值的辐角。因此可取从恒等到 \(g\) 的道路，给出 \(G\) 上恒等映射与 \(V\mapsto gV\) 的同伦。后一个映射保持各开胞腔的复定向，并把对应的 Schubert 簇相互送到，所以它们给出相同同调类及对偶类。
\end{proof}

几何相交时，还可以用一个实际的闭光滑流形来代表这个类。令 \(\widetilde\Omega_\lambda\) 由所有旗
\[
	0=W_0\subset W_1\subset\cdots\subset W_k,
	\qquad\dim W_i=i,\quad W_i\subset F_{d_i}
\]
组成，并令 \(\rho_\lambda(W_\bullet)=W_k\)。已选 \(W_{i-1}\) 后，选择 \(W_i\) 等价于在 \(F_{d_i}/W_{i-1}\) 中选一条线，因而每一步都是复射影丛，纤维为 \(\mathbb{CP}^{d_i-i}\)。这说明 \(\widetilde\Omega_\lambda\) 是紧光滑复流形，复维数为 \(\sum_i(d_i-i)=N-|\lambda|\)。

\ProseParagraphBreak

映射 \(\rho_\lambda\) 的像恰为 \(\Omega_\lambda\)：给定满足秩条件的 \(V\)，可在嵌套的 \(V\cap F_{d_i}\) 中逐次选择 \(W_i\)。若 \(V\in C_\lambda\)，则维数迫使 \(W_i=V\cap F_{d_i}\)，选择唯一；在这个开集上，常秩核空间随矩阵坐标光滑且复解析地变化。因此 \(\rho_\lambda\) 在开胞腔上是保持复定向的同构，而
\begin{equation}\label{g2:eq:schubert-resolution-class}
	(\rho_\lambda)_*[\widetilde\Omega_\lambda]=[\Omega_\lambda].
\end{equation}
最后一个等式可限制到最高维胞腔核对：两边都给出其正生成元，上一命题的唯一性使它们相等。这个旗模型将闭包低维处的多重选择留在纤维中，同时为相交理论提供一个光滑源。

\subsection{通用子丛、商丛和特殊 Schubert 类}
\label{g2:subsec:2105:03}
在 \(G\) 上有通用子丛 \(S\) 与商丛 \(Q\)：它们在 \(V\) 处的纤维分别为 \(V\) 和 \(\mathbb C^n/V\)。这是第19章典范丛的有限维模型，构成短正合列
\begin{equation}\label{g2:eq:grassmann-tautological-sequence}
	0\longrightarrow S\longrightarrow\underline{\mathbb C}^{\,n}\longrightarrow Q\longrightarrow0.
\end{equation}
选择 Hermitian 内积后可把 \(Q\) 识别为 \(S^\perp\)，所以 Whitney 和公式给出 \(c(S)c(Q)=1\)。这个等式已经把两张丛的陈类联系起来；以下识别它们所记录的几何条件。

\begin{theorem}\label{g2:thm:special-schubert-chern}
	在 \(G=\operatorname{Gr}(k,n)\) 上，令 \(q=n-k\)。对 \(1\le j\le q\)，有
	\begin{equation}\label{g2:eq:special-schubert-quotient}
		\sigma_j:=\sigma_{(j,0,\ldots,0)}=c_j(Q).
	\end{equation}
	特别地，\(\sigma_1=c_1(Q)=c_1(S^*)=c_1(\det S^*)\)。
\end{theorem}
\begin{proof}
	取 \(A=F_{q+1-j}\)，记 \(a=\dim A=q+1-j\)。对于分拆 \((j,0,\ldots,0)\)，除第一个条件外，其余条件由维数不等式自动成立。因此 \(\Omega_j=\{V:V\cap A\ne0\}\)。为把这个非零交条件写成一个横截零集，在 \(G\times\mathbb P(A)\) 上记 \(\ell\) 为从第二因子拉回的重言线丛，考虑秩 \(q\) 丛 \(\ell^*\otimes Q\) 的截面
	\[
		s(V,L):L\hookrightarrow\mathbb C^n\longrightarrow\mathbb C^n/V.
	\]
	其零集为 \(Z=\{(V,L):L\subset A\cap V\}\)。固定 \(L\) 后，\(V/L\) 可任取 \(\mathbb C^n/L\) 中的 \((k-1)\) 维子空间，所以 \(Z\) 是 \(\mathbb P(A)\) 上的光滑 Grassmann 丛。

	还可直接核对截面横截：在 \(L\subset V\) 处，\(G\) 的切空间由图坐标识别为 \(\operatorname{Hom}_{\mathbb C}(V,\mathbb C^n/V)\)。截面在这些方向的微分，除一个无关紧要的负号外，就是限制到 \(L\) 的映射
	\[
		\operatorname{Hom}(V,\mathbb C^n/V)\longrightarrow\operatorname{Hom}(L,\mathbb C^n/V),
	\]
	它满射，因为 \(L\) 上的线性映射可沿任意补空间延拓。故 \(s\) 横截于零截面。由\cref{g2:thm:top-class-euler}以及\cref{g2:thm:euler-transverse-zero-duality}，\(Z\) 的对偶类是 \(c_q(\ell^*\otimes Q)\)，定向就是复定向。

	令 \(h=c_1(\ell^*)\)。分裂原理把 \(Q\) 的形式根记为 \(x_1,\ldots,x_q\)，从而
	\begin{equation}\label{g2:eq:schubert-incidence-top-chern}
		c_q(\ell^*\otimes Q)=\prod_{b=1}^q(h+x_b)
		=\sum_{i=0}^q h^{q-i}c_i(Q).
	\end{equation}
	这是在旗丛拉回后验证的对称多项式恒等式；拉回单射使它在原空间成立。沿 \(\mathbb P(A)=\mathbb{CP}^{a-1}\) 求值时，只有 \(h^{a-1}=h^{q-j}\) 留下系数一，所以式\eqref{g2:eq:schubert-incidence-top-chern}给出 \(c_j(Q)\)。

	把这个系数提取写成同调推前，便可直接使用已证的乘积求值公式。记投影为 \(\pi:G\times\mathbb P(A)\to G\)。对任意 \(\beta\in H^{2(N-j)}(G;\mathbb Z)\)，顶陈类与零点类的比较给出
	\[
	\begin{aligned}
		\langle\beta,\pi_*[Z]\rangle
		&=\langle\pi^*\beta,[Z]\rangle\\
		&=\left\langle\pi^*\beta\smile\sum_i h^{q-i}c_i(Q),[G\times\mathbb{CP}^{a-1}]\right\rangle\\
		&=\langle\beta\smile c_j(Q),[G]\rangle.
	\end{aligned}
	\]
	全部次数均为偶数，交换因子没有符号。另一方面，\(\pi|_Z\) 的像为 \(\Omega_j\)，在其开胞腔上 \(V\cap A\) 恰为一维，故该处投影是保持复定向的同构。最高维胞腔的归一化给出 \(\pi_*[Z]=[\Omega_j]\)。\(G\) 的整数上同调自由，Poincar\'e 配对非退化，因此上述全部求值确定 \(\sigma_j=c_j(Q)\)。最后对式\eqref{g2:eq:grassmann-tautological-sequence}取一次陈类，并用对偶与行列式的公式，得到 \(\sigma_1\) 的其余表达。
\end{proof}

这个证明的几何步骤是先把“至少有一条交线”改成“连这条线也一起选出”。增加 \(\mathbb P(A)\) 参数后，退化条件成为一张向量丛的横截截面零集；随后沿这些参数求值，消去辅助选择，得到所求陈类。

\subsection{利用旗的移动及相交配对计算低阶乘法}
\label{g2:subsec:2105:04}
同一个 Schubert 类可以由许多不同旗给出的簇代表。移动旗的用途，是让相交发生在光滑开胞腔中并且横截；这时每个复交点贡献 \(+1\)，杯积便成为实际计数。

\begin{proposition}\label{g2:prop:schubert-moving-flags}
	设分拆 \(\lambda^{(1)},\ldots,\lambda^{(m)}\) 满足 \(\sum_b|\lambda^{(b)}|=N\)。可以选择处于一般位置的完全旗 \(F_\bullet^{(b)}\)，使各 \(\Omega_{\lambda^{(b)}}(F_\bullet^{(b)})\) 的公共交集是有限个点，全部位于各自开胞腔中且横截，并且
	\[
		\#\bigcap_{b=1}^m\Omega_{\lambda^{(b)}}(F_\bullet^{(b)})
		=\left\langle\prod_{b=1}^m\sigma_{\lambda^{(b)}},[G]\right\rangle.
	\]
	这里一般位置可取为旗参数空间中的一个开稠密集合；交集允许为空。
\end{proposition}
\begin{proof}
	群 \(\operatorname{GL}_n(\mathbb C)\) 在 \(G\) 上传递作用。作用映射 \((g,V)\mapsto gV\) 是淹没：在图坐标 \(\mathrm T_VG=\operatorname{Hom}(V,\mathbb C^n/V)\) 中，矩阵的一阶变化可以任意指定这张线性映射。因此将一个固定光滑胞腔随 \(g\) 移动，\cref{g2:thm:parametric-transversality}保证对一般参数，它横截于任意另一个固定光滑子流形。

	每个 Schubert 簇只有有限个胞腔。依次移动这些旗，同时对全部胞腔组合应用参数横截性，就使任意组合的相交都横截。若所有胞腔都取开胞腔，预期复维数为 \(N-\sum_b|\lambda^{(b)}|=0\)。若至少一个取闭包中的低维胞腔，其复维数至少再低一，预期维数变成负数，横截性迫使交集为空。所以全部交点落在开胞腔中。簇的闭性与 \(G\) 的紧性使总交集紧，而零维横截交集离散，故只有有限点。

	用式\eqref{g2:eq:schubert-resolution-class}的光滑旗模型代表各同调类。其乘积映射到 \(G^m\)，公共交点对应于这个映射对对角子流形的逆像；由于交点都在开胞腔中，那里各旗模型与 Schubert 簇相同，逆像仍然横截。将\cref{g2:thm:geometric-intersection-pairing}用于对角子流形，再用帽积的自然性及各个基本类的推前，所得求值就是所列杯积在 \([G]\) 上的值。这也把带奇点闭包的计数归结为闭光滑源与光滑对角的相交。

	在交点附近，所有切空间和法向商空间都是复向量空间，横截比较映射为复线性同构，其实行列式为复行列式绝对值的平方，严格为正。因此每个点的符号为 \(+1\)。最后，光滑旗模型紧，横截性在参数的小变动下保持；同时交点避开各闭包的低维部分，这一避开条件也由紧性保持。故所选好参数构成开集，参数横截性给出的稠密性则使它成为开稠密集。
\end{proof}

为了从这些计数读出一个类的系数，需要知道哪两个 Schubert 类相互对偶。定义矩形中的补分拆
\[
	\lambda^c=(q-\lambda_k,q-\lambda_{k-1},\ldots,q-\lambda_1).
\]
它把矩形中余下的方格旋转半周再左对齐，满足 \(|\lambda|+|\lambda^c|=N\)。

\begin{proposition}\label{g2:prop:schubert-complement-pairing}
	若 \(|\lambda|+|\mu|=N\)，则
	\begin{equation}\label{g2:eq:schubert-complement-pairing}
		\langle\sigma_\lambda\smile\sigma_\mu,[G]\rangle
		=\begin{cases}1,&\mu=\lambda^c,\\0,&\mu\ne\lambda^c.\end{cases}
	\end{equation}
\end{proposition}
\begin{proof}
	采用相反的两条坐标旗
	\[
		F_d=\operatorname{span}(e_1,\ldots,e_d),\qquad
		F'_d=\operatorname{span}(e_{n-d+1},\ldots,e_n).
	\]
	若某个 \(i\) 满足 \(\lambda_i+\mu_{k+1-i}>q\)，则相应两个旗子空间的维数和不超过 \(n\)，且它们在相反坐标旗中不相交。然而若 \(V\) 同时满足两组秩条件，它在这两个不相交子空间中分别至少有 \(i\) 维和 \(k+1-i\) 维，总计至少 \(k+1\) 维，矛盾。因此交集非空要求全部 \(\lambda_i+\mu_{k+1-i}\le q\)。将这些不等式求和，左边为 \(N=kq\)，所以每一项都必须取等号，即 \(\mu=\lambda^c\)。

	在补分拆情形，记 \(d_i=q+i-\lambda_i\)。上面的两子空间分别是 \(F_{d_i}\) 与 \(\operatorname{span}(e_{d_i},\ldots,e_n)\)，交集恰为 \(\mathbb Ce_{d_i}\)。由于 \(V\) 在二者中的维数之和至少为 \(k+1\)，它必须包含 \(e_{d_i}\)。对全部 \(i\) 使用此论证，唯一可能的交点便是
	\[
		V_0=\operatorname{span}(e_{d_1},\ldots,e_{d_k}).
	\]
	它确实满足两组条件，并位于两个开胞腔中。

	最后核对横截性。取补空间由非枢轴坐标向量张成，图坐标中的切向量为 \(A:V_0\to\mathbb C^n/V_0\)。由式\eqref{g2:eq:schubert-echelon-coordinates}，第一片的切方向是 \(A(e_{d_i})\) 在非枢轴行 \(a<d_i\) 的系数；相反旗的第二片恰好允许非枢轴行 \(a>d_i\) 的系数。这两组参数不交且合起来覆盖全部图坐标，所以切空间互补。复定向使唯一交点的符号为正，由前一命题的相交解释得到式\eqref{g2:eq:schubert-complement-pairing}。
\end{proof}

因此，若某个齐次类写成 \(\alpha=\sum_\lambda a_\lambda\sigma_\lambda\)，则
\[
	a_\lambda=\langle\alpha\smile\sigma_{\lambda^c},[G]\rangle.
\]
一个乘积的系数就能通过增加一项互补几何条件来求出。下面在最小的非射影空间例子中把这些条件逐项算完。

\subsection{\texorpdfstring{\(\operatorname{Gr}(2,4)\) 的杯积计算与四条直线问题}{Gr(2,4) 的杯积计算与四条直线问题}}
\label{g2:subsec:2105:05}
把 \(\operatorname{Gr}(2,4)\) 的点 \(V\) 看成 \(\mathbb{CP}^3\) 中的射影直线 \(\mathbb P(V)\)。固定射影旗“点 \(P\) 位于直线 \(L\)，而 \(L\) 位于平面 \(H\)”后，六个分拆的含义如下。
\begin{center}
	\begin{tabular}{ccl}
		\hline
		分拆 & 复余维 & 所参数化的射影直线 \\
		\hline
		\(\varnothing\) & \(0\) & 全部直线 \\
		\((1)\) & \(1\) & 与 \(L\) 相交 \\
		\((2)\) & \(2\) & 经过 \(P\) \\
		\((1,1)\) & \(2\) & 包含于 \(H\) \\
		\((2,1)\) & \(3\) & 经过 \(P\) 且包含于 \(H\) \\
		\((2,2)\) & \(4\) & 等于给定直线 \(L\) \\
		\hline
	\end{tabular}
\end{center}
其中 \(\Omega_2\) 与 \(\Omega_{11}\) 都同构于 \(\mathbb{CP}^2\)，\(\Omega_{21}\) 是一束直线组成的 \(\mathbb{CP}^1\)，而 \(\Omega_{22}\) 是一点。\(\Omega_1\) 的闭包却不必光滑，因而前面的光滑代表与移动旗论证在此已有实际用途。为简化式子，省略下标分拆中的逗号。

\begin{proposition}\label{g2:prop:gr24-schubert-products}
	\(G=\operatorname{Gr}(2,4)\) 的非零正次乘积由以下等式及交换律决定：
	\begin{equation}\label{g2:eq:gr24-schubert-products}
	\begin{gathered}
		\sigma_1^2=\sigma_2+\sigma_{11},\qquad
		\sigma_1\sigma_2=\sigma_1\sigma_{11}=\sigma_{21},\\
		\sigma_1\sigma_{21}=\sigma_{22},\qquad
		\sigma_2^2=\sigma_{11}^2=\sigma_{22},\qquad
		\sigma_2\sigma_{11}=0.
	\end{gathered}
	\end{equation}
	特别地，\(\sigma_1^3=2\sigma_{21}\)、\(\sigma_1^4=2\sigma_{22}\)，且 \(\langle\sigma_{22},[G]\rangle=1\)。
\end{proposition}
\begin{proof}
	先求余维互补的乘积。经过两个不同点的射影直线唯一，所以 \(\sigma_2^2=\sigma_{22}\)。同时位于两个不同平面的直线是它们的交线，所以 \(\sigma_{11}^2=\sigma_{22}\)。点若不在给定平面中，就没有既经过该点又包含于该平面的直线，故 \(\sigma_2\sigma_{11}=0\)。这也与补分拆配对相符：\((2)\) 与 \((1,1)\) 各自为自身的补。

	在平面 \(H\) 中经过点 \(P\) 的直线组成一束。另取一条不包含于 \(H\) 的直线 \(L\)，并使交点 \(Q=L\cap H\) 不等于 \(P\)。这束直线中恰有连接 \(P,Q\) 的一条与 \(L\) 相交，所以 \(\sigma_1\sigma_{21}=\sigma_{22}\)。一般位置保证交点横截，复定向使系数为一。

	六次群只由 \(\sigma_{21}\) 生成。若 \(P\notin L\)，经过 \(P\) 且与 \(L\) 相交的直线组成平面 \(\langle P,L\rangle\) 中的直线束，正是某条旗的 \(\Omega_{21}\)。在一般位置下相交横截于光滑部分，因而 \(\sigma_1\sigma_2=\sigma_{21}\)。同理，平面 \(H\) 内与一条横穿 \(H\) 的直线相交，等价于经过 \(L\cap H\)，所以 \(\sigma_1\sigma_{11}=\sigma_{21}\)。也可把这两个等式分别与 \(\sigma_1\) 配对，由刚才的唯一交点计数及非退化性确定系数。

	最后写 \(\sigma_1^2=a\sigma_2+b\sigma_{11}\)。分别乘 \(\sigma_2\) 与 \(\sigma_{11}\) 并在 \([G]\) 上求值，前面已经算得的乘积给出 \(a=b=1\)。其余两式由继续乘 \(\sigma_1\) 得到。次数超过八的类为零，因此没有其他正次乘积需要决定。
\end{proof}

这里的两种余维二条件不能混为一类：经过一个点与位于一个平面都给出一份 \(\mathbb{CP}^2\)，但它们在 \(G\) 中有不同同调类。它们各自的自交为一，彼此的交数为零，正是四次群秩为二的几何表现。

\ProseParagraphBreak

也可以将乘法表写成陈类的环表示。记 \(h=\sigma_1\)、\(s=\sigma_{11}\)。取一个非零线性函数 \(\alpha:\mathbb C^4\to\mathbb C\)，它限制到各 \(V\) 上给出 \(S^*\) 的截面。零点集是 \(V\subset\ker\alpha\)，即 \(\Omega_{11}\)。在这个零集处，图坐标微分为 \(A\mapsto\alpha\circ A\)，满射到 \(V^*\)，所以截面横截；于是
\[
	s=c_2(S^*),\qquad c(S)=1-h+s,
	\qquad c(Q)=1+h+(h^2-s).
\]
其中 \(c_2(Q)=\sigma_2=h^2-s\)，与特殊 Schubert 类公式一致。乘法表等价于
\begin{equation}\label{g2:eq:gr24-cohomology-ring}
	H^*(\operatorname{Gr}(2,4);\mathbb Z)
	\cong\mathbb Z[h,s]/(h^3-2hs,\ s^2-h^2s),
	\qquad\deg h=2,\quad\deg s=4.
\end{equation}
为核对没有漏掉整数关系，先用两条关系推出 \(h^3s=0\)、\(s^3=0\)、\(h^5=0\)：例如 \(h^3s=2hs^2=2h^3s\)，相减即得第一式。故商环由 \(1,h,h^2,s,hs,s^2\) 加法生成；这些元素的像分别为 \(1,\sigma_1,\sigma_2+\sigma_{11},\sigma_{11},\sigma_{21},\sigma_{22}\)，构成同调所给自由基的另一组整数基，所以映射既满又单。

\begin{theorem}\label{g2:thm:four-lines-two-transversals}
	在 \(\mathbb{CP}^3\) 中，对处于一般位置的四条射影直线，恰有两条射影直线同时与它们相交；两条解不同，每条解的局部相交重数为一。
\end{theorem}
\begin{proof}
	每条给定直线产生一份 \(\Omega_1\) 条件。由\cref{g2:prop:schubert-moving-flags}，一般位置下四份条件横截相交于各自开胞腔，故都是孤立的光滑交点，局部符号为 \(+1\)。总数于是为
	\[
		\langle\sigma_1^4,[G]\rangle
		=2\langle\sigma_{22},[G]\rangle=2.
	\]
	横截性的局部方程具有可逆复 Jacobian 矩阵，故每个零点都是单零点，其局部重数为一。
\end{proof}

这个数字还可以在一个直观的二次曲面模型中看出。先取三条两两不交的直线。前两条来自分解 \(\mathbb C^4=A\oplus B\) 的两个二维因子；第三条由于与它们都不交，可以写成某个同构 \(T:A\to B\) 的图。借助 \(T\) 识别 \(B\cong A\)，三条线便为
\[
	L_1=\mathbb P(A\oplus0),\quad L_2=\mathbb P(0\oplus A),\quad
	L_3=\{[a:a]:[a]\in\mathbb P(A)\}.
\]
一条同时遇到前两条线的直线必连接 \([a:0]\) 与 \([0:b]\)。它再与第三条相交，当且仅当 \(a,b\) 成比例。因此三条线的公共截线恰为
\[
	\ell_{[a]}=\{[sa:ta]:[s:t]\in\mathbb{CP}^1\},\qquad[a]\in\mathbb P(A).
\]
这些直线扫出二次曲面
\begin{equation}\label{g2:eq:three-lines-ruled-quadric}
	Q=\{[x_1:x_2:y_1:y_2]:x_1y_2-x_2y_1=0\}
	\cong\mathbb{CP}^1\times\mathbb{CP}^1.
\end{equation}
这个同构就是 \(([a],[s:t])\mapsto[sa:ta]\)：秩一的二乘二矩阵可分解为列向量乘行向量，两个因子分别只差非零倍数，所以参数唯一。非零秩一矩阵处方程的梯度不为零，故曲面光滑。固定 \([s:t]\) 得到包含原三条线的一个线族，固定 \([a]\) 得到公共截线的另一个线族。

\ProseParagraphBreak

一般的第四条线不包含于 \(Q\)，其参数代入二次方程后得到一个有两个不同根的齐次二次式。因而它与 \(Q\) 相交于两个点，每点唯一确定一条 \(\ell_{[a]}\)。这两条截线不同：若两个点属于同一条 \(\ell_{[a]}\)，则第四条线就是这条截线，因而包含于 \(Q\)，与选择矛盾。这给出了两条解的具体几何位置。若第四条线与 \(Q\) 相切，两个交点可能合并；若第四条线本身属于适当线族，还可能出现整族解。一般位置在这里排除的正是这些退化。

\subsection{与 V.27.5 中 Chow 环及 Plücker 方程计算的比较}
\label{g2:subsec:2105:06}
Grassmann 流形也有一个具体的射影方程模型。定义 Pl\"ucker 映射\footnote{Julius Pl\"ucker（1801--1868），德国数学家，研究解析几何、直线几何及数学物理。}
\[
	\iota:G\longrightarrow\mathbb P(\textstyle\bigwedge^k\mathbb C^n),
	\qquad\operatorname{span}(v_1,\ldots,v_k)\longmapsto[v_1\wedge\cdots\wedge v_k].
\]
换基只将外积乘以行列式，所以射影点不变。非零可分解外积 \(\xi=v_1\wedge\cdots\wedge v_k\) 又能恢复原子空间为 \(\{v:v\wedge\xi=0\}\)，故映射单射。在一个非零 \(k\) 阶子式对应的图域中，图矩阵的各项都是其他子式与这一子式之比，所以逆坐标光滑且微分单射。\(G\) 紧，因而 \(\iota\) 是到其像的光滑复嵌入。

\ProseParagraphBreak

其拉回的射影重言线在 \(V\) 处正是 \(\bigwedge^kV\)，所以
\[
	\iota^*\mathcal O(1)=\det S^*,\qquad\iota^*c_1(\mathcal O(1))=\sigma_1.
\]
这里 \(\mathcal O(1)\) 表示复射影空间的超平面线丛。若固定一个 \(q\) 维子空间 \(A\)，投影 \(\mathbb C^n\to\mathbb C^n/A\) 的行列式在 \(\bigwedge^k\mathbb C^n\) 上给出线性函数。它在 \(V\) 的 Pl\"ucker 点处为零，恰好当 \(V\cap A\ne0\)，所以特殊除子 \(\Omega_1\) 确为这一嵌入下的超平面截面。杯积中的每个 \(\sigma_1\)，因此也可以理解为加入一条这样的线性方程。

\ProseParagraphBreak

对 \(\operatorname{Gr}(2,4)\)，令 \(p_{ij}\) 为外积的六个坐标。由 \(\xi\wedge\xi=0\) 得到
\begin{equation}\label{g2:eq:gr24-plucker-quadric}
	p_{12}p_{34}-p_{13}p_{24}+p_{14}p_{23}=0.
\end{equation}
反过来，若某个坐标如 \(p_{12}\ne0\)，把它归一化为一，便可取矩阵
\[
	\begin{pmatrix}
		1&0&-p_{23}&-p_{24}\\
		0&1&p_{13}&p_{14}
	\end{pmatrix}.
\]
它的前五个相应子式是指定坐标，最后一个为 \(p_{13}p_{24}-p_{14}p_{23}=p_{34}\)。其他非零坐标图只需重新排列基，所以式\eqref{g2:eq:gr24-plucker-quadric}正好刻画 Pl\"ucker 像。因而这个 Grassmann 流形是 \(\mathbb{CP}^5\) 中的光滑二次超曲面，\(\int_G\sigma_1^4=2\) 与其二次方程相呼应。

第五卷第27章第5节将沿这个同一模型组织 Chow 环与方程计算：子簇先按有理等价形成代数循环类，再通过复几何的循环类映射与这里的整数上同调比较，并处理具体方程的消去、局部重数和退化参数。当前已经建立的是胞腔基本类、特殊陈类与横截相交计数；一般位置下每个点的局部贡献已由复横截性确定为一。转到代数模型后，重数还要在点合并或交集不横截时继续记录总计数，两种描述因此共享同一算例而承担不同的计算任务。
% 纲要条目排为编号小节标题，并进入目录与书签；说明段仍保留于纲要注释中。


\begin{chaptersummary}
	实线丛的第一 Stiefel--Whitney 类记录方向沿闭路的变化，复线丛的第一陈类是其底层定向实平面丛的 Euler 类。对偶重言复线丛在复射影直线上的第一类取值为 \(1\)，这固定了本章的符号。过渡函数的局部辐角把第一陈类表示为整数二上闭链，实线丛复化则给出整数 Bockstein 类。

	\ProseParagraphBreak

	Leray--Hirsch 定理将整体纤维基转成总空间的自由上同调模。射影丛关系的系数定义各阶陈类与 Stiefel--Whitney 类；迭代射影化所得的旗丛使原丛分裂为线丛，且上同调拉回单射。自然性、Whitney 公式及根的对称计算因而成立。顶陈类等于复定向 Euler 类，顶 Stiefel--Whitney 类等于模二 Euler 类；Thom 平方得到同一组实示性类，对角嵌入又给出 \(w(\mathrm TM)=\operatorname{Sq}(v(M))\)。

	\ProseParagraphBreak

	复化定义整数 Pontryagin 类。它们稳定、自然，但一般直和的整数公式需要保留奇次陈类乘积造成的二阶扭修正。射影空间的切丛稳定同构使三类计算具体化；球面的稳定平凡性则说明这些稳定类与依赖原秩的 Euler 类各有作用。示性类还能给出附加结构的条件：\(w_1\) 检测方向，\(c_1\) 检测行列式平凡性，而已定向丛的 spin 提升恰由 \(w_2\) 控制。

	\ProseParagraphBreak

	复 Grassmann 流形的 Schubert 胞腔给出整数加性基，几何入射条件的对偶类提供乘法解释。特殊 Schubert 类等于商丛的陈类，因而线性子空间之间的相交问题可以转成示性类运算。在 \(\operatorname{Gr}_2(\mathbb C^4)\) 中，四次超平面型 Schubert 条件的乘积求值为 \(2\)，对应三维复射影空间中同时与四条一般位置直线相交的两条直线。
\end{chaptersummary}

\BookExercises
\begin{exercise}\label{g2:ex:line-sections-positive-negative-indices}
	取整数 \(p,q\ge0\)，\(p+q=d>0\)。在 \(\mathbb CP^1\) 上，齐次单项式 \(z_0^pz_1^q\) 给出 \(\mathcal O(d)\) 的一个截面。找出全部零点，计算局部指标，并判断这些零点何时横截。再构造 \(\mathcal O(-d)\) 的光滑截面，使零点仍位于同样的位置，且非零局部指标全部为负。
\end{exercise}
\begin{solution}
	\(\mathcal O(d)=(\gamma^*)^{\otimes d}\)。在每条复直线 \(\ell\subset\mathbb C^2\) 上，齐次多项式限制为一个齐次 \(d\) 次函数，即 \((\ell^*)^{\otimes d}\) 的元素，所以题给式确实定义一个整体截面。

	当 \(p>0\) 时，\([0:1]\) 是零点。在附近以 \(w=z_0/z_1\) 作坐标，以 \((w,1)\) 及其对偶的张量幂作标架，截面系数为 \(w^p\)。在小圆周上归一化后，它是 \(e^{\mathrm i\theta}\mapsto e^{\mathrm i p\theta}\)，故局部指标为 \(p\)。同理，当 \(q>0\) 时，\([1:0]\) 是零点，以 \(v=z_1/z_0\) 作坐标，局部系数为 \(v^q\)，指标为 \(q\)。指数为零的那一端不是零点，此外没有零点。因此指标和是 \(p+q=d\)，与 \(c_1(\mathcal O(d))=dh\) 的求值相同。

	局部映射 \(w\mapsto w^p\) 在零点的实微分可逆恰好当 \(p=1\)；若 \(p>1\)，微分为零。另一端同理。所以所有实际零点横截恰好当 \(p,q\) 中每个非零指数都为一。指标和计入重数，而非只数零点位置。

	将此截面作纤维复共轭，得到 \(\overline{\mathcal O(d)}\) 的光滑截面。Hermitian 内积给出复线性同构 \(\overline{\mathcal O(d)}\cong\mathcal O(d)^*=\mathcal O(-d)\)。在共轭标架中的局部系数分别为 \(\overline w^p\)、\(\overline v^q\)；其归一化圆周映射的度分别为 \(-p\)、\(-q\)。随后施加的丛同构在各纤维中保持复方向，故不改变这些指标。所得截面的指标和为 \(-d\)，零点位置仍与原截面相同。
\end{solution}

\begin{exercise}\label{g2:ex:circle-bundle-torsion-line}
	在 \(\mathbb CP^1\) 上取 \(\mathcal O(m)\)，其中 \(m\ge2\)，令 \(Y_m=S(\mathcal O(m))\) 为其单位圆丛，投影为 \(\pi\)。对 \(Y_m\) 上的线丛 \(L=\pi^*\mathcal O(1)\)，计算 \(c_1(L)\) 的阶，并确定哪些整数 \(k\) 使 \(L^{\otimes k}\) 平凡。说明 \(c_1(L)\) 换成实系数后为何为零，并直接写出 \(L^{\otimes m}\) 的非零截面。
\end{exercise}
\begin{solution}
	\(Y_m\) 是紧光滑流形，第17章的 Morse 胞腔模型使它具有有限 CW 同伦型；结合丛的同伦不变性，可以在它上面使用线丛的上同调分类。记 \(h=c_1(\mathcal O(1))\)。给底层实二维丛取复方向，其 Euler 类为 \(mh\)。\cref{g2:thm:thom-gysin}的相应一段为
	\[
		H^0(\mathbb CP^1;\mathbb Z)=\mathbb Z
		\xrightarrow{\ \times m\ }H^2(\mathbb CP^1;\mathbb Z)=\mathbb Z
		\xrightarrow{\pi^*}H^2(Y_m;\mathbb Z)\longrightarrow0.
	\]
	于是 \(H^2(Y_m;\mathbb Z)\cong\mathbb Z/m\)，\(c_1(L)=\pi^*h\) 是其生成元，恰有阶 \(m\)。线丛的张量分类给出
	\[
		L^{\otimes k}\text{ 平凡}\quad\Longleftrightarrow\quad k\,c_1(L)=0
		\quad\Longleftrightarrow\quad m\mid k.
	\]
	负整数 \(k\) 按对偶解释，结论不变。换成实系数后，乘 \(m\) 是 \(\mathbb R\) 上的同构，同一段序列给出 \(H^2(Y_m;\mathbb R)=0\)，所以该非零整数扭类的实像为零。

	还可以直接看出第 \(m\) 个张量幂平凡。由定义，\(L^{\otimes m}=\pi^*\mathcal O(m)\)。一个点 \(y\in Y_m\) 本身就是 \(\mathcal O(m)_{\pi(y)}\) 中的单位向量，因此
	\[
		s(y)=(y,y)\in\bigl(\pi^*\mathcal O(m)\bigr)_y
	\]
	是处处非零的连续截面。它给出完整线标架，从而直接产生这个张量幂的平凡化。
\end{solution}

\begin{exercise}\label{g2:ex:klein-line-complexification}
	给 Klein 瓶 \(K\) 取一个零胞腔、两个有向一胞腔 \(a,b\) 及一个按字 \(aba^{-1}b\) 附着的二胞腔。令 \(\alpha,\beta\in H^1(K;\mathbb F_2)\) 分别在 \(a,b\) 上取值为 \((1,0)\)、\((0,1)\)。对 \(\varepsilon,\eta\in\mathbb F_2\)，记实线丛 \(L_{\varepsilon,\eta}\) 满足 \(w_1(L_{\varepsilon,\eta})=\varepsilon\alpha+\eta\beta\)。计算这些实线丛复化后的第一陈类，判断复化何时平凡，并说明四个不同实线丛会得到多少种复线丛。
\end{exercise}
\begin{solution}
	附着字中 \(a\) 的代数次数为零，\(b\) 的代数次数为二，所以整数胞腔边界为
	\[
		\partial_2(1)=2b,\qquad\partial_1=0.
	\]
	因此整数胞腔上链的非零余边界是
	\[
		\mathrm d^1:\mathbb Z^2\longrightarrow\mathbb Z,\qquad(u,v)\longmapsto2v,
	\]
	从而 \(H^2(K;\mathbb Z)=\mathbb Z/2\)。记二胞腔对偶上链所代表的生成元为 \(z\)。模二余边界全为零，故 \(\alpha,\beta\) 确实构成一次群的一组基。

	按整数 Bockstein 的除二公式，将 \(\alpha\) 提升为整数上链 \((1,0)\)，其余边界为零，所以 \(\widetilde\beta_2(\alpha)=0\)。将 \(\beta\) 提升为 \((0,1)\)，余边界为 \(2\)，除以二后得到二次上链 \(1\)，所以 \(\widetilde\beta_2(\beta)=z\)。由\cref{g2:prop:real-line-complexification-bockstein}，
	\[
		c_1(L_{\varepsilon,\eta}\otimes_{\mathbb R}\mathbb C)=\eta z.
	\]
	因此恰在 \(\eta=0\) 时复化平凡。原来的四个实线丛由四个不同的 \(w_1\) 区分；复化后只有两种同构类，\(\varepsilon\) 的差别消失，\(\eta\) 的差别保留为整数二阶扭类。这个计算在同一底空间上同时展示了实线丛非平凡性被复化消去和被复化保留的两种情况。
\end{solution}

\begin{exercise}\label{g2:ex:leray-hirsch-orientation-torus}
	令 \(a:S^2\to S^2\) 为对径映射，并取映射环面
	\[
		Y=(S^2\times[0,1])/\bigl((x,1)\sim(a(x),0)\bigr),\qquad p:Y\to S^1.
	\]
	判断纤维的整数二次生成元能否延拓为 \(Y\) 上的整数类，计算 \(H^*(Y;\mathbb Z)\)。再证明模二系数下存在满足 Leray--Hirsch 条件的整体纤维基，并计算模二上同调环。
\end{exercise}
\begin{solution}
	对径映射在 \(H^2(S^2;\mathbb Z)\) 上作用为 \(-1\)，因为\cref{g2:prop:reflection-antipodal-degree}给出它的度为 \((-1)^3=-1\)。若整数生成元来自整体类，沿基圆一周运输后必须保持不变，这与生成元变号矛盾。

	用基圆的零胞腔及一胞腔计算局部系数，\(q=0\) 行是常值 \(\mathbb Z\)，\(q=2\) 行的上链微分是 \(-2:\mathbb Z\to\mathbb Z\)。因此 Serre 谱序列的非零第二页项只有
	\[
		E_2^{0,0}=\mathbb Z,\qquad E_2^{1,0}=\mathbb Z,\qquad E_2^{1,2}=\mathbb Z/2.
	\]
	底空间只有两个可能的列，全部高阶微分为零；每个总次数也至多有一个非零项，故没有群延拓选择。于是
	\[
		H^k(Y;\mathbb Z)=
		\begin{cases}
			\mathbb Z,&k=0,1,\\
			\mathbb Z/2,&k=3,\\
			0,&\text{其他次数}.
		\end{cases}
	\]
	特别地，\(H^2(Y;\mathbb Z)=0\) 也直接排除了整数纤维生成元的延拓。

	模二系数下，\(-1=1\)，两行都是常值 \(\mathbb F_2\)。同一谱序列仍无非零微分，且边缘映射
	\[
		H^2(Y;\mathbb F_2)\longrightarrow H^2(S^2;\mathbb F_2)
	\]
	为同构。选 \(u\) 映到纤维生成元，则 \(1,u\) 满足 Leray--Hirsch 条件。令 \(x=p^*\eta\)，其中 \(\eta\) 是基圆的一次非零类。于是 \(1,x,u,xu\) 为加性基。\(x^2=p^*(\eta^2)=0\)，而 \(u^2=0\)，因为四次群为零；\(xu\ne0\) 由 Leray--Hirsch 的唯一表达保证。因此
	\[
		H^*(Y;\mathbb F_2)\cong\mathbb F_2[x,u]/(x^2,u^2),\qquad |x|=1,\quad |u|=2.
	\]
	同一纤维丛在不同系数下是否具有整体纤维基，取决于运输作用是否保留所需生成元。
\end{solution}

\begin{exercise}\label{g2:ex:projectivization-line-twist}
	设 \(E\to B\) 为复秩 \(r\ge1\) 丛，\(L\to B\) 为复线丛。构造覆盖底空间恒等映射的自然同构 \(\mathbb P(E)\cong\mathbb P(E\otimes L)\)，并比较两张重言线丛及其对偶的一阶类。若 \(r=2\)，将两条射影丛关系的系数分别记为 \(a_1,a_2\) 和 \(a'_1,a'_2\)。用 \(\ell=c_1(L)\) 表示 \(a'_1,a'_2\)。
\end{exercise}
\begin{solution}
	在每个 \(b\in B\) 上，将一条线 \(V\subset E_b\) 送到 \(V\otimes L_b\subset E_b\otimes L_b\)。局部选取 \(L\) 的非零标架后，这是通常的线子空间识别；更换标架只乘一个非零复数，不改变线，所以各局部映射粘合。张量 \(L_b^*\) 给出连续逆，因此得到自然同构 \(f:\mathbb P(E)\to\mathbb P(E\otimes L)\)。按每个点的纤维定义，有
	\[
		f^*\lambda_{E\otimes L}\cong\lambda_E\otimes p^*L.
	\]
	令 \(h=c_1(\lambda_E^*)\)、\(h'=f^*c_1(\lambda_{E\otimes L}^*)\)。线丛的张量及对偶公式给出
	\[
		h'=h-p^*\ell.
	\]
	所以射影丛作为空间相同，其用于记录原丛的重言线类仍可能不同。

	当 \(r=2\) 时，原关系为 \(h^2+p^*a_1h+p^*a_2=0\)。代入 \(h=h'+p^*\ell\)，整理得
	\[
		(h')^2+p^*(a_1+2\ell)h'+p^*(a_2+a_1\ell+\ell^2)=0.
	\]
	由射影关系系数的唯一性，
	\[
		a'_1=a_1+2\ell,\qquad a'_2=a_2+a_1\ell+\ell^2.
	\]
	这也说明原向量丛的信息不仅在射影总空间中，还在其重言线丛及投影中。
\end{solution}

\begin{exercise}\label{g2:ex:exterior-square-rank-three-roots}
	设 \(E\to B\) 为复秩三向量丛，记本章构造的陈类为 \(c_1,c_2,c_3\)。用分裂原理计算 \(\bigwedge^2E\) 的三个陈类，说明计算中的每个整数系数如何出现。
\end{exercise}
\begin{solution}
	在使 \(E\) 分裂的旗空间上写 \(q^*E=L_1\oplus L_2\oplus L_3\)，令 \(x_i=c_1(L_i)\)。外积的双线性给出实际丛分解
	\[
		\bigwedge^2(q^*E)\cong(L_1\otimes L_2)\oplus(L_1\otimes L_3)\oplus(L_2\otimes L_3).
	\]
	线丛张量的一阶类可加，所以新丛的三个形式根是
	\[
		y_1=x_1+x_2,\qquad y_2=x_1+x_3,\qquad y_3=x_2+x_3.
	\]
	它们的一次和为 \(2(x_1+x_2+x_3)\)，因为每个 \(x_i\) 出现两次。两两乘积之和为
	\[
	\begin{aligned}
		y_1y_2+y_1y_3+y_2y_3
		&=x_1^2+x_2^2+x_3^2+3(x_1x_2+x_1x_3+x_2x_3)\\
		&=(x_1+x_2+x_3)^2+(x_1x_2+x_1x_3+x_2x_3).
	\end{aligned}
	\]
	每个平方出现一次，每个混合项出现三次，所以最后恰剩一个二阶初等对称项。三个根的乘积为
	\[
		(x_1+x_2)(x_1+x_3)(x_2+x_3)
		=(x_1+x_2+x_3)(x_1x_2+x_1x_3+x_2x_3)-x_1x_2x_3.
	\]
	右边第一项中三重乘积出现三次，减去一次后留下左边的两次；其余六个三次单项式各出现一次。把初等对称多项式识别为 \(q^*c_i\)，再由 \(q^*\) 单射返回底空间，得到
	\[
		c_1\!\left(\bigwedge^2E\right)=2c_1,\qquad
		c_2\!\left(\bigwedge^2E\right)=c_1^2+c_2,\qquad
		c_3\!\left(\bigwedge^2E\right)=c_1c_2-c_3.
	\]
\end{solution}

\begin{exercise}\label{g2:ex:endomorphism-chern-low-classes}
	设 \(E\to B\) 为 CW 复形上的复秩 \(r\ge1\) 向量丛。计算复向量丛 \(\operatorname{End}(E)=E^*\otimes E\) 的第一、第二陈类，用 \(c_1(E),c_2(E)\) 表示结果。再说明将 \(\operatorname{End}(E)\) 换成无迹自同态丛 \(\operatorname{End}_0(E)\) 时，所有正次陈类为何不变。
\end{exercise}
\begin{solution}
	将 \(E\) 拉回到分裂空间，记其形式根为 \(x_1,\ldots,x_r\)。由\cref{g2:prop:chern-operations}，\(\operatorname{End}(E)\) 的根为全部 \(x_j-x_i\)。当 \(i=j\) 时根为零；其余根可按 \((i,j)\) 与 \((j,i)\) 配对，因此
	\[
		c(\operatorname{End}(E))
		=\prod_{i<j}\bigl(1+(x_j-x_i)\bigr)\bigl(1-(x_j-x_i)\bigr)
		=\prod_{i<j}\bigl(1-(x_j-x_i)^2\bigr)
	\]
	成立于分裂空间。每个因子的第一个正次项都是四次，所以 \(c_1(\operatorname{End}(E))=0\)。四次项给出
	\[
	\begin{aligned}
		c_2(\operatorname{End}(E))
		&=-\sum_{i<j}(x_j-x_i)^2\\
		&=-(r-1)\sum_i x_i^2+2\sum_{i<j}x_ix_j\\
		&=-(r-1)\bigl(c_1(E)^2-2c_2(E)\bigr)+2c_2(E)\\
		&=2r\,c_2(E)-(r-1)c_1(E)^2.
	\end{aligned}
	\]
	第二行中，每个 \(x_i^2\) 在含有下标 \(i\) 的 \(r-1\) 对中出现；混合项来自平方展开的 \(-2x_ix_j\)。旗丛拉回单射使这些等式返回原基空间。配对计算始终使用整数多项式恒等式，没有对上同调类除以二，故包含整数扭信息的基空间也适用。

	迹映射在换架共轭下不变，故 \(\operatorname{End}_0(E)=\ker\operatorname{tr}\) 是复秩 \(r^2-1\) 子丛。映射
	\[
		A\longmapsto\left(\frac{\operatorname{tr}A}{r},\ A-\frac{\operatorname{tr}A}{r}I\right)
	\]
	给出 \(\operatorname{End}(E)\cong\underline{\mathbb C}\oplus\operatorname{End}_0(E)\)，逆映射为 \((z,A_0)\mapsto zI+A_0\)。这里的除以 \(r\) 发生在复纤维的线性代数中。平凡线因子的总陈类为一，所以 Whitney 公式说明两张丛的全部正次陈类相同。
\end{solution}

\begin{exercise}\label{g2:ex:three-real-lines-projective-product}
	设 \(m,n\ge1\)，取 \(B=\mathbb RP^m\times\mathbb RP^n\)。记从两因子拉回的典范实线丛为 \(L_a,L_b\)，令 \(a=w_1(L_a)\)、\(b=w_1(L_b)\)，并置
	\[
		E=L_a\oplus L_b\oplus(L_a\otimes L_b).
	\]
	计算 \(w_1(E),w_2(E),w_3(E)\)，判断 \(E\) 是否可定向、是否平凡，并确定哪些 \((m,n)\) 使它存在处处非零截面。
\end{exercise}
\begin{solution}
	射影空间模二环及域系数 Künneth 公式给出
	\[
		H^*(B;\mathbb F_2)=\mathbb F_2[a,b]/(a^{m+1},b^{n+1}).
	\]
	第三条线的第一类为 \(a+b\)，所以 Whitney 公式得到
	\[
	\begin{aligned}
		w(E)&=(1+a)(1+b)(1+a+b),\\
		w_1(E)&=0,\\
		w_2(E)&=a^2+ab+b^2,\\
		w_3(E)&=a^2b+ab^2.
	\end{aligned}
	\]
	由\cref{g2:prop:first-classes-determinants}，\(w_1=0\) 说明 \(E\) 可定向。由于 \(m,n\ge1\)，单项式 \(ab\) 在上述截断多项式环中始终非零，并且与尚未消失的 \(a^2,b^2\) 线性无关。因此 \(w_2(E)\ne0\)，\(E\) 对所有给定 \((m,n)\) 都不平凡。

	当至少一个 \(m,n\) 大于一时，\(a^2b\)、\(ab^2\) 中至少有一个非零；若两项均非零，它们是不同的基单项式，不能相消。所以此时 \(w_3(E)\ne0\)。\cref{g2:thm:top-class-euler}把它识别为模二 Euler 类，而非零截面会使该类为零，故这些情形没有处处非零截面。

	剩下 \(m=n=1\) 时，底空间是二维环面，\(w_3(E)=0\)。给 \(E\) 选内积，单位球丛的纤维为 \(S^2\)。先在零骨架选单位向量；纤维道路连通，使该选择能沿每条一胞腔延拓。每张二维胞腔上的丛平凡，边界截面对应一条 \(S^1\to S^2\) 的映射，而 \(\pi_1(S^2)=0\)，故还能延入圆盘。底空间没有更高胞腔，得到整体单位截面。因此恰在 \((m,n)=(1,1)\) 时存在非零截面。此时 \(w_2(E)=ab\ne0\)，所以一条非零截面仍不足以形成整个三维丛的标架。
\end{solution}

\begin{exercise}\label{g2:ex:real-projective-wu-formula-all-dimensions}
	设 \(n\ge1\)，令 \(a\in H^1(\mathbb RP^n;\mathbb F_2)\) 为非零类。直接由 Wu 类的求值定义与射影空间上的平方公式证明
	\[
		v_i(\mathbb RP^n)=\binom{n-i}{i}a^i\qquad(0\le i\le n),
	\]
	其中二项式系数取模二，且下标大于非负上标时约定为零。说明 \(i>n/2\) 时为何消失，并分别计算 \(n=4,5\) 时的总 Wu 类及由 Wu 公式得到的切丛总 Stiefel--Whitney 类。
\end{exercise}
\begin{solution}
	对 \(0\le i\le n\)，\(H^i(\mathbb RP^n;\mathbb F_2)\) 由 \(a^i\) 生成，因而可写 \(v_i=\lambda_i a^i\)。模二基本类满足 \(\langle a^n,[\mathbb RP^n]_2\rangle=1\)，所以把 Wu 类的定义用于 \(a^{n-i}\)，得到
	\[
	\begin{aligned}
		\lambda_i
		&=\langle v_i a^{n-i},[\mathbb RP^n]_2\rangle\\
		&=\langle\operatorname{Sq}^i(a^{n-i}),[\mathbb RP^n]_2\rangle
		=\binom{n-i}{i}.
	\end{aligned}
	\]
	最后一步使用式\eqref{g2:eq:projective-steenrod-formula}。若 \(i>n/2\)，则 \(i>n-i\)，平方次数超过输入类的次数，不稳定条件使它为零；这也与二项式系数的约定一致。\(i>n\) 时所在上同调群本来为零。

	当 \(n=4\) 时，可能的正次项只有一次和二次，系数分别为 \(\binom31=1\)、\(\binom22=1\)（模二），故
	\[
		v(\mathbb RP^4)=1+a+a^2.
	\]
	由平方公式，\(\operatorname{Sq}(a)=a+a^2\)，\(\operatorname{Sq}(a^2)=a^2+a^4\)，中间的 \(\operatorname{Sq}^1(a^2)=2a^3=0\)。因此\cref{g2:thm:wu-formula}给出
	\[
		w(\mathrm T\mathbb RP^4)=\operatorname{Sq}(1+a+a^2)=1+a+a^4.
	\]
	当 \(n=5\) 时，\(\binom41=0\)、\(\binom32=1\)（模二），所以
	\[
		v(\mathbb RP^5)=1+a^2,\qquad
		w(\mathrm T\mathbb RP^5)=1+a^2+a^4.
	\]
	两种情形中 Wu 类的高次部分都先消失，随后对低次 Wu 类施加平方，仍能产生切丛较高次的非零类。两组结果也分别等于截断后的 \((1+a)^5\)、\((1+a)^6\)，与射影空间切丛的直接计算一致。
\end{solution}

\begin{exercise}\label{g2:ex:projective-stiefel-binary}
	对 \(n\ge1\)，证明 \(\mathbb RP^n\) 的全部正次 Stiefel--Whitney 类为零，当且仅当 \(n+1\) 是二的幂。再对 \(\mathbb RP^3\)、\(\mathbb RP^5\)、\(\mathbb RP^7\) 分别判断可定向性与 spin 结构的存在性，并计算存在时 spin 结构的个数。
\end{exercise}
\begin{solution}
	令 \(m=n+1\)，在 \(\mathbb F_2[a]/(a^m)\) 中有 \(w(\mathrm T\mathbb RP^n)=(1+a)^m\)。将 \(m\) 写成二进制 \(m=\sum_{j\in J}2^j\)。因为特征二下 \((1+a)^{2^j}=1+a^{2^j}\)，所以
	\[
		(1+a)^m=\prod_{j\in J}(1+a^{2^j}).
	\]
	若 \(m=2^k\)，右端为 \(1+a^m=1\)。反过来，若 \(J\) 至少有两个元素，令 \(j_0=\min J\)。右端 \(a^{2^{j_0}}\) 的系数为一：只能从第 \(j_0\) 个因子选择正次项，其他正次项的次数都更大。而 \(2^{j_0}<m\)，这一项在商环中不为零。因此总类不等于一，得到所求判据。

	三个维数都是奇数，所以相应实射影空间都可定向。\(w_2\) 的系数依次是 \(\binom42=6\)、\(\binom62=15\)、\(\binom82=28\)，故 \(\mathbb RP^3\)、\(\mathbb RP^7\) 有 spin 结构，\(\mathbb RP^5\) 没有。前两者的一次模二上同调均为 \(\mathbb F_2\)，所以各有两份 spin 结构。正次 Stiefel--Whitney 类全零是示性类层面的结论；判定实际切丛是否平凡，还需研究完整的标架延拓。
\end{solution}

\begin{exercise}\label{g2:ex:product-projective-chern-numbers}
	给 \(M=\mathbb CP^1\times\mathbb CP^2\) 乘积复结构。记两个因子的正生成元拉回为 \(x,y\)，所以 \(x^2=0\)、\(y^3=0\)，且 \(\langle xy^2,[M]\rangle=1\)。求 \(c_1,c_2,c_3\)、\(p_1(\mathrm TM)\)，计算三个陈数 \(\langle c_1^3,[M]\rangle\)、\(\langle c_1c_2,[M]\rangle\)、\(\langle c_3,[M]\rangle\)，并判断 \(M\) 是否有 spin 结构。
\end{exercise}
\begin{solution}
	乘积切丛分裂为两个因子的切丛拉回，所以
	\[
		c(\mathrm TM)=(1+x)^2(1+y)^3
		=(1+2x)(1+3y+3y^2).
	\]
	按次数收集得到
	\[
		c_1=2x+3y,\qquad c_2=6xy+3y^2,\qquad c_3=6xy^2.
	\]
	因此
	\[
		p_1=c_1^2-2c_2=(12xy+9y^2)-(12xy+6y^2)=3y^2.
	\]
	只有 \(xy^2\) 项在顶次求值中有贡献，于是
	\[
		c_1^3=54xy^2,\qquad c_1c_2=24xy^2,\qquad c_3=6xy^2.
	\]
	三个数分别为 \(54,24,6\)。最后 \(w_2(\mathrm TM)=\rho_2c_1=\overline y\ne0\)，其非零性可通过限制到第二个因子的射影直线检验。因此 \(M\) 没有 spin 结构。
\end{solution}

\begin{exercise}\label{g2:ex:two-spin-lifts-circle}
	设 \(E=S^1\times\mathbb R^3\) 带标准方向与度量。用 \(\operatorname{Spin}(3)=S^3\) 及四元数覆盖 \(\rho:S^3\to\operatorname{SO}(3)\)，对 \(\epsilon\in\{1,-1\}\) 定义
	\[
		Q_\epsilon=(\mathbb R\times S^3)/\!\sim,
		\qquad(t+1,q)\sim(t,\epsilon q).
	\]
	证明 \(\eta_\epsilon([t,q])=(e^{2\pi\mathrm i t},\rho(q))\) 给出 \(E\) 的 spin 结构。证明二者作为覆盖固定标架丛的 spin 结构不同构，却作为抽象主 \(S^3\) 丛都平凡，并说明二者为何穷尽所有 spin 结构。
\end{exercise}
\begin{solution}
	\(\epsilon\) 位于 \(S^3\) 的中心，因此右乘作用 \([t,q]g=[t,qg]\) 有定义，短弧上的实参数提升给出主丛局部平凡化。又 \(\rho(\epsilon q)=\rho(q)\)，故 \(\eta_\epsilon\) 尊重商关系，且逐纤维为 \(\rho\)、满足所需等变性。

	取固定标架丛中的截面 \(s(z)=(z,I_3)\)。其在 \(Q_\epsilon\) 中的原像是
	\[
		(\mathbb R\times\{1,-1\})/\bigl((t+1,a)\sim(t,\epsilon a)\bigr).
	\]
	\(\epsilon=1\) 时，这是两个互不相交的圆；\(\epsilon=-1\) 时，两张局部覆盖沿一周互换，得到一个连通的二重圆周覆盖。任何与 \(\eta\) 交换的同构都须把这两份原像同构地对应，连通分支数不同使它不可能存在。

	作为抽象主丛，\(Q_1\) 已经是乘积。对 \(Q_{-1}\)，取单位纯虚四元数 \(u\)，令 \(h(t)=\cos(\pi t)+u\sin(\pi t)\)，则 \(h(t+1)=-h(t)\)。映射
	\[
		[t,q]\longmapsto(e^{2\pi\mathrm i t},h(t)q)
	\]
	尊重接缝、保持右作用，并在局部具有连续逆，所以给出主丛平凡化。它改变了到固定 \(\operatorname{SO}(3)\) 标架丛的映射，因而不与前面不同构的结论冲突。由\cref{g2:thm:spin-obstruction}，所有 spin 结构同构类受 \(H^1(S^1;\mathbb F_2)=\mathbb F_2\) 自由传递作用，恰有两份，故上述构造已经穷尽。
\end{solution}

\begin{exercise}\label{g2:ex:smooth-projective-surface-numbers}
	设 \(F(z_0,z_1,z_2,z_3)\) 为次数 \(d\ge1\) 的复齐次多项式，其射影零点集 \(X\subset\mathbb CP^3\) 是光滑复曲面；等价地，所定义的齐次方程沿零点集的法向微分处处满射。记 \(h=c_1(\lambda^*)\)，并用同一字母表示其在 \(X\) 上的限制。证明
	\[
		c_1(\mathrm TX)=(4-d)h,\qquad
		c_2(\mathrm TX)=(d^2-4d+6)h^2,\qquad
		\langle h^2,[X]\rangle=d.
	\]
	据此计算 \(X\) 的两个陈数及第一 Pontryagin 数，并证明 \(X\) 有 spin 结构当且仅当 \(d\) 为偶数。
\end{exercise}
\begin{solution}
	在一条直线 \(\ell\subset\mathbb C^4\) 上，\(F\) 的限制是 \(\ell\) 上的次数 \(d\) 齐次函数，所以是 \((\ell^*)^{\otimes d}\) 的元素。因此 \(F\) 给出线丛 \((\lambda^*)^{\otimes d}\) 的光滑截面，零集正是 \(X\)。题设满秩条件说明该截面横截于零，并将复法丛识别为 \((\lambda^*)^{\otimes d}|_X\)。选 Hermitian 度量分裂切丛的短正合列，得到复丛同构
	\[
		\mathrm T\mathbb CP^3|_X\cong\mathrm TX\oplus(\lambda^*)^{\otimes d}|_X.
	\]
	线丛张量公式给出后一项的第一陈类为 \(dh\)，因此
	\[
		c(\mathrm TX)=\frac{(1+h)^4}{1+dh}
		=(1+4h+6h^2)(1-dh+d^2h^2)
		=1+(4-d)h+(d^2-4d+6)h^2,
	\]
	这里只保留实次数不超过四的项。

	由\cref{g2:thm:euler-transverse-zero-duality}，\(X\) 的对偶类是这条复线丛的 Euler 类，即 \(dh\)。复法向实维数为二，交换法向与切向的方向不会产生符号，所以该零点方向与 \(X\) 的复方向一致。于是
	\[
		\langle h^2,[X]\rangle=\langle h^2\smile dh,[\mathbb CP^3]\rangle=d.
	\]
	从而
	\[
	\begin{aligned}
		\langle c_1^2,[X]\rangle&=d(4-d)^2,\\
		\langle c_2,[X]\rangle&=d(d^2-4d+6),\\
		\langle p_1(\mathrm TX),[X]\rangle&=\langle c_1^2-2c_2,[X]\rangle=d(4-d^2).
	\end{aligned}
	\]
	最后 \(w_2(\mathrm TX)=(4-d)\overline h\)。\(d\) 偶数时它为零。\(d\) 奇数时，\(\langle\overline h^{\,2},[X]_2\rangle=1\)，故 \(\overline h\ne0\)，于是 \(w_2\ne0\)。由 spin 提升准则得到所求充要条件。比如 \(d=4\) 时两种陈数为 \(0,24\)，第一 Pontryagin 数为 \(-48\)；这些数均由上述切丛与零点类计算得到。
\end{solution}

\begin{exercise}\label{g2:ex:gr25-cells-pairing}
	列出 \(\operatorname{Gr}(2,5)\) 的各次整数 Betti 数与 Euler 示性数。给出四次、八次上同调的 Schubert 基，并写出两组基之间的相交配对矩阵。
\end{exercise}
\begin{solution}
	此时 \(k=2\)、\(q=3\)、\(N=6\)。矩形中的分拆按大小排列为
	\[
	\begin{array}{c|l}
		|\lambda|&\lambda\\\hline
		0&(0,0)\\
		1&(1,0)\\
		2&(2,0),(1,1)\\
		3&(3,0),(2,1)\\
		4&(3,1),(2,2)\\
		5&(3,2)\\
		6&(3,3).
	\end{array}
	\]
	故从零次到十二次的偶数 Betti 数依次为 \(1,1,2,2,2,1,1\)，全部奇数 Betti 数为零。Euler 示性数为这些数之和，即 \(10\)。这里每个分拆贡献一个偶维胞腔，所以也可直接数得 \(\binom52=10\)。

	四次上同调取基 \((\sigma_2,\sigma_{11})\)，八次取基 \((\sigma_{31},\sigma_{22})\)。补分拆公式为 \(\lambda^c=(3-\lambda_2,3-\lambda_1)\)，所以 \((2,0)^c=(3,1)\)、\((1,1)^c=(2,2)\)。按上述次序，相交配对矩阵为
	\[
		\begin{pmatrix}1&0\\0&1\end{pmatrix}.
	\]
	这里第二组基与第一组基所在次数不同，因为 Grassmann 流形实维数为十二，互补次数是四与八。
\end{solution}

\begin{exercise}\label{g2:ex:gr24-tangent-chern-numbers}
	在 \(G=\operatorname{Gr}(2,4)\) 上令 \(h=\sigma_1\)、\(s=\sigma_{11}\)。证明复切丛同构于 \(S^*\otimes Q\)，计算 \(c_1(\mathrm TG)\)、\(c_2(\mathrm TG)\)、\(c_4(\mathrm TG)\)，并将顶陈数与胞腔计算得到的 Euler 示性数比较。
\end{exercise}
\begin{solution}
	在 \(V\) 附近，另一个二维子空间可写成从 \(V\) 到补空间的线性映射的图。对这些图在零映射处求微分，得到 \(\mathrm T_VG=\operatorname{Hom}_{\mathbb C}(V,\mathbb C^4/V)\)。更换补空间不改变商值的一阶映射，所以这些识别粘成 \(\mathrm TG\cong S^*\otimes Q\)。

	在共同分裂空间上记 \(S^*\) 的根为 \(x_1,x_2\)，\(Q\) 的根为 \(y_1,y_2\)。由正文计算，
	\[
		x_1+x_2=y_1+y_2=h,\qquad x_1x_2=s,\qquad y_1y_2=h^2-s.
	\]
	切丛四个根为 \(x_i+y_j\)。它们之和为 \(4h\)，所以 \(c_1(\mathrm TG)=4h\)。平方和为
	\[
	\begin{aligned}
		\sum_{i,j}(x_i+y_j)^2
		&=2(h^2-2s)+2\bigl(h^2-2(h^2-s)\bigr)+2h^2\\
		&=2h^2.
	\end{aligned}
	\]
	由根的对称多项式恒等式 \(2c_2=c_1^2-\sum_{i,j}(x_i+y_j)^2\)，得到 \(2c_2=14h^2\)。整数上同调没有二扭，故 \(c_2(\mathrm TG)=7h^2\)。

	顶类是四根之积。先对 \(y_j\) 求积，并利用 \(x_i^2=hx_i-s\)，有
	\[
	\begin{aligned}
		c_4(\mathrm TG)
		&=\prod_{i=1}^2\bigl(x_i^2+hx_i+h^2-s\bigr)\\
		&=\prod_{i=1}^2(2hx_i+h^2-2s)\\
		&=3h^4-4h^2s+4s^2=3h^4.
	\end{aligned}
	\]
	最后一步使用 \(s^2=h^2s\)。共同旗丛拉回单射，因此这些根计算给出原空间上的等式。由 \(\int_Gh^4=2\)，顶陈数为 \(6\)。\(G\) 恰有六个偶维胞腔，Euler 示性数也为 \(6\)；两者相等是复切丛顶陈类、实 Euler 类与\cref{g2:cor:tangent-euler-characteristic}的具体核验。
\end{solution}

\begin{exercise}\label{g2:ex:schubert-divisor-singularity}
	令 \(A\subset\mathbb C^4\) 为固定二维子空间，\(\Omega_1=\{V\in\operatorname{Gr}(2,4):V\cap A\ne0\}\)。证明 \(\Omega_1\) 在点 \(V=A\) 处不光滑，并在该点附近给出它的方程。再说明正文中的关联空间 \(Z=\{(V,L):L\subset A\cap V\}\) 在这个点上的投影纤维是什么。
\end{exercise}
\begin{solution}
	选择补空间 \(B\)，把 \(A\) 附近的 \(V\) 写成线性映射 \(T:A\to B\) 的图。则 \(V\cap A\) 对应 \(\ker T\)，所以局部方程为
	\[
		\det T=t_{11}t_{22}-t_{12}t_{21}=0
	\]
	于 \(\operatorname{Hom}(A,B)\cong\mathbb C^4\) 中。在 \(T\ne0\) 且行列式为零时，\(T\) 的秩为一，行列式的至少一个偏导数非零，故该处为复三维光滑超曲面。

	在 \(T=0\) 处，集合是所有秩至多一矩阵组成的锥。若它在那里是复三维光滑子流形，则其切空间必须包含每条曲线 \(t\mapsto tE_{ij}\) 在零点的速度。四个矩阵单位 \(E_{ij}\) 都秩一，且线性张成全部 \(\mathbb C^4\)，这与切空间只有复维数三矛盾。因此原点确实是奇点，而不仅是这个方程的一阶导数恰好同时消失。

	投影到 \(V=A\) 时，条件为任意直线 \(L\subset A\)，所以纤维是 \(\mathbb P(A)\cong\mathbb{CP}^1\)。在其余 \(\Omega_1\) 的点上，\(V\cap A\) 一维，纤维只有一个点。关联空间把奇点处不唯一的交线选择保留下来；正文通过其基本类推前定义同一个 Schubert 类，因此闭包有奇点并不妨碍类的存在。
\end{solution}

\begin{exercise}\label{g2:ex:real-four-lines-count}
	在 \(\mathbb C^4=\mathbb C^2\oplus\mathbb C^2\) 中取射影直线
	\[
		L_1=\mathbb P(\mathbb C^2\oplus0),\quad L_2=\mathbb P(0\oplus\mathbb C^2),\quad L_3=\mathbb P\{(a,a):a\in\mathbb C^2\}.
	\]
	第四条直线取 \(L_T=\mathbb P\{(a,Ta):a\in\mathbb C^2\}\)。分别令
	\[
		T_+=\begin{pmatrix}0&2\\3&0\end{pmatrix},\qquad
		T_-=\begin{pmatrix}0&-1\\2&0\end{pmatrix}.
	\]
	证明每组四条线两两不交，并求所有复公共截线。两组数据都由实矩阵给出，它们各有多少条实公共截线？
\end{exercise}
\begin{solution}
	前面三条直线两两不交。图直线与 \(L_2\) 不交是因为第一分量为零会迫使 \(a=0\)；与 \(L_1\) 不交等价于 \(T\) 可逆，与 \(L_3\) 不交等价于 \(T-I\) 可逆。两个所给矩阵分别满足
	\[
		\det T_+=-6,\quad\det(T_+-I)=-5,\qquad
		\det T_-=2,\quad\det(T_--I)=3,
	\]
	所以四条线两两不交。

	前三条线的公共截线由正文给出，为 \(\ell_{[a]}=\{[sa:ta]:[s:t]\in\mathbb{CP}^1\}\)。它与 \(L_T\) 相交，恰好当 \(Ta\) 与 \(a\) 成比例，即 \([a]\) 是 \(T\) 的特征线。对 \(T_+\)，特征值为 \(\pm\sqrt6\)，可取特征向量 \(a_\pm=(2,\pm\sqrt6)\)。所以两条公共截线为 \(\ell_{[a_+]}\)、\(\ell_{[a_-]}\)，均由实向量张成，是两条不同的实射影直线。

	对 \(T_-\)，特征值为 \(\pm\mathrm i\sqrt2\)，对应特征向量可取 \((1,\mp\mathrm i\sqrt2)\)。仍有两条不同的复公共截线，它们互为复共轭，且都不由实二维子空间定义。因为实矩阵的实特征线必有实特征值，而这里没有实特征值，所以没有实公共截线。两个例子说明复交数 \(2\) 可以对应两条实解，也可以对应一对非实共轭解，不能直接把复计数当作实计数。
\end{solution}

% !TeX root = ../../Book2.tex
% 纲要标题：## 第22章　Chern–Weil 与次级示性形式
% 纲要位置：几何学计划纲要.md，第 1848 行。
% 正文状态：未撰写；本文件目前是框架。

% 纲要内容开始
% 
% > 本章目标与证明范围：在实或复系数上完整比较拓扑类与曲率代表。Chern–Simons 的转移构造给出明确选择依赖；Euler 形式为第三卷高维 Gauss–Bonnet 提供规范一致的输入。
% 
% 纲要内容结束

\chapter{Chern–Weil 与次级示性形式}
\label{g2:ch:22}
\BookChapterNumber{g2:ch:22}

\begin{chapterabstract}
待撰写
\end{chapterabstract}

\begin{learninggoals}
待撰写
\end{learninggoals}

待撰写

% !TeX root = ../../Book2.tex
% 纲要标题：### 22.1 不变多项式与曲率
% 纲要位置：几何学计划纲要.md，第 1852 行。
% 正文状态：未撰写；本文件目前是框架。

% 纲要内容开始
% 
% * 对称不变多项式及极化
% * 丛值外乘与迹
% * Bianchi 恒等式给出的闭性
% 
% 纲要内容结束

\section{不变多项式与曲率}
\label{g2:sec:2201}

待撰写

\subsection{对称不变多项式及极化}
\label{g2:subsec:2201:01}
待撰写

\subsection{丛值外乘与迹}
\label{g2:subsec:2201:02}
待撰写

\subsection{Bianchi 恒等式给出的闭性}
\label{g2:subsec:2201:03}
待撰写
% 纲要条目排为编号小节标题，并进入目录与书签；说明段仍保留于纲要注释中。

% !TeX root = ../../Book2.tex
% 纲要标题：### 22.2 转移形式与联络无关性
% 纲要位置：几何学计划纲要.md，第 1858 行。
% 正文状态：初稿；尚未正式审核。

% 纲要内容开始
% 
% * 联络路径及参数积分
% * 陈–Weil 转移公式
% * 自然性和陈–Weil 同态
% * 规范变换下的变化
% 
% 纲要内容结束

\section{转移形式与联络无关性}
\label{g2:sec:2202}

同一张丛可以选择许多联络，它们的曲率通常也不同。上一节证明了不变多项式作用于曲率后得到闭形式；要让这个形式代表丛本身的类，还须解释改变联络为何只增加一个恰当形式。办法是让联络随一个参数连续变化，把曲率代表的变化率写成外微分，再沿参数积分。所得的原始形式还保留了两条联络之间的比较信息，将成为次级示性形式的来源。

\subsection{联络路径及参数积分}
\label{g2:subsec:2202:01}
固定光滑主右矩阵 Lie 群丛 \(Q\to M\)，其结构群为 \(G\)，Lie 代数为 \(\mathfrak g\)。本节的流形有限维、Hausdorff、第二可数；参数族在闭区间的端点光滑，意思是局部可以光滑延拓到稍大的区间。第一卷第18.2节“主联络的仿射空间”命题说明，两条联络之差属于 \(\Omega^1(M;\operatorname{ad}Q)\)。恢复它在局部的含义：若联络矩阵为 \(A_0,A_1\)，则换架时
\[
	A_j'=g^{-1}A_jg+g^{-1}\mathrm dg\quad(j=0,1),
	\qquad a'=g^{-1}ag,\quad a=A_1-A_0.
\]
相减后微分项消去，故 \(a\) 是一张整体伴随丛值一形式。反过来，将任意这样的 \(a\) 加到一条联络上，仍满足联络的换架规律。因此
\begin{equation}\label{g2:eq:affine-connection-path}
	\nabla_t=\nabla_0+t a,\qquad A_t=A_0+t a,\qquad 0\le t\le1
\end{equation}
给出连接两条联络的路径。对酉联络或正交联络，差分别取值于斜 Hermitian 或反对称矩阵，这条路径也保持相应相容性。

\begin{proposition}\label{g2:prop:curvature-path-variation}
	设 \(Q\to M\) 为光滑主右矩阵群丛，\(\nabla_t\) 是其光滑联络路径，局部矩阵与曲率分别为 \(A_t\)、\(F_t=\mathrm dA_t+A_t\wedge A_t\)。记 \(\dot A_t=\partial_tA_t\)，并以 \(\mathrm D_t\) 表示伴随丛值形式的协变外微分。则 \(\dot A_t\) 是整体伴随丛值一形式，而且
	\begin{equation}\label{g2:eq:curvature-path-variation}
		\partial_tF_t=\mathrm D_t\dot A_t.
	\end{equation}
	若 \(P:\mathfrak g^k\to\mathbb K\) 是对称 \(\operatorname{Ad}(G)\) 不变 \(k\) 线性式，\(k\ge1\)、\(\mathbb K=\mathbb R\) 或 \(\mathbb C\)，则
	\begin{equation}\label{g2:eq:chern-weil-form-variation}
		\partial_tP(F_t,\ldots,F_t)
		=k\,\mathrm dP(\dot A_t,F_t,\ldots,F_t),
	\end{equation}
	右侧有 \(k-1\) 个曲率因子。
\end{proposition}
\begin{proof}
	取不随参数变化的局部截面。过渡函数 \(g\) 与 \(t\) 无关，对换架公式求导即得 \(\dot A_t'=g^{-1}\dot A_tg\)，所以它们确实粘成伴随丛值形式。逐项求导给出
	\[
		\partial_tF_t=\mathrm d\dot A_t+\dot A_t\wedge A_t+A_t\wedge\dot A_t.
	\]
	一形式的协变外微分为 \(\mathrm D_t\dot A_t=\mathrm d\dot A_t+A_t\wedge\dot A_t+\dot A_t\wedge A_t\)，这证明第一式。两个乘积项都是加号，是因为被微分的形式次数为一。

	每个 \(F_t\) 的次数为二，交换两个曲率因子不产生符号。因此 \(P\) 的对称性给出
	\[
		\partial_tP(F_t,\ldots,F_t)
		=kP(\mathrm D_t\dot A_t,F_t,\ldots,F_t).
	\]
	另一方面，\cref{g2:lem:invariant-form-differential}将普通外微分改写为各因子的协变外微分。作用于第一个因子得到上式右边除去 \(k\) 的部分；作用于其余因子的项都含有 \(\mathrm D_tF_t\)，由\cref{g2:prop:cw-bianchi}为零。故
	\[
		\mathrm dP(\dot A_t,F_t,\ldots,F_t)
		=P(\mathrm D_t\dot A_t,F_t,\ldots,F_t),
	\]
	完成第二式的证明。
\end{proof}

沿参数积分是对形式系数作通常的一元积分。更准确地，对光滑族 \(\alpha_t\in\Omega^r(M;\mathbb K)\)，定义
\[
	\left(\int_0^1\alpha_t\,\mathrm dt\right)_x(v_1,\ldots,v_r)
	=\int_0^1(\alpha_t)_x(v_1,\ldots,v_r)\,\mathrm dt.
\]
紧参数区间上的光滑依赖允许在局部坐标中逐项求导，因此积分仍为光滑 \(r\) 形式，并且与 \(\mathrm d_M\) 交换。这个积分不改变底空间形式的次数。

\ProseParagraphBreak

也可以把上述计算放到 \([0,1]\times M\) 上。令拉回丛在参数方向没有联络分量，则它的曲率为
\[
	\widehat F=F_t+\mathrm dt\wedge\dot A_t.
\]
我们将参数微分写在前面。由于 \((\mathrm dt)^2=0\)，有
\begin{equation}\label{g2:eq:chern-weil-cylinder-form}
	P(\widehat F,\ldots,\widehat F)
	=P(F_t,\ldots,F_t)+\mathrm dt\wedge kP(\dot A_t,F_t,\ldots,F_t).
\end{equation}
总空间上的闭性在 \(\mathrm dt\) 分量中恰好给出式\eqref{g2:eq:chern-weil-form-variation}。这样，联络变化率成为参数方向的一份曲率，而转移形式将由含有这份曲率的项积分得到。

\subsection{陈–Weil 转移公式}
\label{g2:subsec:2202:02}
\begin{definition}\label{g2:def:chern-weil-transgression}
	设 \(\Gamma=(\nabla_t)_{0\le t\le1}\) 是固定光滑主 \(G\) 丛上的光滑联络路径，\(P\) 是次数 \(k\ge1\) 的对称不变多项式。形式
	\begin{equation}\label{g2:eq:chern-weil-transgression}
		T_P(\Gamma)=k\int_0^1P(\dot A_t,F_t,\ldots,F_t)\,\mathrm dt
		\in\Omega^{2k-1}(M;\mathbb K)
	\end{equation}
	称为这条路径的\term[zhuanyixingshi]{转移形式}[transgression form]。对式\eqref{g2:eq:affine-connection-path}的仿射路径，简记为 \(T_P(\nabla_0,\nabla_1)\)。
\end{definition}

这里每个被积项都是整体标量形式：\(\dot A_t\) 和 \(F_t\) 都按共轭变换，而 \(P\) 不变。其次数为 \(1+2(k-1)=2k-1\)，比曲率代表低一次。

\begin{theorem}[陈--Weil 转移公式]\label{g2:thm:chern-weil-transgression}
	设 \(Q\to M\) 为光滑主右矩阵群丛，\(P\) 为次数 \(k\ge1\) 的实或复对称不变多项式，\(\Gamma=(\nabla_t)\) 为连接 \(\nabla_0,\nabla_1\) 的光滑联络路径。则
	\begin{equation}\label{g2:eq:chern-weil-transgression-differential}
		\mathrm dT_P(\Gamma)=P(F_1,\ldots,F_1)-P(F_0,\ldots,F_0).
	\end{equation}
	因此闭形式 \(P(F_\nabla,\ldots,F_\nabla)\) 的 de Rham 上同调类与联络无关。
\end{theorem}
\begin{proof}
	参数积分与外微分交换，故\cref{g2:prop:curvature-path-variation}给出
	\[
	\begin{aligned}
		\mathrm dT_P(\Gamma)
		&=\int_0^1 k\,\mathrm dP(\dot A_t,F_t,\ldots,F_t)\,\mathrm dt\\
		&=\int_0^1\partial_tP(F_t,\ldots,F_t)\,\mathrm dt\\
		&=P(F_1,\ldots,F_1)-P(F_0,\ldots,F_0).
	\end{aligned}
	\]
	最后一步是逐个形式系数使用微积分基本定理。任意两条联络可由仿射路径连接，右侧总为恰当形式；结合\cref{g2:thm:chern-weil-closed}的闭性，得到上同调类与联络无关。
\end{proof}

例如复线丛上取 \(P(X)=\mathrm iX/(2\pi)\)。由于 \(k=1\)，转移形式就是 \(\mathrm i(\nabla_1-\nabla_0)/(2\pi)\)，曲率之差是它的外微分。高次数中的额外曲率因子，使同一思路能比较所有陈--Weil 形式。

\ProseParagraphBreak

转移形式本身仍可能随连接路径改变。其可控程度需要一个二参数计算；仅知道两份转移形式的外微分相同，只能推出差闭，还不能推出差恰当。

\begin{proposition}\label{g2:prop:transgression-path-homotopy}
	设 \(Q\to M\) 为光滑主右矩阵群丛，\(P\) 为次数 \(k\ge1\) 的对称不变多项式。若两条光滑联络路径 \(\Gamma_0,\Gamma_1\) 具有相同起点与终点，则 \(T_P(\Gamma_1)-T_P(\Gamma_0)\) 为恰当形式。特别地，对三条联络 \(\nabla_0,\nabla_1,\nabla_2\)，
	\begin{equation}\label{g2:eq:transgression-triangle-mod-exact}
		T_P(\nabla_0,\nabla_2)
		\equiv T_P(\nabla_0,\nabla_1)+T_P(\nabla_1,\nabla_2)
		\pmod{\mathrm d\Omega^{2k-2}(M;\mathbb K)}.
	\end{equation}
	仿射转移还严格满足 \(T_P(\nabla_1,\nabla_0)=-T_P(\nabla_0,\nabla_1)\)。
\end{proposition}
\begin{proof}
	由于联络空间是仿射空间，可用
	\[
		\nabla_{s,t}=(1-s)\nabla_t^{\,0}+s\nabla_t^{\,1}
	\]
	连接两条路径，而且在 \(t=0,1\) 处不随 \(s\) 变化。记局部矩阵为 \(A_{s,t}\)，令 \(a=\partial_tA\)、\(b=\partial_sA\)、\(F=\mathrm d_MA+A\wedge A\)。在方形与 \(M\) 的乘积上取无参数分量的联络，其曲率为
	\[
		\widehat F=F+\mathrm ds\wedge b+\mathrm dt\wedge a.
	\]
	展开 \(P(\widehat F^k)\)，写成
	\[
		\omega+\mathrm ds\wedge U+\mathrm dt\wedge V
		+\mathrm ds\wedge\mathrm dt\wedge W,
	\]
	其中
	\[
	\begin{aligned}
		\omega&=P(F^k),&
		U&=kP(b,F^{k-1}),\\
		V&=kP(a,F^{k-1}),&
		W&=k(k-1)P(a,b,F^{k-2}).
	\end{aligned}
	\]
	\(k=1\) 时约定 \(W=0\)。这里 \(F^j\) 表示连续放入 \(j\) 个曲率因子。核对最后一项的符号：\((\mathrm ds\wedge b)\wedge(\mathrm dt\wedge a)=-\mathrm ds\wedge\mathrm dt\wedge b\wedge a\)，再交换两个一形式 \(a,b\) 又产生一个负号，故得到上述 \(W\)。

	总外微分为零。比较 \(\mathrm ds\wedge\mathrm dt\) 的系数，得到
	\[
		\partial_sV-\partial_tU+\mathrm d_MW=0.
	\]
	在方形上积分，并用固定端点条件 \(b(s,0)=b(s,1)=0\)，有
	\[
	\begin{aligned}
		T_P(\Gamma_1)-T_P(\Gamma_0)
		&=\int_0^1\int_0^1\partial_sV\,\mathrm dt\,\mathrm ds\\
		&=\int_0^1\bigl(U(s,1)-U(s,0)\bigr)\,\mathrm ds
		-\mathrm d_M\int_0^1\int_0^1W\,\mathrm dt\,\mathrm ds\\
		&=-\mathrm d_M\int_0^1\int_0^1W\,\mathrm dt\,\mathrm ds.
	\end{aligned}
	\]
	这给出了差的显式原始形式。

	为证明三角式，先沿 \(\nabla_0\) 到 \(\nabla_1\) 的路径走，再沿 \(\nabla_1\) 到 \(\nabla_2\) 走。作保向参数变换不会改变式\eqref{g2:eq:chern-weil-transgression}，因为 \(\dot A_t\,\mathrm dt\) 按一元换元公式变换。可将两段在端点附近重新参数化为常值，再光滑拼接；积分是两段积分之和。将这条路径与 \(\nabla_0\) 到 \(\nabla_2\) 的仿射路径比较，刚才的二参数计算给出式\eqref{g2:eq:transgression-triangle-mod-exact}。最后，反向走仿射路径使速度换号，换元后直接得到严格的反对称性。
\end{proof}

因此，两条固定联络之间有一个与路径无关的元素
\[
	[T_P(\nabla_0,\nabla_1)]
	\in\Omega^{2k-1}(M;\mathbb K)/\mathrm d\Omega^{2k-2}(M;\mathbb K).
\]
这是形式模去恰当形式的商；只有当两端的曲率代表相等时，转移形式才闭，才能进一步称为 de Rham 上同调类。

\subsection{自然性和陈–Weil 同态}
\label{g2:subsec:2202:03}
将不变多项式的次数加倍，恰好对应曲率形式的次数。记
\[
	I^*(G;\mathbb K)=\bigoplus_{k\ge0}I^k(G;\mathbb K)
\]
为对称不变多项式代数，其中 \(I^k\) 的元素是次数 \(k\) 的齐次多项式，并与其极化后的对称 \(k\) 线性式对应。乘积取普通多项式乘积；极化给出的组合系数已包含在乘积定义中。

\begin{theorem}[陈--Weil 同态]\label{g2:thm:chern-weil-homomorphism}
	设 \(G\) 为矩阵 Lie 群，\(Q\to M\) 为有限维、Hausdorff、第二可数光滑流形上的光滑主右 \(G\) 丛，\(\mathbb K=\mathbb R\) 或 \(\mathbb C\)。对任意光滑联络 \(\nabla\)，令
	\begin{equation}\label{g2:eq:chern-weil-homomorphism}
		\operatorname{cw}_Q(P)=[P(F_\nabla,\ldots,F_\nabla)].
	\end{equation}
	这定义了与联络无关的保单位代数同态
	\[
		\operatorname{cw}_Q:I^*(G;\mathbb K)\longrightarrow
		H_{\mathrm{dR}}^{\mathrm{even}}(M;\mathbb K),
		\qquad I^k\longrightarrow H_{\mathrm{dR}}^{2k}.
	\]
	对光滑映射 \(f:N\to M\)，它满足
	\begin{equation}\label{g2:eq:chern-weil-pullback}
		\operatorname{cw}_{f^*Q}(P)=f^*\operatorname{cw}_Q(P).
	\end{equation}
	对结构群同态 \(\rho:H\to G\) 与主 \(H\) 丛 \(R\to M\)，若 \(Q=R\times_\rho G\)，则
	\[
		\operatorname{cw}_{Q}(P)=\operatorname{cw}_{R}(\rho^*P),
		\qquad
		(\rho^*P)(X_1,\ldots,X_k)=P(\rho_*X_1,\ldots,\rho_*X_k).
	\]
\end{theorem}
\begin{proof}
	光滑主丛存在联络：局部平凡丛先取零联络，用从属光滑单位分解在总空间上加权；权重之和为一保持竖直再现条件，底空间上的权重保持等变性。这是第一卷第18.2节“主联络的存在”定理的构造。闭性由\cref{g2:thm:chern-weil-closed}保证，联络无关性由\cref{g2:thm:chern-weil-transgression}保证，故式\eqref{g2:eq:chern-weil-homomorphism}有定义。

	固定一条联络。对多项式相加、乘标量，代入曲率的运算相应相加、乘标量。对次数分别为 \(k,\ell\) 的 \(P_1,P_2\)，在极化后的乘积中，每项都把 \(k+\ell\) 个二次曲率因子分给两组。因这些因子均为偶次，重排不产生符号；极化的平均系数又使全部相同项的系数和为一。因此
	\[
		(P_1P_2)(F^{k+\ell})=P_1(F^k)\wedge P_2(F^\ell).
	\]
	常数多项式 \(1\) 给出常数零形式 \(1\)，从而得到保单位代数同态。

	拉回联络在局部为 \(\mathrm d+f^*A\)。因为拉回与外微分、外乘交换，其曲率是 \(f^*F\)，所以
	\[
		P(F_{f^*\nabla}^k)=f^*P(F_\nabla^k).
	\]
	在源流形上若另选一条联络，转移公式使所得上同调类不变，于是得到自然性。

	最后，沿 \(\rho\) 扩张结构群后的局部联络矩阵为 \(\rho_*A\)。Lie 代数同态保持括号，故相应曲率为 \(\rho_*F\)。将它代入 \(P\) 得到 \((\rho^*P)(F^k)\)，证明最后一式。
\end{proof}

自然性也在转移形式的层次成立。若拉回整条路径，则
\begin{equation}\label{g2:eq:transgression-pullback}
	T_P(f^*\Gamma)=f^*T_P(\Gamma),
\end{equation}
因为速度、曲率和参数积分都与拉回交换。这是一条形式等式，比上同调类自然性更精细。下一节将选取确定的行列式及 Pfaffian 多项式，把式\eqref{g2:eq:chern-weil-homomorphism}与上一章的拓扑类逐一比较。

\subsection{规范变换下的变化}
\label{g2:subsec:2202:04}
规范变换是覆盖底空间恒等映射的主丛自同构。沿用第一卷第18.4节“联络的规范变换”命题，在固定局部截面中写成
\[
	A^g=g^{-1}Ag+g^{-1}\mathrm dg,\qquad F^g=g^{-1}Fg.
\]
这里的局部群值函数满足伴随换架规律，因而表示同一个整体规范变换；平凡丛中它就是一张整体映射 \(g:M\to G\)。不变多项式立即给出 \(P((F^g)^k)=P(F^k)\)，是代表形式的严格相等。

\begin{proposition}\label{g2:prop:transgression-gauge-naturality}
	设 \(Q\to M\) 为光滑主右矩阵群丛，\(P\) 为次数 \(k\ge1\) 的对称不变多项式，\(g\) 为固定规范变换。对任意联络路径 \(\Gamma\)，同时变换路径上的所有联络，有
	\[
		T_P(\Gamma^g)=T_P(\Gamma).
	\]
	因此仿射转移满足
	\begin{equation}\label{g2:eq:transgression-gauge-invariant-pair}
		T_P(\nabla_0^g,\nabla_1^g)=T_P(\nabla_0,\nabla_1).
	\end{equation}
	只比较一条联络及其变换时，\(T_P(\nabla,\nabla^g)\) 是闭形式，其上同调类与 \(\nabla\) 无关。
\end{proposition}
\begin{proof}
	\(g\) 不随路径参数改变，所以 \(\partial_tA_t^g=g^{-1}\dot A_tg\)，而 \(F_t^g=g^{-1}F_tg\)。将这些式子代入 \(P\)，共轭由不变性消去，被积形式逐项不变。规范变换还把仿射路径送到变换后两端之间的仿射路径，因为附加项 \(g^{-1}\mathrm dg\) 的系数为一，故得到式\eqref{g2:eq:transgression-gauge-invariant-pair}。

	转移公式给出
	\[
		\mathrm dT_P(\nabla,\nabla^g)=P((F^g)^k)-P(F^k)=0.
	\]
	再取另一联络 \(\nabla'\)。用\cref{g2:prop:transgression-path-homotopy}沿四个端点依次比较，模去恰当形式有
	\[
	\begin{aligned}
		T_P(\nabla',(\nabla')^g)
		&\equiv T_P(\nabla',\nabla)+T_P(\nabla,\nabla^g)
		+T_P(\nabla^g,(\nabla')^g)\\
		&=-T_P(\nabla,\nabla')+T_P(\nabla,\nabla^g)
		+T_P(\nabla,\nabla')\\
		&=T_P(\nabla,\nabla^g).
	\end{aligned}
	\]
	两个外侧项由同时规范不变性及反对称性抵消，所以闭转移类与所选联络无关。
\end{proof}

\begin{example}\label{g2:exm:circle-gauge-transgression}
	取 \(S^1=\mathbb R/\mathbb Z\) 上的平凡酉线丛，\(\nabla=\mathrm d\)，并令 \(g(x)=e^{2\pi\mathrm i m x}\)，其中 \(m\in\mathbb Z\)。角度 \(x\) 本身不是圆周上的实函数，但 \(\mathrm dx\) 是整体一形式，且正向积分为 \(1\)。于是
	\[
		A^g=g^{-1}\mathrm dg=2\pi\mathrm i m\,\mathrm dx.
	\]
	对 \(P(X)=\mathrm iX/(2\pi)\)，有
	\[
		T_P(\nabla,\nabla^g)=\frac{\mathrm i}{2\pi}(A^g-0)
		=-m\,\mathrm dx,\qquad
		\int_{S^1}T_P(\nabla,\nabla^g)=-m.
	\]
	\(m\ne0\) 时，这个闭形式不是恰当形式，因为闭曲线上恰当一形式的积分为零。两端曲率都为零，但规范变换的绕数仍留在转移形式中。
\end{example}

因此，曲率代表的规范不变性与转移形式的选择依赖各自有准确的含义：同时改变一对联络，转移形式保持相等；固定一端而用规范变换改变另一端，可以留下非零闭类。后面的三维例子将把同样的现象写成整数改变的积分，从而得到模整数的不变量。
% 纲要条目排为编号小节标题，并进入目录与书签；说明段仍保留于纲要注释中。

% !TeX root = ../../Book2.tex
% 纲要标题：### 22.3 拓扑类的比较
% 纲要位置：几何学计划纲要.md，第 1815 行。
% 正文状态：未撰写；本文件目前是框架。

% 纲要内容开始
% 
% * 线丛上的归一化
% * 分裂原理及一般 Chern 类
% * Pontryagin 形式与 Euler 的 Pfaffian
% * 曲率形式看不到的扭信息
% 
% 纲要内容结束

\section{拓扑类的比较}
\label{g2:sec:2203}

待撰写

\subsection{线丛上的归一化}
\label{g2:subsec:2203:01}
待撰写

\subsection{分裂原理及一般 Chern 类}
\label{g2:subsec:2203:02}
待撰写

\subsection{Pontryagin 形式与 Euler 的 Pfaffian}
\label{g2:subsec:2203:03}
待撰写

\subsection{曲率形式看不到的扭信息}
\label{g2:subsec:2203:04}
待撰写
% 纲要条目排为编号小节标题，并进入目录与书签；说明段仍保留于纲要注释中。

% !TeX root = ../../Book2.tex
% 纲要标题：### 22.4 陈–Simons 与几何应用
% 纲要位置：几何学计划纲要.md，第 1872 行。
% 正文状态：初稿；尚未正式审核。

% 纲要内容开始
% 
% * 陈–Simons 形式
% * 陈–Simons 积分与规范变换
% * 边界、延拓和相对不变量
% * 陈–Simons 变分与陈特征
% 
% ---
% 
% 纲要内容结束

\section{陈–Simons 与几何应用}
\label{g2:sec:2204}

两个联络给出同一个示性上同调类，二者的曲率形式却未必相等。\secref{g2:sec:2202}的转移公式把这个差写成一个外微分；现在保留那个次数低一的形式，便得到能随联络变化的几何量。在三维，二次不变多项式的四次曲率形式因维数消失，三次转移形式仍可以积分。问题随之变成：这个积分依赖哪些选择，改变标架时为何会出现整数？

\subsection{陈–Simons 形式}
\label{g2:subsec:2204:01}
设 \(P\) 是次数为 \(k\ge1\) 的对称不变多项式，\(\nabla_0,\nabla_1\) 是同一张光滑丛上的两个联络。沿直线路径 \(\nabla_t=\nabla_0+t a\)，其中 \(a=\nabla_1-\nabla_0\)，定义
\begin{equation}\label{g2:eq:relative-chern-simons-form}
	\operatorname{CS}_P(\nabla_0,\nabla_1)
	\coloneqq k\int_0^1P(a,F_t,\ldots,F_t)\,\mathrm dt.
\end{equation}
这就是\cref{g2:thm:chern-weil-transgression}中的转移形式，称为相对于 \(\nabla_0\) 的\term[chen-Simons-xingshi]{陈--Simons 形式}[Chern--Simons form]\footnote{James Harris Simons（1938--2024），美国数学家，研究微分几何，并与陈省身共同发展次级示性不变量。陈省身的简介见第21章。}。它是底空间上的整体 \((2k-1)\) 次形式，并满足
\[
	\mathrm d\operatorname{CS}_P(\nabla_0,\nabla_1)=P(F_1^k)-P(F_0^k).
\]
若右边为零，则它给出 de Rham 上同调类。这个条件可能来自两个曲率代表相等，也可能来自底空间维数小于 \(2k\)。

\ProseParagraphBreak

这里的参照联络是定义的一部分。\cref{g2:prop:transgression-path-homotopy}说明，保持两端点而更换联络路径只改变一个恰当形式；对三个联络，还有
\begin{equation}\label{g2:eq:chern-simons-reference-change}
	\operatorname{CS}_P(\nabla_0,\nabla_2)
	\equiv\operatorname{CS}_P(\nabla_0,\nabla_1)+\operatorname{CS}_P(\nabla_1,\nabla_2)
	\pmod{\mathrm d\Omega^{2k-2}}.
\end{equation}
因此换参照联络时，变化由另一份转移形式决定。若原丛已有一个全局标架，可以取该标架下的平凡联络 \(\mathrm d\) 为参照；换一个随点变化的标架，也就换了这个参照联络。这正是下面规范变换公式的来源。次级形式的原始构造及选择依赖可参阅\cite[第3节]{ChernSimons1974}。

\begin{proposition}\label{g2:prop:su-two-chern-simons-form}
	设 \(A\in\Omega^1(M;\mathfrak{su}(2))\) 为光滑流形 \(M\) 上平凡 \(\operatorname{SU}(2)\) 丛在一个固定标架中的联络矩阵，\(F=\mathrm dA+A\wedge A\)。定义
	\begin{equation}\label{g2:eq:su-two-chern-simons-density}
		\operatorname{CS}(A)=\frac1{8\pi^2}\operatorname{tr}\left(A\wedge\mathrm dA+\frac23A\wedge A\wedge A\right).
	\end{equation}
	则 \(\operatorname{CS}(A)=\operatorname{CS}_{c_2}(\mathrm d,\mathrm d+A)\)，且
	\[
		\mathrm d\operatorname{CS}(A)=\frac1{8\pi^2}\operatorname{tr}(F\wedge F)=c_2(\nabla).
	\]
\end{proposition}
\begin{proof}
	\(\mathfrak{su}(2)\) 由无迹反 Hermitian 矩阵组成，故 \(\operatorname{tr}F=0\)。本章采用 \(c(\nabla)=\det(I+\mathrm iF/(2\pi))\)，其二次项为
	\[
		c_2(\nabla)=\frac1{8\pi^2}\left(\operatorname{tr}(F^2)-(\operatorname{tr}F)^2\right)=\frac1{8\pi^2}\operatorname{tr}(F^2).
	\]
	相应极化为 \(P(X,Y)=\operatorname{tr}(XY)/(8\pi^2)\)。从零联络到 \(A\) 的路径是 \(A_t=tA\)，曲率为 \(F_t=t\mathrm dA+t^2A^2\)。代入式\eqref{g2:eq:relative-chern-simons-form}，得到
	\[
		2\int_0^1\frac1{8\pi^2}\operatorname{tr}\bigl(A(t\mathrm dA+t^2A^2)\bigr)\,\mathrm dt
		=\frac1{8\pi^2}\operatorname{tr}\left(A\mathrm dA+\frac23A^3\right).
	\]
	转移恒等式给出它的外微分。这里的乘法均含外乘；迹的分次循环律保证该计算与矩阵的非交换乘法相容。
\end{proof}

若 \(\dim M=3\)，四形式 \(c_2(\nabla)\) 因维数恒为零，因而 \(\operatorname{CS}(A)\) 闭。其积分却可以随 \(A\) 连续变化：四次拓扑类在这个维数不提供数值，三次次级形式还记得联络。

\subsection{陈–Simons 积分与规范变换}
\label{g2:subsec:2204:02}
现在固定一个连通闭定向光滑三维流形 \(M\)。在平凡 \(\operatorname{SU}(2)\) 丛的一张标架中，令
\[
	S(A)=\int_M\operatorname{CS}(A)\in\mathbb R.
\]
为了解它在换标架后的变化，先在群本身上计算一个规范化三次形式。

\begin{lemma}\label{g2:lem:su-two-integral-three-form}
	将 \(\operatorname{SU}(2)\) 与单位三维球面识别为
	\[
		q=(x_0,x_1,x_2,x_3)\longmapsto
		\begin{pmatrix}x_0+\mathrm i x_1&x_2+\mathrm i x_3\\-x_2+\mathrm i x_3&x_0-\mathrm i x_1\end{pmatrix},
		\qquad\sum_{j=0}^3x_j^2=1.
	\]
	给 \(S^3\) 使用 \(\mathbb R^4\) 中单位球的外法向优先边界定向，记由左平移定义的 Lie 代数值一形式为 \(\vartheta=q^{-1}\mathrm dq\)。则
	\begin{equation}\label{g2:eq:su-two-normalized-volume}
		\omega=-\frac1{24\pi^2}\operatorname{tr}(\vartheta^3),\qquad \int_{S^3}\omega=1.
	\end{equation}
	因此对任意光滑映射 \(g:M\to\operatorname{SU}(2)\)，有 \(\int_Mg^*\omega=\deg g\in\mathbb Z\)。
\end{lemma}
\begin{proof}
	在单位元处，与坐标 \(x_1,x_2,x_3\) 对应的切向量为
	\[
		I=\begin{pmatrix}\mathrm i&0\\0&-\mathrm i\end{pmatrix},\qquad
		J=\begin{pmatrix}0&1\\-1&0\end{pmatrix},\qquad
		K=\begin{pmatrix}0&\mathrm i\\\mathrm i&0\end{pmatrix}.
	\]
	它们满足 \(I^2=J^2=K^2=-1\)、\(IJ=K\)、\(JK=I\)、\(KI=J\)，逆序乘积变号。三次外乘在 \((I,J,K)\) 上取值是六个交错矩阵乘积之和；每个交错乘积都等于 \(-1\)，所以
	\[
		\operatorname{tr}(\vartheta^3)(I,J,K)=6\operatorname{tr}(-1)=-12.
	\]
	在点 \((1,0,0,0)\)，这组切基为球面的正向标准正交基。左乘单位四元数保持 \(\mathbb R^4\) 的 Euclid 范数，因而保持球面圆度量；它随单位四元数连续变化，且恒等左乘保向，所以全部左乘保向。\(\vartheta\) 与圆体积形式均左不变，刚才在单位元的计算遂给出
	\[
		-\frac1{24\pi^2}\operatorname{tr}(\vartheta^3)=\frac1{2\pi^2}\,\mathrm{vol}_{S^3}.
	\]
	为核对常数，球坐标 \((\cos\chi,\sin\chi\,u)\)，\(u\in S^2\)、\(0\le\chi\le\pi\)，给出体积
	\[
		\operatorname{vol}(S^3)=\operatorname{area}(S^2)\int_0^\pi\sin^2\chi\,\mathrm d\chi=4\pi\cdot\frac\pi2=2\pi^2.
	\]
	故 \(\omega\) 的积分为一。最后使用\cref{g2:thm:degree-comparison}：该定理在闭连通同维定向流形之间将归一化顶次形式的拉回积分识别为基本类所定义的整数映射度。这里两端恰为三维，故所有条件满足。
\end{proof}

\begin{theorem}\label{g2:thm:chern-simons-gauge-integer}
	设 \(M\) 为连通闭定向光滑三维流形，\(A\) 为平凡 \(\operatorname{SU}(2)\) 丛的一张标架下的光滑联络矩阵。对任意光滑规范变换 \(g:M\to\operatorname{SU}(2)\)，令
	\[
		A^g=g^{-1}Ag+g^{-1}\mathrm dg,\qquad \theta=g^{-1}\mathrm dg,\qquad B=g^{-1}Ag.
	\]
	则有形式恒等式
	\begin{equation}\label{g2:eq:chern-simons-gauge-formula}
		\operatorname{CS}(A^g)-\operatorname{CS}(A)
		=-\frac1{8\pi^2}\mathrm d\operatorname{tr}(\theta\wedge B)+g^*\omega,
	\end{equation}
	并且
	\begin{equation}\label{g2:eq:chern-simons-integer-shift}
		S(A^g)-S(A)=\deg g.
	\end{equation}
	于是 \(S(A)\bmod\mathbb Z\in\mathbb R/\mathbb Z\) 与全局标架无关，并在规范变换下不变。
\end{theorem}
\begin{proof}
	记 \(R=g^{-1}(\mathrm dA)g\)。微分乘积法则和 \(\mathrm d(g^{-1})=-g^{-1}(\mathrm dg)g^{-1}\) 给出
	\[
		\mathrm d\theta=-\theta^2,\qquad \mathrm dB=R-\theta B-B\theta.
	\]
	为把所有符号看清，先省略公因子 \(1/(8\pi^2)\)。迹中三个一次形式作循环置换时符号为正，因此展开 \(B+\theta\) 的三次项后，三个含两个 \(B\) 的项相同，三个含两个 \(\theta\) 的项也相同。于是
	\[
	\begin{aligned}
		&\operatorname{tr}\left((B+\theta)\mathrm d(B+\theta)+\frac23(B+\theta)^3\right)
		-\operatorname{tr}\left(A\mathrm dA+\frac23A^3\right)\\
		&\qquad=\operatorname{tr}(\theta R-B\theta^2)-\frac13\operatorname{tr}(\theta^3).
	\end{aligned}
	\]
	另一方面，
	\[
		\mathrm d\operatorname{tr}(\theta B)
		=\operatorname{tr}(-\theta^2B-\theta(R-\theta B-B\theta))
		=-\operatorname{tr}(\theta R)+\operatorname{tr}(B\theta^2).
	\]
	代回即得式\eqref{g2:eq:chern-simons-gauge-formula}，因为 \(-\operatorname{tr}(\theta^3)/(24\pi^2)=g^*\omega\)。在闭流形上，恰当三次形式的积分由 Stokes 定理为零；\cref{g2:lem:su-two-integral-three-form}将剩下的一项化为 \(\deg g\)。因此变换前后的实数相差整数，模整数的值不变。
\end{proof}

取 \(M=S^3=\operatorname{SU}(2)\)、\(A=0\) 和 \(g(q)=q\)，便得到
\[
	A^g=\vartheta,\qquad F_{A^g}=0,\qquad S(0)=0,\qquad S(\vartheta)=1.
\]
两个平坦联络在规范变换下等价，但在各自矩阵表达中得到的实数不同。整数模糊由换标架的同伦类产生；把数值取到 \(\mathbb R/\mathbb Z\) 后，这两份表达才一致。

\subsection{边界、延拓和相对不变量}
\label{g2:subsec:2204:03}
有边界时，式\eqref{g2:eq:chern-simons-gauge-formula}中的恰当项会留下边界积分。另一方面，在四维内部积分 \(c_2(\nabla)\)，也自然要求在三维边界记录一个陈--Simons 修正。下面把这两件事组织为相对不变量。所有边界都取外法向优先定向，所用 Stokes 定理是第一卷第14.4节的版本：对紧有向光滑流形及光滑至边界的形式，\(\int_W\mathrm d\eta=\int_{\partial W}\eta\)。

\begin{proposition}\label{g2:prop:relative-chern-simons-integer}
	设 \(W\) 为紧定向光滑四维流形，\(M=\partial W\)，\(E\to W\) 为具有 \(\operatorname{SU}(2)\) 结构的复秩二光滑丛，并给边界丛选一张 \(\operatorname{SU}(2)\) 标架 \(s\)。对任意相容光滑联络 \(\nabla\)，记其边界矩阵为 \(A_s\)。则
	\begin{equation}\label{g2:eq:relative-second-chern-number}
		N(W,E,s)=\int_W c_2(\nabla)-\int_M\operatorname{CS}(A_s)
	\end{equation}
	是一个只依赖 \((W,E,s)\) 的整数。若 \(s'=sg\)，其中 \(g:M\to\operatorname{SU}(2)\)，则
	\[
		N(W,E,s')=N(W,E,s)-\sum_{M_\alpha\subset M}\deg(g|_{M_\alpha}).
	\]
\end{proposition}
\begin{proof}
	\cref{g2:thm:collar-neighborhood}在紧光滑边界附近给出乘积领圈。先把边界标架延到这条领圈：沿领圈的法向区间作平行移动，便把边界上的丛识别为领圈上的乘积。于是 \(s\) 在领圈上确定平凡联络。联络构成仿射空间，用只依赖领参数的光滑截断将给定联络与这个平凡联络插值，可得到一个联络 \(\nabla^0\)，使它在较小领圈上恰为 \(\mathrm d\)。

	对 \(\nabla^0\) 与 \(\nabla\) 使用转移公式，再用 Stokes 定理，得到
	\[
	\begin{aligned}
		\int_Wc_2(\nabla)-\int_Wc_2(\nabla^0)
		&=\int_M\operatorname{CS}_{c_2}(\nabla^0|_M,\nabla|_M)\\
		&=\int_M\operatorname{CS}(A_s).
	\end{aligned}
	\]
	因此式\eqref{g2:eq:relative-second-chern-number}等于 \(\int_Wc_2(\nabla^0)\)。将 \(W\) 与反向副本 \(-W\) 沿边界粘合为闭定向四维流形 \(X\)。在第二份 \(-W\) 上取平凡复秩二丛，并用 \(s\) 与第一份 \(E\) 粘合，得到 \(E_X\)。两边联络在领圈标架中都为零，故光滑粘合为 \(\nabla_X\)。于是
	\[
		\int_Wc_2(\nabla^0)=\int_Xc_2(\nabla_X)
		=\langle c_2(E_X),[X]\rangle\in\mathbb Z.
	\]
	最后一个等式是\cref{g2:thm:chern-weil-chern-comparison}，其中整数陈类先映到实上同调，再由积分与基本类配对。闭合后的丛 \(E_X\) 由边界粘合数据确定；改变领圈识别只给出同构的丛，因为两种向内延拓可沿领参数收缩回同一边界标架。因此右侧与联络及辅助选择无关。

	改用 \(s'=sg\) 时，内部曲率积分不变，边界矩阵由 \(A_s\) 变为 \(A_s^g\)。对每个边界分支应用式\eqref{g2:eq:chern-simons-integer-shift}，便得到所列变化公式。边界方向已进入各分支映射度的定义。
\end{proof}

特别地，如果边界上的联络有一份四维延拓，那么
\begin{equation}\label{g2:eq:chern-simons-extension-mod-integer}
	S(A_s)\equiv\int_Wc_2(\nabla)\pmod{\mathbb Z}.
\end{equation}
两份延拓的内部积分相差整数。也可把两份填充沿共同边界粘合，第二份反向，差便等于闭四维流形上的陈数。若需先使联络在领圈成为边界联络的乘积，可用保持边界值的插值；转移形式在边界的限制为零，所以此调整不改变四维积分。这个论证说明，整性来自闭合后的整数拓扑类，而边界联络仍决定模整数的连续数值。

\ProseParagraphBreak

相对上同调也直接记录同一个修正。对包含 \(j:M\hookrightarrow W\)，采用复形
\[
	\mathcal C^r=\Omega^r(W)\oplus\Omega^{r-1}(M),\qquad
	\mathrm D(\alpha,\beta)=(\mathrm d\alpha,j^*\alpha-\mathrm d\beta).
\]
有 \(\mathrm D^2=0\)。这个复形与边界限制为零的相对形式复形计算同一上同调：边界形式可沿领圈延拓并用截断补零，因而限制映射 \(\Omega^*(W)\to\Omega^*(M)\) 逐次满射；由核的短正合列与上述复形的长正合列比较，恒等映射在外侧给出同构，中间也为同构。结合\cref{g2:thm:de-rham-support}，它就是 \(H^*(W,M;\mathbb R)\) 的一个模型。

\ProseParagraphBreak

成对形式 \((c_2(\nabla),\operatorname{CS}(A_s))\) 是闭元，因为边界上 \(\mathrm d\operatorname{CS}(A_s)=j^*c_2(\nabla)\)。其相对求值为
\[
	\langle(\alpha,\beta),[W,M]\rangle=\int_W\alpha-\int_M\beta.
\]
对一个余边界 \(\mathrm D(\eta,\xi)\)，右端由两次 Stokes 公式为零，所以只依赖相对类。式\eqref{g2:eq:relative-second-chern-number}正是这个求值；边界的平凡化说明了为什么这里需要相对数据。

\subsection{陈–Simons 变分与陈特征}
\label{g2:subsec:2204:04}
三维陈--Simons 积分首先给出一个具体的变分问题。将所有 \(\mathfrak{su}(2)\) 值一次形式看成平凡丛的联络空间，允许它们作任意光滑变化。

\begin{proposition}\label{g2:prop:chern-simons-variation}
	设 \(M\) 为闭定向光滑三维流形，\(A\) 为平凡 \(\operatorname{SU}(2)\) 丛的联络矩阵，\(a\in\Omega^1(M;\mathfrak{su}(2))\)。则
	\begin{equation}\label{g2:eq:chern-simons-variation}
		\left.\frac{\mathrm d}{\mathrm dt}\right|_{t=0}S(A+ta)
		=\frac1{4\pi^2}\int_M\operatorname{tr}(a\wedge F_A).
	\end{equation}
	因此 \(S\) 的临界联络恰为平坦联络。
\end{proposition}
\begin{proof}
	对式\eqref{g2:eq:su-two-chern-simons-density}求导，用迹的分次循环律，得到
	\[
	\begin{aligned}
		\delta\operatorname{tr}\left(A\mathrm dA+\frac23A^3\right)
		&=\operatorname{tr}(a\mathrm dA+A\mathrm da+2aA^2)\\
		&=2\operatorname{tr}(aF_A)-\mathrm d\operatorname{tr}(Aa).
	\end{aligned}
	\]
	闭流形上的最后一项积分为零，故得到变分公式。若 \(F_A=0\)，所有变分都为零。反过来，给 \(M\) 临时选一个 Riemann 度量，并对 \(\mathfrak{su}(2)\) 取正定内积 \(\langle X,Y\rangle=-\operatorname{tr}(XY)\)。对实二次形式定义一次形式 \(*\gamma\)，使 \(\beta\wedge *\gamma=\langle\beta,\gamma\rangle\mathrm{vol}_M\) 对所有二次形式 \(\beta\) 成立；在正向正交余标架中，它将 \(e^1\wedge e^2\) 送到 \(e^3\)，其余两项循环处理，故逐点存在且唯一。对 Lie 代数值形式逐个系数应用它。若处处驻定，可取 \(a=*F_A\)。因为一次与二次形式交换时符号为正，
	\[
		0=\int_M\operatorname{tr}((*F_A)\wedge F_A)=-\int_M|F_A|^2\,\mathrm{vol}_M.
	\]
	非负连续函数的积分为零迫使 \(F_A=0\)。辅助度量只用于检验变分泛函非退化，原来的作用量不含这个选择。
\end{proof}

实数作用量在规范变换下可增加整数，因此对整数 \(\ell\)，相位
\[
	\exp\bigl(2\pi\mathrm i\ell S(A)\bigr)
\]
在规范变换下不变。整数条件可以直接在 \(S^3\) 上检验：恒等映射产生一次单位平移，要求相位不变就有 \(e^{2\pi\mathrm i\ell}=1\)。如果流形有边界，变分中还保留 \(-\int_{\partial M}\operatorname{tr}(Aa)/(8\pi^2)\)；固定边界联络使 \(a\) 的切向限制为零，从而消去这一项。边界条件因此是变分问题实际数据的一部分。

\ProseParagraphBreak

指标理论会使用另一种方便相加的曲率组合。对复向量丛的联络定义有限次数意义下的形式
\begin{equation}\label{g2:eq:chern-character-curvature-form}
	\operatorname{ch}(\nabla)=\operatorname{tr}\exp\left(\frac{\mathrm iF}{2\pi}\right)
	=\operatorname{rank}E+\sum_{j\ge1}\frac1{j!}\left(\frac{\mathrm i}{2\pi}\right)^j\operatorname{tr}(F^j).
\end{equation}
在给定有限维底流形上，超过维数的形式自动为零，所以这里只需有限和。每一项来自不变多项式，因而闭，且其上同调类与联络无关。它称为\term[chen-tezheng-xingshi]{陈特征形式}[Chern character form]。

\ProseParagraphBreak

直和曲率是分块对角矩阵，所以 \(\operatorname{ch}(\nabla^E\oplus\nabla^F)=\operatorname{ch}(\nabla^E)+\operatorname{ch}(\nabla^F)\)。张量丛曲率为 \(F_E\otimes I+I\otimes F_F\)，两个偶次矩阵形式可交换，指数因而分解，迹也分解，得到张量乘法公式。在分裂空间上，若陈根为 \(x_j\)，则其上同调类为 \(\sum_j e^{x_j}\)。分裂拉回单射使这个表达返回原空间；前三项为
\[
	\operatorname{ch}(E)=\operatorname{rank}E+c_1(E)+\frac12\bigl(c_1(E)^2-2c_2(E)\bigr)+\cdots
\]
的实或复系数像。特别在刚才的 \(\operatorname{SU}(2)\) 模型中，四次陈特征分量等于 \(-c_2\)，其转移形式就是 \(-\operatorname{CS}(A)\)。

\ProseParagraphBreak

第24章将把这种可加表达放进拓扑 K 理论。另一方面，解析指标是适当 Fredholm 算子的 \(\dim\ker D-\dim\operatorname{coker}D\)；把它与示性形式积分联系起来，需要构造算子、证明 Fredholm 性并建立相应指标定理，第三卷第17章将给出所选版本及其证明范围。本章已经提供其中的曲率表达、拓扑归一化和改变联络时的边界项。对 Euler 形式，\cref{g2:thm:chern-weil-euler-comparison}与上一章的切丛 Euler 求值直接给出闭定向偶维流形上的 Euler 示性数积分；一般椭圆算子的指标则还需要相应的解析输入。
% 纲要条目排为编号小节标题，并进入目录与书签；说明段仍保留于纲要注释中。


\begin{chaptersummary}
待撰写
\end{chaptersummary}

\BookExercises
待撰写

